% ============================================================
%  R. Nielsen, Evangelietroen og Theologien
%  Tolv Forelæsninger, holdte ved Universitetet i Kjøbenhavn
%  i Vinteren 1849--50.
%  Kjøbenhavn: C. A. Reitzel, 1850.  (VIII + 174 pp., Fraktur)
%  TRANSCRIPTION (Danish)
%  Source: KB scan, bibliotek/Nielsen, Rasmus/1850-evangelietroen-theologien.pdf
%          (KB e-mat/dod/11030800271D.pdf).
%  PAGE MAP:  see pagemap.py — the offset is NOT uniform.
%             printed   1-- 83  =  PDF + 13
%             [PDF 97--98 duplicate printed 82--83 in the scan; skip]
%             printed  84--174  =  PDF + 15
%             printed roman III--VIII  =  PDF 8--13
% ============================================================
\documentclass[12pt,a4paper]{book}
\usepackage[utf8]{inputenc}
\usepackage[T1]{fontenc}
\usepackage{libertinus}
\usepackage{libertinust1math}
\usepackage{textalpha}   % Greek typed directly
\usepackage{graphicx}    % for \reflectbox (the old Danish »ɔ:« id-est mark)
\DeclareUnicodeCharacter{0254}{\reflectbox{c}}  % ɔ (U+0254) = reversed c
\usepackage[danish]{babel}
\usepackage[protrusion=true,expansion=true]{microtype}
\usepackage{setspace}
\usepackage[margin=1.2in]{geometry}
\usepackage{enumitem}
\usepackage{fancyhdr}
\setlength{\headheight}{28pt}
\addtolength{\topmargin}{-13.5pt}
\usepackage{hyperref}

\hypersetup{
  pdftitle={Evangelietroen og Theologien},
  pdfauthor={R. Nielsen},
  hidelinks
}

\pagestyle{fancy}
\fancyhf{}
\fancyhead[LE,RO]{\thepage}
\fancyhead[RE]{\textit{Evangelietroen og Theologien}}
\fancyhead[LO]{\textit{\leftmark}}

\setstretch{1.3}

\title{\textbf{Evangelietroen og Theologien}\\[6pt]
  \large Tolv Forelæsninger, holdte ved Universitetet\\
  i Kjøbenhavn i Vinteren 1849--50}
\author{R.\ Nielsen}
\date{Kjøbenhavn: C.\ A.\ Reitzels Forlag, 1850}

\begin{document}
\maketitle
\thispagestyle{empty}
\newpage

\tableofcontents
\newpage

%% ============================================================
%%  FORORD (printed pp. III--VIII = PDF 8--13)
%% ============================================================
\chapter*{Forord}
\addcontentsline{toc}{chapter}{Forord}
\label{ch:forord}

% --- p. III ---
„Dersom man antager, at alle de mange Præster her og i Udlandet, som holde
og skrive Prædikener, ere troende Christne, hvorledes lader det sig saa
forklare, at man aldrig hører eller læser den Bøn, som især i vore Tider
ligger saa nær: Gud i Himlene, jeg takker Dig, at Du ikke har fordret af et
Menneske, at han skal begribe Christendommen; thi, hvis det fordredes, da var
jeg den Elendigste af Alle. Jo mere jeg søger at begribe den, jo mere
ubegribelig forekommer den mig, jo mere opdager jeg Forargelsens Mulighed.
Derfor takker jeg Dig, at Du ene fordrer Troen, og jeg beder, at Du fremdeles
vil forøge mig den. Denne Bøn vilde orthodox være aldeles correct, og antaget,
at den var sand i den Bedende, vilde den tillige være correct Ironie over hele
Speculationen. Men mon Troen findes paa Jorden!“

Disse Ord af en christelig Psycholog, der kjender „Sygdommen til Døden“,
benytter jeg som Forord for nærværende Skrift, deels fordi de have grebet mig
dybt, deels ogsaa fordi de maaskee bedst kunne forhjælpe mig til en
Adskillelse, uden hvilken mine Forelæsninger, trods alle Protester og
afværgende Forsikkringer, neppe ville undgaae Mistydning, jeg mener
Adskillelsen imellem det, at stride for Troen, og det, at strides om Troen.

At stride for Troen er guddommeligt; at strides om Troen er menneskeligt, ja
tildeels verdsligt. Læseren vil da allerede heraf kunne skjønne, hvorfor
ovenstaaende Ord ikke anbringes som Motto, men ligefrem opstilles som Forord:
et sindrigt Motto vækker en Forventning; men dette Forord staaer just her, for
at afvise en Forventning.

Hvormeget jeg end i disse Forelæsninger har bestræbt mig for at undgaae Alt,
hvad der paa nogen Maade kunde have Udseende af en blot personlig Strid mellem
Mand og Mand, har jeg dog ikke kunnet undgaae, at komme det Bestaaende saa
nær, at Læseren jo fra sit Synspunkt turde have et Slags Ret i at udtyde Noget
personligt. Hvormeget jeg end i disse Forelæsninger har bestræbt mig for fra
% --- p. IV ---
Først til Sidst, at gjøre det begribeligt, at Christendommen ikke vil begribes;
har jeg dog ikke ganske kunnet undgaae Skinnet af at stille et Begreb mod et
Begreb, en Theorie mod en Theorie. Thi at forklare, hvorfor man ikke kan
forklare, er dog altid at betroe sig til en Forklaring; at give Grunde, hvorfor
der ingen Grunde kan gives, er jo dog alligevel at indlade sig paa en
Begrundelse.

Nærværende Forelæsninger maae altsaa, ifølge deres Anlæg og Udførelse,
naturligviis finde sig i at blive henregnede til den Art Forsøg, hvori man ikke
saameget strider for Troen, som man, egentlig talt, strides om Troen.

Men at strides om Troen er dog kun menneskeligt; at stride for Troen, --- ja
det er i Sandhed guddommeligt. Det Sidste tør paa ingen Maade forvexles med det
Første; for at forebygge en saadan Forvexling, bør Adskillelsen fremhæves saa
bestemt, som muligt.

Som jeg nu sad og grundede paa denne Adskillelse, og søgte et betegnende
Udtryk, randt hine Ord mig ihu: „\emph{jeg takker Dig, Gud i Himlene, at Du
ikke har fordret af et Menneske, at han skal begribe Christendommen!}“; og
Ordene rørte sig i min Erindring, og dannede mig et Billede, et talende
Billede, Billedet af en luthersk Præst, der holder Bøn i sin Menighed.

Idet jeg saa vilde fremstille som i et Exempel, hvilken Forskjel der dog er
mellem det, at strides med Grunde mod Grunde, og det, at glemme alle endelige
Grunde i en eneste afgjørende \emph{personlig} Grund; mellem det, at strides
paa Forklaringer mod Forklaringer, og det, at gjemme alle Forklaringer i en
eneste Alt forklarende Tanke; kort sagt: hvilken Forskjel der er mellem det, at
strides om Troen, og det, at stride for Troen: da syntes det mig, at jeg ikke
kunde vælge noget bedre Exempel, end Exemplet af denne bedende Præst, der
visselig ikke vil strides med Nogen om Troen, men stride i al Inderlighed ---
og dog med Alverden --- for Troen.

At den Stridendes Bøn maa vise sig for en udenfor staaende Betragter som en
„correct Ironie over hele Speculationen“, er begribeligt; det er jo ingen
Kirkebøn for „den christelige Videnskab“, ingen Bededagsbøn for „den kritiske
Indledning“: Ironien er da heller ikke i Ordene, men i Gjerningen, ikke for den
Bedende, men for den reflecterende Betragter. Bønnen selv er en eenfoldig og
forsaavidt da ogsaa „uvidenskabelig“ Bøn i Troen for Troen; men den
% --- p. V ---
reflecterende Iagttager hører det dog strax, hvor den Bedendes Ord ligesom
uvilkaarlig forraade et hemmeligt Bekjendtskab med „hele Speculationen“.
Hvorledes skulde Præsten vel ellers vide med sig selv, at han vilde være „den
Elendigste af Alle“, hvis det fordredes af ham, at han skulde begribe
Christendommen; hvorledes skulde just denne Præst iblandt saa mange
raisonnerende, kritiserende, æsthetiserende Præster, der ellers bære paa
Kirkens „dogmatiske“ Speculationer, uden at Nogen skal mærke det paa dem, at
Byrden tynger; hvorledes --- spørger jeg --- skulde just denne Ene kunne føle
saa dybt, at han paa Speculationens Vilkaar maatte blive --- „den Elendigste!“
hvis han ikke selv, og det paa egen Risico, dog ogsaa engang havde tilladt sig
at see i „Begrebet“, og da seet saa dybt, at han med Forfærdelse saae ---
„Forargelsens Mulighed“.

Hvis der da nu virkelig overalt af alle Præster og i alle Menigheder blev
saaledes stridt i Troen for Troen; hvis hine Dage, da Spørgsmaalet om det
salige Liv var Slægtens høieste Spørgsmaal, vendte tilbage; hvis Kirkens store
Jubelaar nu var saa nær, at det sang for vore Øren, som hørte vi allerede
Klangen af de Klokker, der skulle ringe det ind: ja, da kunde vore Opvakte jo
sagtens afvise en blot „theoretisk Skolestrid“, og give dem Anviisning paa
Livets „Ildprøve“, som --- „istedetfor at gjøre derefter“! --- kun skrive
derom, forsaavidt de skrive om Troen. „Men“ --- tilføier den christelige
Psycholog --- „mon Troen findes paa Jorden!“

Vil man imidlertid, for dog at blive „Skolestriden“ med samt dens
Undersøgelser qvit, give Sagen en sindigere Vending, og erklære, at det med den
her anmældte Principstrid dog alligevel ikke har stort at betyde, efterdi
Seiren --- naar vi da endelig skulle være ganske oprigtige --- jo i Grunden
allerede forud er afgjort paa et andet Sted: kan jeg fra mit nuværende
Standpunct naturligviis ikke have det Mindste derimod, men kun glæde mig over,
at man med en saa skjøn Oprigtighed nu endelig er tilsinds at give Sandheden
Æren.

Den Strid, hvori Troesprincipet med den meest intensive Reflexionskraft kan
bekrige den dogmatiserende Speculation, er visselig udstridt i en ganske anden
og høiere Productivitet, end den, der kan tilbydes i disse Foredrag. Den Seier,
der under de forhaanden værende Betingelser kan vindes over Nutidens
æsthetiske, metaphysiske, systematiske, verdenshistoriske „Sandsebedrag“, er
--- saavidt jeg skjønner --- virkelig allerede vunden af en Anden, og det i
dybeste Stilhed, og det saa alsidig, saa fuldstændig og efter en saa omfattende
% --- p. VI ---
Maalestok, at Den, der vandt en saadan Seier, vist med Rette kunde sige ved sig
selv: „vi have meer end seiret“.

Alligevel er her dog nok endnu --- saavidt jeg kan see --- „ein Aber dabei“.

Med Krige og Seire forholder det sig nemlig noget anderledes i den aandelige
end i den sandselige Verden. Er en Mand i Sandsernes Verden bleven legemlig
ihjelslaaet, kan Enhver jo let overbevise sig om, at den Faldne virkelig er død.
Men i Aandens Verden er det anderledes. I Aandens Verden kan en Mand stundom
blive ramt saa dræbende, og dog saa behændig, at han ikke engang selv (end sige
da saa mange Andre) ret mærker, at der egentlig er hændet ham Noget. Dette vil
just være Tilfældet, jo større den Overlegenhed er, for hvilken den Overvundne
er falden. I legemlig Forstand kan det jo være latterligt nok, at behandle en
lyslevende Mand som et Gjenfærd, om end Rygtet har løiet ham død; da den
sandselige Optræden naturligviis er det meest slaaende Beviis, at Personen
endnu er ilive. Men i Aandens Verden er det anderledes. I Aandens Verden staaer
den udvortes Optræden og det phænomenale Velbefindende endda just ikke saa
sjeldent i et --- desværre! aldeles omvendt Forhold til Individets indre
Sundhed og væsentlige Livskraft; her gjælder det da især, hvad Erfaringen jo
ogsaa vil lære os, at Skinnet bedrager.

Da nu det skuffende Skin unegtelig er Sandhedens farligste Fjende; saa bliver
det i Ideernes Verden undertiden nødvendigt, at der ikke blot som i lydløs
Stilhed strides for Sandhed og Ret, men at der ogsaa med en lidt lydeligere (om
end just derfor lidt grovere) Vaabenart tillige strides om, hvo det da nu
egentlig er, der har seiret i Striden. Forsømmer man dette Sidste formeget,
kunde det let hænde sig, at den høiere ideelle Magtfuldkommenhed ---
idetmindste for en Tid --- blev forqvaklet, og en Forvirring reiste sig i den
usynlige Republik, en Forvirring, der er de høieste og helligste Interesser
endnu langt farligere, end den Forvirring, der opkommer i det Synlige, naar en
beleiret Fæstning mangler en Commandant.

Saa slutter jeg da heraf, at det endnu ikke turde være saa aldeles
betydningsløst at tilbyde en Strid om Troeserkjendelsens Princip, og beder
desaarsag den velvillige Læser at tage til Takke med mine Bidrag.

Hvor mislig min Opgave imidlertid viser sig, maa vist Enhver indrømme, der
% --- p. VII ---
tænker lidt nøiere efter. Skal her virkelig strides om, hvo det er, der har
seiret i Striden, hvor nær ligger da ikke den Mistydning, at den foregivne
Interesse for Opfattelsen af et Princip dog maaskee alligevel kun tjener som et
af de Paaskud hvormed en selvisk Lidenskab vil besmykke en anmassende
Forfængelighed.

Skulde denne Mistydning (hvad jeg dog ikke antager) formaae at trænge igjennem
og sætte sig fast som authentisk Fortolkning: maa jeg da bede om, at
Beskyldningen fremsættes saa aabenlyst og saa ganske uden al
\textit{reservatio mentalis}, at den udtrykkelig kan falde paa mig alene.

Jeg frembærer denne Anmodning, ikke for at affectere et Heltemod, jeg ikke
besidder, men fordi jeg anseer det for min Pligt at forsøge Alt, hvad der paa
nogen Maade kan bidrage til at holde denne Sag saa reen, som muligt.

At en Person, der hidtil har levet fredsommelig, stille, tilfreds, som Den, der
havde Nok i sit Eget, pludselig farer ud og begynder at støie; at en Mand, der
ellers ikke vides at have overtraadt Velanstændigheds Grændse, pludselig møder
i spraglet Dragt, og med fremmede Fjer, og som En, der vil til at triumphere;
at en ellers ædruelig Forfatter, hvem Litteraturens Myndige hidtil ikke have
frakjendt en vis Dannelse, Besindighed og Tanksomhed, med Eet faaer isinde at
gebærde sig, som var han enten bleven stukken af en fix Idee, eller ogsaa i den
Grad besat af Forfængelighedens Dæmon, at han, i Mangel af at være Noget selv,
maatte skabe sig efter en Anden: at Sligt kan hændes, hvor sørgeligt, hvor
beklageligt! at Sligt virkelig hændes, hvor sørgeligt, hvor beklageligt for
Den, det hændes; men --- naar det virkelig hændes --- fordærver det dog
alligevel ikke den, iøvrigt alvorlige, Sag, en saadan Prostitueret indbilder
sig at forfægte.

Hvis derimod en Forfatter med lidt sundere Sands gav sig til at behandle en
Troessag, og nu dog alligevel (være sig af Anmasselse eller Ubehjælpsomhed) fik
det Høiere saaledes indspundet i det Lavere, det Religiøse i det Verdslige, det
Principielle i det Personlige, at han ikke blot ved sine høirøstede Indsigelser
fortrædigede det Bestaaende (som i de gamle Former maaskee dog ellers endnu
kunde bevare meget baade Ædelt og Godt), men til denne relative, for det
Religiøse selv høist underordnede, Tvistighed endog forsøgte at drage „en
Anden“ med, „en Anden“, som dog ellers aldrig er med, --- hvor der blæses i
Basune, „en Anden“, som dog ellers frabeder sig enhver udvortes Forfremmelse,
„en Anden“, der ellers ingenlunde attraaer at vorde udmærket i Krige, naar han
% --- p. VIII ---
saa blot maa være ubemærket i Fred, „en Anden“, der jo ellers saa indstændig
beder, at han som Uvedkommende altid maa være glemt, og Tanken, kun Tanken,
erindres: da burde en saa uhyggelig Fremfærd vistnok fra begge Sider afvises
med energisk Protest, og hvert Ord i det uhjemlede Forsøg mortificeres, som
Retten mortificerer et Uqvemsord.

Dette og meget Andet har jeg nu drøftet og overveiet paa bedste Maade; men med
alt Dette maa jeg dog endnu (idetmindste foreløbig) bekjende mig til den
Mening, at man dog endnu --- \textit{propter iniquitatem temporum}! dog endnu
--- som Sagerne staae \textit{in republica litterarum}! dog endnu burde gjøre
et Forsøg og strides om Troen.

Skulde det omsider komme saavidt, at man ikke engang fandt det Umagen værd at
strides om Troen, mon det da ogsaa var et godt Tegn? mon det ikke meget mere
var et Tegn paa, at man foragtede denne menneskelige Strid, fordi man i Grunden
heller ikke agtede den guddommelige; et Tegn paa, at man endnu mindre fandt det
Umagen værd at stride for Troen: et Tegn paa, at Troen ikke mere fandtes paa
Jorden?

Forfatteren af nærværende Skrift tilstaaer oprigtig, at han endnu hverken i Ord
eller Gjerning ret forstaaer at stride for Troen, som det sig bør; men i det
Haab, at det Guddommelige engang maatte vorde ham til Deel, naar han kun ikke i
altfor sandselig Sikkerhed altfor meget forsømmer det Menneskelige, har han, om
just ikke til Manges Belæring, saa dog idetmindste til egen Bestyrkelse, villet
udtale sig i disse Forelæsninger angaaende den Maade, hvorpaa man for
nærværende Tid --- strides om Troen.

\bigskip
\noindent Kjøbenhavn, d.\ 17 Februar 1850.
\hfill R.\ Nielsen.

%% ============================================================
%%  FØRSTE FORELÆSNING (pp. 1--13 = PDF 14--26)
%% ============================================================
\chapter*{Første Forelæsning}
\addcontentsline{toc}{chapter}{Første Forelæsning. Indledende Betragtninger}
\label{ch:forel1}
\markboth{Første Forelæsning}{}

\begin{center}\itshape
Indledende Betragtninger. Christendommen er høiere end Videnskaben,
og Evangelietroen ueensartet med Theologien.
\end{center}

% --- p. 1 ---
Det er Forholdet mellem Evangelietroen og Theologien, der gjøres til
Gjenstand for disse Forelæsninger. Istedetfor at sige: „Evangelietroen“,
kunne vi strax sige: „Christendommen“; derimod kunne vi ikke uden videre
gjøre „Theologien“ eensbetydig med „Videnskaben“. Spørgsmaalet om
Evangelietroens Forhold til Theologien viser da foreløbig hen til et andet
og mere almindeligt Spørgsmaal, nemlig om Christendommens Forhold til
Videnskaben.

Som bekjendt, bliver dette Forhold til forskjellige Tider og i de
forskjellige Skoler opfattet høist forskjelligt, snart som et Stridens,
snart igjen som et Forsoningens Forhold.

I Striden mellem Christendommen og Videnskaben ere der da Nogle, der gaae
saa vidt, at de i Videnskabens Navn fordre Christendommen udryddet, medens
Andre derimod i deres Iver for at fremme Christendommen prædike Korstog
mod den frie Videnskab.

I Følelsen af, at der dog umulig kan findes Strid i Sandhedens Væsen, og
at Sandhedens Rige ikke kan være splidagtigt med sig selv, har man saa
igjen arbeidet paa at tale de stridende

% --- p. 2 ---
Parter tilrette og tilveiebringe en Forsoning ved lidt efter lidt at gjøre
Christendommen videnskabelig og --- hvad der kommer ud paa det Samme ---
Videnskaben christelig.

Hvilken af disse Bestræbelser skulle vi nu skjænke vort Bifald? Det er
just ikke saa let at sige. Den hele Sag synes saa indviklet, at vi ikke
kunne bifalde noget enkelt af de her antydede Forsøg, uden at bifalde dem
alle, ei heller forkaste noget enkelt, uden at forkaste dem alle. Thi hvis
De, der ville Striden, have Ret, saa maae jo ogsaa De have Ret, der ville
Forsoningen. At indbyde til Forsoning er her det Samme, som at
forudsætte Strid og ægge til Strid; og omvendt, at forkynde Strid er her
det Samme, som at indbyde til Forsoning.

Denne dunkle Tale vil maaskee klare sig og blive forstaaelig, naar jeg nu
fremsætter min Thesis, den dobbeltsidige Grundtanke, hvorom disse Foredrag
skulle bevæge sig, den Thesis nemlig, at \emph{Christendommen er høiere
end Videnskaben, og Evangelietroen ueensartet med Theologien.}

Det første Led i denne Thesis er hypothetisk; det sidste er kategorisk.

At Christendommen er høiere end Videnskaben, lader sig ikke godtgjøre ved
nogetsomhelst almeengyldigt Beviis. Et almeengyldigt Beviis herfor maatte
jo være et videnskabeligt; men et videnskabeligt Beviis, for at
Christendommen er høiere end Videnskaben, er en Modsigelse.

I hypothetisk Form kan Sætningen derimod udtrykkes saaledes: „hvis
Christendommen er, hvad den siger sig at være, en aabenbaret Sandhed, da
maa denne Sandhed tilhøre en Sphære, der er høiere end den, hvortil
Videnskaben kan hæve sig“; og omvendt: „er Videnskabens Sphære den
høieste Sandheds høieste Sphære, saa er Christendommen tabt, og Troen
afskaffet“.

At nu Videnskabens Sphære nødvendigviis maa være Sandhedserkjendelsens
høieste Sphære, kan Videnskaben ikke bevise;

% --- p. 3 ---
thi, idet den paatager sig at føre Beviset, begynder den netop med at
forudsætte Det, den egentlig vilde bevise, og kommer ikke ud over sin egen
Cirkel.

Da nu Videnskaben altsaa hverken kan bevise eller modbevise
Christendommens Sandhed; saa kunne vi med al Grund forbinde vore to
hypothetiske Sætninger i den kategoriske Paastand, at Christendommen i
ethvert Tilfælde falder udenfor Videnskaben, eller med andre Ord, at
Christendommen er ueensartet med Videnskaben.

Skulde denne Paastand --- endnu er det jo kun en Paastand --- lade sig
forsvare, saa følge strax heraf trende Ting:

a) at den heftige Strid mellem Christendommen og Videnskaben maa bero paa
en Misforstaaelse; thi imellem ueensartede Magter kan der jo egentlig
talt ingen ligefrem Strid være. Videnskaben kan som Videnskab umulig
ville krænke Christendommen, saalidt som Christendommen kan ville hemme
eller fornærme Videnskaben, hvis det er sandt, at de bevæge sig i ganske
forskjellige Sphærer.

b) at den tilstræbte Forsoning mellem Christendommen og Videnskaben maa
hidrøre fra en stor Misforstaaelse. Er Striden en Indbildning, maa
Forsoningen ogsaa være en Indbildning; thi en virkelig Forsoning
forudsætter en virkelig Strid.

c) at den eneste Betingelse, hvorunder Videnskaben kan komme til at
krænke og fornærme Christendommen, er, at den drager den ind i sin
Sphære, og gjør den eensartet med sig; og at Theologien blandt alle
Videnskaber følgelig er den eneste Videnskab, der ret for Alvor kan
krænke og fornærme Christendommen.

Saa anstødelig denne sidste Paastand end maaskee vil være for Mange, saa
fremgaaer den dog med umiskjendelig Conseqvens af den opstillede Thesis.
Vor Undersøgelse fæster sig altsaa til det ene Spørgsmaal, om
Christendommen da nu ogsaa virkelig

% --- p. 4 ---
er ueensartet med enhver Videnskab, af hvad Navn nævnes kan, eller ikke.

Christendommen forkynder sig, ligesom Jødedommen, at være en overnaturlig
Aabenbaring. Gud taler, Underet skeer: dermed begynder Aabenbaringen.
Denne Begyndelse er af uvidenskabelig eller, om man saa vil, af
overvidenskabelig Natur, og kan aldrig blive Gjenstand for rationel
Begrundelse.

Hvad Jehova siger og gjør i det Gl. Tst., er incommensurabelt med den
menneskelige Forstand, og kan ikke begrundes videnskabeligt; hvad Sønnen,
„den Eenbaarne“, siger og gjør i det N. Tst., er incommensurabelt med den
menneskelige Forstand, og kan ikke begrundes videnskabeligt. Moses og
Propheterne, Christus og Apostlene ere uvidenskabelige, det vil sige:
overvidenskabelige Skikkelser. Ingen af dem kommer i Videnskabens Navn,
men de komme jo allesammen frem, og tale og handle i den aabenbarede Guds
Navn.

Underet er i sit Væsen over og derfor i sin Tilsyneladelse imod Naturens
Lov, og kan ikke blive Gjenstand for videnskabelig Behandling. Moses
slaaer paa Klippen, og Kilden udvælder: dette Factum kan Naturlæren ikke
bruge, det vedkommer den slet ikke. Christus, „Menneskens Søn“, der
gaaer sin Lidelse og Korsfæstelse imøde, siger i Bønnen, at han havde
Herlighed hos Faderen, før Verden blev; denne Tale kan Philosophien ikke
udgrunde, men vel see og forstaae, at, hvis Talen er sand, da maa det
være en overvidenskabelig, en guddommelig Tale.

Er Aabenbaringen i sit Princip ueensartet med Videnskaben: hvorledes
skulle de da komme i Strid med hinanden?

Aabenbaringen mødes med Naturlæren. Naturlæren forsikkrer os nu, at,
hvis Mirakler ere skete, saa maae de være skete mod Naturens Orden; men
det Selvsamme lærer jo ogsaa Aabenbaringen. Lad saa Naturlæren udforske
Tingene og opdage Lovene; lad den i fortsat Undersøgelse overbevise sig
om, at der i

% --- p. 5 ---
Naturens Rige intet Spring er, i Lovenes System ingen Plads for de
overnaturlige Kjendsgjerninger; lad den kun gaae videre, lad den med
sindrig Klogt samle sin rige Erfarings mangfoldige Analogier, for at
gjøre os disse Kjendsgjerninger saa urimelige, saa usandsynlige, som
muligt: den kommer dog alligevel ikke til at krænke Christendommen; thi
Christendommen paastaaer jo ingenlunde, at dens Kjendsgjerninger skulde
rime sig med den sandselige Erfarings Analogier. Naturlæren kan saavist
gjerne forsikkre, at Christendommens Factum ikke lader sig forklare af
Naturkræfternes Sammenhæng; naar den kun ikke forsikkrer Mere (og Mere er
den da vel, som Videnskab, ikke berettiget til at forsikkre): saa krænker
den ikke Christendommen. Den kan endogsaa, hvis den finder sig opfordret
dertil, gjøre et Skridt endnu, og sige, at Aabenbaringens Phænomener ere
saa stridende mod Naturens Love, at de endog vise sig absurde i Forhold
til de naturlige Phænomener; naar den kun ikke siger Mere (og Mere er
den, som Videnskab, dog vist ikke berettiget til at sige): den krænker
dog alligevel ikke Christendommen; tvertimod! her er nu netop i denne
Henseende den allerbedste Forstaaelse mellem Naturlæren og
Christendommen, eftersom de jo, hver paa sit Maal, i Grunden sige det
Selvsamme.

Aabenbaringen mødes dernæst med Philosophien, og Synskredsen udvides.
Naturlæren og Christendommen kunne dog endnu lettere bevare
Ueensartetheden; kun som i Forbigaaende vil Naturlæren streife det
Overnaturliges Enemærker, og føler sig sjælden fristet til at trænge
dybere ind. Anderledes med Philosophien, der, idet den prøver Alt, endog
gjør Christendommen selv til Gjenstand for videnskabelig Prøvelse. Men er
dette da tilladt? Naturligviis! Saa liberal er Christendommen, at den end
ikke vægrer sig ved at blive Gjenstand for Philosophiens allerstrengeste
Undersøgelse (hvis den vægrede sig herved, vilde den jo krænke
Videnskaben). Saa liberal er Christendommen, at Philosophien

% --- p. 6 ---
altsaa faaer Lov til at gjøre med den, hvad den vil, erklære den for
begribelig, hvis den kan forsvare det for Videnskaben, forkaste den som
ubegribelig, hvis den kan forsvare det for Videnskaben, men --- vel at
mærke --- paa een Betingelse: at Philosophien med alt Dette saa ikke
tillige falder paa, at ville gjøre sig eensartet med Christendommen.

Bliver denne Betingelse ubrødelig holdt, da har det vist ingen Fare,
hverken med Philosophien eller med Christendommen. Erklærer Philosophien,
at den nu endelig har begrebet eller dog er paa gode Veie til at begribe
Christendommen; har det dog ingen Fare, naar det kun siges i
Videnskabens Navn. Fordi Christendommen, \emph{videnskabelig forstaaet},
engang imellem viser sig som begribelig (eller dog næstendeels
begribelig), derfor kan den jo, \emph{christelig forstaaet}, ligefuldt
være, hvad den er, og hvad den altid har været, underfuld, ufattelig, for
os syndige Mennesker ubegribelig. Og, fordi Philosophien engang imellem
farer op, og vil paa Videnskabens Vegne foreholde den dannede Verden, at
Christendommen nu endelig har viist sig i sit rette Lys, som ubegribelig,
absurd og meningsløs, er her dog alligevel ingen virkelig Fare, thi her
er jo i Grunden ingen virkelig Strid. Nei, langtfra! Philosophien er nu,
uden at den maaskee i Heden selv ret mærker det, just nu paa Veien til
den allerbedste Forstaaelse med Christendommen. Christendommen har jo
aldrig lovet, at den nogensinde vilde lade sig begribe (eller dog
næstendeels begribe) af Philosophien. Christendommen er saa langt fra at
ængstes for denne philosophiske Opdagelse, at den tvertimod ønsker
Videnskaben til Lykke med Opdagelsen, idet den saa gjerne seer, at
Philosophien, der søger Sandheden, dog endelig engang maatte komme til
Sandheds Erkjendelse, d. e. til Erkjendelse af, at den er ueensartet med
Christendommen.

Saa er det da forunderligt nok med dette Forhold mellem Videnskaben og
Christendommen. De strides, forsones,

% --- p. 7 ---
og strides igjen. Ifølge den herskende Synsmaade, ængstes man meget for
Striden og grunder dybt paa Forsoningen; og Striden vilde være endnu
haardere, og Forsoningen forventes endnu inderligere, hvis ikke Verden
allerede var bleven saa klog, at den, i Følelsen af, at der vel i Grunden
ikke kommer synderligt ud enten af det Ene eller af det Andet, allerede
var begyndt at blive lidt flau ved det Hele.

Her stiller Sagen sig anderledes. Istedetfor at ængstes for Striden og
haabe paa Forsoningen, som om det forstod sig af sig selv, at baade
Striden og Forsoningen havde sand Gyldighed og overvættes Betydning, ere
vi nu tvertimod stillede saaledes, at vi ikke ret kunne see, hvad alle
disse Ophævelser egentlig føre til, da det Hele synes at bero paa en
Misforstaaelse.

Her maa dog være en Grund til denne Misforstaaelse; her maa være et
Medium, hvori de ueensartede Magter mødes med hinanden, og gribe ind i
hinanden saaledes, at de idetmindste faae Skin af at være eensartede:
hvor er nu denne Grund? hvilket er dette Medium? --- det existerende
Menneske!

I den menneskelige Existens, d. e. i det existerende Menneske ere
ueensartede Størrelser forbundne til Eenhed. Det Sandselige og det
Aandelige, det Timelige og det Evige, det Naturlige og det Overnaturlige
mødes her. Det er Menneskets almindelige Opgave at forholde sig
forskjelligt til det Forskjellige. Men, naar Aabenbaringen træder til,
bliver Forskjellen og Modsætningen saa stærk, at Mennesket, om han end
føler med sig selv, at Sandheden kun er een, og kun kan være een, dog
alligevel føres ind i en saa dyb Sandhedsdualisme, at han i denne
Existens maa forholde sig aldeles forskjellig til tvende i Existensen saa
ueensartede Sandheder, som: \emph{Aabenbaringssandheden og den
almindelige Fornuftsandhed.}

Heraf forklares det, hvorledes Striden imellem Christendommen og
Videnskaben opkommer i Christenheden,

% --- p. 8 ---
Den Christne bliver naturkyndig; den Naturkyndige er en Christen. Begge
ueensartede Bestemmelser forbindes i een Bevidsthed. Opgaven er, at det
paagjældende Individ nu skal forstaae at forholde sig christeligt til
Christendommen og videnskabeligt til Videnskaben, saa at han af
Christendommen ikke udleder Slutninger for Naturlæren, eller omvendt fra
Naturlæren udrager Conseqvenser for Christendommen. Men, da det dog er
det samme Menneske, den samme Bevidsthed, der her skal forene det
Ueensartede i existentiel Eenhed: hvor nær ligger saa ikke den Fristelse,
at gjøre de ueensartede Størrelser lidt eensartede, at behandle
Videnskaben christeligt og Christendommen videnskabeligt, at vække
Spliden ved at tilstræbe Forsoningen.

Den Christne bliver Philosoph; Philosophen er en Christen. Det gjælder
nu om at forholde sig christeligt til Christendommen og videnskabeligt
til Videnskaben. Men, da Christendommen jo forkynder sig som Sandheden,
og da Videnskaben, \textit{in specie} Philosophien, undersøger og prøver
Sandheden: hvor nær ligger det da ikke, at gjøre Christendommens
religiøse Erkjendelse eensartet med den philosophiske Erkjendelse af
Christendommen, at forsone Troen med Viden, og just derved vække Splid
mellem Tro og Viden!

Naturlæren og Philosophien ere imidlertid med Hensyn til Princip og
Udgangspunct dog endnu saa forskjellige fra Christendommen, at man let
bliver Ueensartetheden vaer, naar Striden blot føres conseqvent. Men der
gives jo endnu en Videnskab, en med Christendommen langt inderligere
beslægtet Videnskab, en Videnskab, der med en vis Selvfølelse kalder sig
selv den christelige, en Videnskab, der paastaaer at have Princip og
Forudsætninger tilfælles med Christendommen, en Videnskab, der absolut
vil forsvare Christendommen, for at forsvare sig selv med Christendommen,
en Videnskab, der trykker sig saa nær op til Christendommen, at det er,
som skulde de leve og døe tilsammen;

% --- p. 9 ---
denne Videnskab er Theologien: er den christelige Theologie da nu ogsaa
ueensartet med Christendommen? --- Ganske vist! Theologien som
Videnskab er tilvisse ligesaa ueensartet med Christendommen, som enhver
anden Videnskab. At Forholdet her er mere tvetydigt, forklares ganske
simpelt deraf, at Theologien selv er en Tvetydighed.

Hvorledes det med denne Tvetydighed nærmere hænger sammen, er ikke
vanskeligt at opdage.

Troen kan ikke være uden Erkjendelse. Den Troende maa vide, hvad det er,
han troer, maa kunne aflægge sig selv og Andre Regnskab for sin Troes
Indhold og det Samfund, hvortil han bekjender sig. Forsaavidt Theologien
nu udelukkende sætter sig den Opgave, at udfinde, hvad det Christelige
er, at udvikle og forklare, hvad det er, Kirken troer, og den Troende
bekjender, forsaavidt er Theologien som Theologie vistnok eensartet med
den objective Lære om Christendommen; men denne objective Lære om
Christendommen er derfor endnu ingenlunde eensartet med Christendommen
selv. Forsaavidt Theologen forholder sig christeligt til Christendommen,
forholder han sig ikke videnskabeligt; forsaavidt han forholder sig
videnskabeligt til Christendommen, forholder han sig ikke christeligt.

Den gamle orthodore Theologie, der i alle Punkter vilde være conform med
Kirkelæren, led af en stor Tvetydighed. Som den alsidige, fuldstændige og
nøiagtige Fremstilling af Kirkens Lære var denne Theologie i Grunden
intet Andet end en articuleret Bekjendelse af Troens Indhold, og som
saadan endnu ikke Videnskab. Den var i det Høieste kun
Troeserkjendelsens videnskabelige Udtryk og Form. Men neppe havde den i
Formen begyndt at antage Videnskabens Præg, og altsaa gjort det første
Skridt til en Forsoning af Christendom og Videnskab, før den ogsaa fandt
sig bestedt i en haard Selvmodsigelse. Indholdet var christeligt, Formen
videnskabelig; dette var Modsigelsen mellem

% --- p. 10 ---
Form og Indhold, en Modsigelse, for hvilken den orthodore Theologie
omsider bukkede under.

Den orthodore Theologie fremstiller enhver Læresætning thetisk,
forsaavidt den fremstiller den i Troens Navn; dette er det Christelige.
Men saa behandler den tillige enhver Læresætning hypothetisk, forsaavidt
den behandler den i Lærdommens, Undersøgelsens og Forskningens Navn;
dette er saa det Videnskabelige. Ifølge heraf bliver da enhver
Læresætning at betragte fra et dobbelt Synspunkt. Christelig forstaaet
tilhører den Kirken, er ophøiet over Videnskaben, og trodser ethvert
videnskabeligt Angreb. Videnskabelig forstaaet stammer den fra theologisk
Forskning, theologisk Discussion og theologisk Prøvelse, hjemfalder
altsaa til Videnskaben, og maa som saadan blive Gjenstand for en ny
Discussion, en ny Forskning, en ny Prøvelse. Under den fortsatte
Discussion staaer nu Forskning mod Forskning, Prøvelse mod Prøvelse.
Skal Forskningen betyde Noget for Christendommen, maa den rette sig
efter Kirken; skal den gjælde Noget i Videnskaben, maa den være fri og
uafhængig af Kirken. Den theologiske Prøvelse maa altsaa paa een Gang
være baade fri og ufri, baade udvortes afhængig og dog udvortes
uafhængig; det er en Modsigelse. Den theologiske Læresætning viser sig
her som en tvetydig Mynt: paa den ene Side: Korsets Tegn! paa den anden:
Minervas Ugle! Her hjælper det ikke engang at sige: giv Keiseren, hvad
Keiserens er; thi denne tvetydige Mynt tilhører jo med lige Ret baade Gud
og Keiseren.

Theologien, der i Begyndelsen vilde være Eet med Kirkelæren, begyndte saa
efterhaanden under den frie Forskning, den frie Discussion og den frie
Prøvelse at frigjøre sig fra sit oprindelige Afhængighedsforhold og
forvandle sig fra orthodox til heterodox: jo mere heterodox den blev,
desto mere videnskabelig, og omvendt: jo mere videnskabelig, desto mere
heterodox. Først

% --- p. 11 ---
i det sidste Aarhundrede har Theologien fattet sig som selvstændig
Videnskab; først i det 19de Aarhundrede sees det klart, at Theologien er
ueensartet med Christendommen.

Men Christendommen er eensartet med Evangelietroen.

Christendommen er kun for den Troende; Theologien er kun for den
Vidende. Evangelietroen lever i Kirken; Theologien lever i Videnskaben.
Ville de vedblive at forholde sig eensartede, maae de arbeide hinanden
imod, og skulle da snart forarbeide deres egen gjensidige Opløsning.

Evangelietroen lever i Kirken. Hvad er da Kirken, og hvor er den? Vi
svare med den augsburgske Confession: „\emph{Kirken er de Helliges
Samfund, i hvilket Evangeliet rettelig forkyndes, og Sacramenterne
rettelig forvaltes}“ (\textit{Ecclesia est congregatio sanctorum, in qua
evangelium recte docetur et recte administrantur sacramenta}).

Men Theologien, hvad er den? hvor er den? Theologien er den Videnskab,
der paatager sig at underkaste Kirkens Evangelium en historisk-kritisk
Prøvelse; den Videnskab, der efter Tid og Leilighed tillader sig at
raisonnere for og imod Herrens Sacramenter; den „kirkelige“ Videnskab,
der vel ikke tør lade sig høre i Kirken, men desto friere tør lade sig
høre i Skolen.

Ueensartetheden er nu iøinefaldende.

Kirken er det Troens Samfund, i hvilket Evangeliet \emph{rettelig
forkyndes}. Men Evangeliet forkyndes kun rettelig, naar det forkyndes som
Guds Ord. Guds Ord er et ufeilbart, bestemt og myndigt Ord. Derfor gaaer
Ordets frie Forkyndelse i Menigheden stedse tilbage til det aabenbarede
Grundord; den evangeliske Prædiken maa have sin Text: den hellige Skrift
afgiver Texten. Den bibelske Text oplæses for Menighedens Andagt;
Andagten er uvidenskabelig; Andagten kjender ingen Kritik: den forelæste
Text er i alle Maader ufeilbar og troværdig, som det sig bør. „Saamange
ere Ordene,“ siger jo Præsten, og derpaa tilføier han en Bøn. Saaledes
begynder den evangeliske

% --- p. 12 ---
Prædiken da ikke med raisonnerende Kritik, men med Bøn og Paakaldelse, og
kun saaledes kan Evangeliet rettelig forkyndes.

Anderledes i Theologien; thi Theologien er jo en Videnskab. Videnskaben
kan, som Videnskab, ikke begynde med Bøn og Paakaldelse. Skal altsaa
Kirkens Evangelium prøves i Videnskaben, maa det behandles paa en ganske
anden Maade. Det skeer. I Videnskaben sidder „den kritiske Indledning“
med Evangeliet for sig, og fritter det ud, om det er ægte eller maaskee
uægte; og jo længere man spørger, og jo mere man fritter, desto større
Mistanke faaer man om, at dette Evangelium dog maaskee er --- uægte. Saa
træder „Hermeneutiken“ til, og vil i det samme Evangelium videnskabelig
adskille Aanden fra Bogstaven, Ideen fra Indklædningen, Sandheden fra
Forestillingen; og jo længere man sondrer, jo skarpsindigere man
adskiller, desto Mindre bliver der af Guds Ord og desto Mere af
Menneskeordet i den bibelske Text.

Skal heraf uddrages nogen Slutning, da maa det vel være den Slutning, at,
hvad Kirken sætter i Tro, omsætter Theologien i Tvivl.

Kirken er det Aandens Samfund, i hvilket \emph{Sacramenterne rettelig
forvaltes}. Men Sacramenterne forvaltes kun rettelig under den ene
afgjørende Forudsætning, at Indstiftelsens Ord er et ufeilbart,
guddommeligt, undergjørende Ord. Kirken bekjender, at Daaben er et
Gjenfødelsens Bad formedelst den Hellig Aand, og Nadveren Christi Legems
og Blods Samfund til Syndernes Forladelse for dem, som troe. Kirken
bygger paa Indstiftelsens Ord, men lægger sig slet ikke efter at begribe
Sacramentets Mysterium. Som Ordet lyder over Daaben, skal det troes ---
mod Forstand; som Ordet lyder i Nadveren, skal det troes --- mod
Forstand. Underet, Naademidlets underfulde Virkning, skal slet ikke
forklares, ikke begribes: Tilegnelsen er kun i Troen.

% --- p. 13 ---
Anderledes i Theologien; thi Theologien er jo Videnskab. Som objectiv
Videnskab maa Theologien naturligviis gjøre Sacramenterne selv til
kritiske Objecter, og behandler da Naadens Midler som Midler til at
forhøie Idee-Erkjendelsens videnskabelige Glands. Saa bliver det igjen
tvivlsomt, høist tvivlsomt, hvorledes man skal forklare sig det
Overnaturliges Forening med det Naturlige. Vanskeligheden er, at man ikke
kan forklare det rigtig, og dog saa inderlig gjerne vilde have det
forklaret. Jo længere man saa forklarer, desto mere tvivlsom bliver
Forklaringen.

Skal heraf uddrages nogen Slutning, da maa det vel være den Slutning, at,
hvad Kirken sætter i Tro, omsætter Theologien i Tvivl.

At nu Theologien paa sin Viis kan have Ret i at indlade sig paa
videnskabelige Undersøgelser og derhos tillige Ret i at indlade sig med
Tvivlen, skal ikke benegtes; men hvad der paa det Bestemteste maa
benegtes, er, at Theologien skulde have Ret til at forvexle sin Opgave
med Kirkens. Det er ingen Tilfældighed, at Kirken lader Alt afgjøres i
Tro; det er heller ingen Tilfældighed, at Theologien opløser Alt i
Tvivl. Kirken nedbryder sig selv, naar den tvivler om sin Myndighed, og
opgiver Troen; Theologien fornegter sig selv, naar den ved et Magtsprog
afbryder sine Undersøgelser, for at standse Tvivlen.

Som Tvivlen ikke kan nære sig i Kirken, saa kan Troen ikke nære sig i
Theologien. Evangeliet i Kirken bekræfter Troen, som Troen Evangeliet;
men Tvivlen i Theologien strider jo imod Evangeliet, som Evangeliet mod
Tvivlen.

Skal heraf uddrages nogen Slutning, da maa det vel være den Slutning, at
Tvivlen er Kirkens Fiende, og \emph{Evangelietroen ueensartet med
Theologien}.

%% ============================================================
%%  ANDEN FORELÆSNING (pp. 14--25 = PDF 27--38)
%% ============================================================
\chapter*{Anden Forelæsning}
\addcontentsline{toc}{chapter}{Anden Forelæsning}
\label{ch:forel2}
\markboth{Anden Forelæsning}{}

\begin{center}\itshape
I Opfattelsen af Forholdet mellem den historiske og religiøse Tro maa
Theologien adskille det i Religionen Uadskillelige, og derfor er
Evangelietroen ueensartet med Theologien.
\end{center}

% --- p. 14 ---
Den positive Aabenbaring --- i Jødedom og Christendom --- er en historisk
Aabenbaring. Som Aabenbaring taler den til det indvortes Menneske og
forudsætter den subjective religiøse Tro (\textit{fides religiosa}); som
historisk Aabenbaring virker den paa det naturlige Menneskes Forestilling
igjennem et Indbegreb af overnaturlige, men dog virkelige Kjendsgjerninger, og
forudsætter da tillige en historisk Tro (\textit{fides historica}).

Den, der nu forholder sig christeligt til Christendommen, kan og tør ikke
adskille det Historiske fra det Religiøse. I den sande og levende Tro maae
begge Factorer virke sammen uden dog at smelte sammen. Da nu den historiske
Tro, naar den betragtes for sig, er ueensartet med den subjective religiøse
Tro, medens den Troende dog alligevel forener disse tvende ueensartede
Bestemmelser i Evangelietroens Eenhed: er det den raisonnerende Forstand som
Følge heraf saare vanskeligt at forstaae den Religiøse.

% NB (p. 15): the second speech of the Religious man is printed with ONE
% opening „ but TWO closing “ (after „tredie Dag“ and again after „derpaa“).
% Transcribed exactly as printed; this is a defect in the 1850 setting.
„De hellige Kjendsgjerninger blive jo ikke virkelige, fordi jeg troer paa dem;
men omvendt, fordi jeg veed, at de ere virkelige, ere historiske, derfor troer
jeg paa dem. Det er skeet, at
% --- p. 15 ---
Gud har aabenbaret sig for Abraham; det er skeet, at Gud fordum mange Gange og
paa mange Maader har talet til Fædrene ved Propheterne; det er Altsammen skeet,
som Skriften vidner, og derfor troer jeg derpaa.“ Saaledes taler den Religiøse.
--- „Det er skeet, at Jesus Christus, Guds eenbaarne Søn, født af Faderen fra
Evighed, blev i Tidens Fylde undfangen af den Hellig Aand og født af Jomfru
Maria; det er skeet, at han blev piint under Pontius Pilatus, at han blev
korsfæstet og døde og blev begraven og nedfoer til Helvede og opstod igjen paa
den tredie Dag“, det er Altsammen skeet, som Skriften vidner, og Kirken
bekjender; og efterdi det visselig er skeet, saa troer jeg derpaa“. Saaledes
taler den Religiøse.

Men for den raisonnerende Forstand er det alligevel meget vanskeligt, at
forstaae denne Tale. Det seer bestandig ud, som om den Religiøse ligefrem
byggede sin Tro paa en ligefrem historisk Vished. Det seer ud, som om
Aabenbaringens Kjendsgjerninger ifølge hans Mening havde saa mange udvortes,
enkelte, endelige Vidnesbyrd, at Ingen kunde forkaste disse Kjendsgjerninger
uden at forkaste al Historie. Den Religiøse maa da forekomme de Dannede (især
de theologisk-videnskabelig Dannede) som et uvidenskabeligt, udannet og
eenfoldigt Menneske, der har saamegen Tro, fordi han kjender saa lidt til
Tvivlen. Mangen En vilde da maaskee sige: „Vidste jeg kun, at det Altsammen var
skeet, som det staaer i Skriften; vidste jeg kun, at dette Ene var skeet, at
Christus, efterat han virkelig var død og begraven, opstod igjen paa den tredie
Dag; vidste jeg nu blot, at dette Ene virkelig var skeet: ja da skulde jeg vel
troe og i Troen blive et nyt og bedre Menneske. Men hvo borger os for, at disse
Fortællinger ere ægte, ere paalidelige og i alle Maader troværdige? Den Troende
er saa vis i sin Sag, fordi han slet ikke kjender Vanskelighederne: han skulde
blot begynde med Undersøgelserne, han skulde kjende Indvendingerne, han skulde
blot prøve
% --- p. 16 ---
paa at høre et theologisk Cursus over den kritiske Indledning, og han vilde da
vel snart komme paa andre Tanker“. --- Saaledes dømmer den raisonnerende
Forstand; og dog er det en Misforstaaelse.

Den Religiøse kan, naar det endelig skal være, meget vel gjøre sig bekjendt med
de theologiske Vanskeligheder uden at tage Skade paa sin Tro. Med den „kritiske
Indledning“ har det slet ingen Fare. Kunde de tre Mænd gaae uskadte igjennem den
gloende Ovn: skulde den Troende da ikke ogsaa kunne gaae uskadt igjennem den
„kritiske Indledning“! Møde vi saa den Religiøse, som han træder ud af
Kritikens Labyrinth, og spørge ham, hvorledes det nu staaer sig med hans
Overbeviisning, da svarer han os, at han jo nu rigtignok har faaet Meget at
vide, som han forhen ikke vidste, men at han dog derved kun ydermere har
overbeviist sig om, hvad han allerede forud vidste, at Menneskenes tvivlsomme
Tanker maae bøie sig for Guds aabenbarede Tanker, at Menneskenes Gjerninger
Intet formaae imod Guds Gjerninger, og alle menneskelige Skrifter Intet imod
Guds egen hellige Skrift; og idet han taler saaledes, taler han just ud af sin
subjective religiøse Tro (\textit{fides religiosa}).

Men ogsaa denne Tale er, naar vi adspørge den raisonnerende Forstand, udsat for
Misforstaaelse. Den Troende, der trodser paa det aabenbarede Ord, viser sig nu
som En, der er saa sikker paa det Historiske, fordi han i Grunden afleder sin
historiske Tro af sin religiøse Tro. Mangen En tænker da maaskee i sit stille
Sind: „Den Troende har Ret. Det har vel heller ikke stort at betyde med den
lange lærde Indledning, naar man blot af Hjertet kunde troe. Kunde jeg kun
troe, at Christus er opstanden, saa er han visselig ogsaa opstanden, opstanden
for mig. Men hvad nytter det at holde sig til Opstandelsens Evangelium, naar
denne Kjendsgjerning dog alligevel staaer som en fjern og fremmed Begivenhed,
der intet Indtryk gjør. Kunde Begivenheden blive mig nærværende, kunde jeg faae
et umiddelbart Indtryk af
% --- p. 17 ---
den opstandne Frelser, som Maria Magdalena, der hun mødte ham i Urtegaarden,
eller som Thomas, der han stak sin Finger i Naglegabet, og udbrød: min Herre og
min Gud! Men det er umuligt; med den ydre Kjendsgjerning kan man ikke begynde.
Det Historiske er et Forbigangent, og kan ikke mere blive nærværende: man maa
begynde med den subjective religiøse Tro.“

Saaledes dømmer den reflecterende Betragter, idet han ligesom bøier sig henimod
Troen uden dog at kunne gribe Troen. Og dog er ogsaa dette en Misforstaaelse.
Man kan virkelig ligesaalidt aflede den historiske Tro af den religiøse, som
omvendt denne af hiin. Naar vi eensidig begynde med den subjective Tro; holder
den Religiøse igjen, kaster sig til den modsatte Side, og lægger et uendeligt
Eftertryk paa det Historiske. „Det er Altsammen skeet, som skrevet staaer, og
efterdi det er skeet, saa troer jeg paa Kjendsgjerningen“. --- Tage vi saa dette
Historiske i og for sig, og lade den theologiske Kritik rokke ved
Kjendsgjerningen: saa stemmer den Religiøse igjen imod af al Magt, kaster sig
til den anden Side, beraaber sig paa Aandens indre Vidnesbyrd, og holder nu med
uendelig Lidenskab det Historiske fast i den subjective, religiøse Tro.

At ville aflede den religiøse Tro af den historiske, er en Misforstaaelse, der
fører til et catholsk Extrem; at priisgive det Historiske og holde sig til den
subjective religiøse Tro, er et moderne protestantisk Extrem.

Den catholske Kirke vil ved et guddommeligt Magtsprog forbyde al friere Kritik,
og kommer i Strid med Videnskaben. Ved sit Magtsprog viser den just, at den
frygter Videnskaben, og er bange for, at det ikke hænger rigtig sammen med det
Historiske. Denne Frygt vækker Mistanke; denne Mistanke undergraver Troen.

Den moderne Protestantisme frigiver Kritiken; Kritiken opløser det Historiske,
og undergraver Troen. Thi, hvis „den kritiske
% --- p. 18 ---
Indledning“ skal have mindste Indflydelse paa den religiøse Overbeviisning;
hvis man fra Kritikens Resultater tør uddrage en eneste Slutning for eller imod
Troen, da er det visselig den subjective Tro, idetmindste i dette Øieblik,
aldeles umuligt at troe religiøst paa den hellige Skrift.

Hvoraf kommer nu dette? Udentvivl deraf, at man saavel i Catholicismen som i
Protestantismen har gjort \emph{Christendommen eensartet med Videnskaben,
Evangelietroen eensartet med Theologien}.

Hvis Catholicismen var sig bevidst, at Christendom og Theologie ere
ueensartede; behøvede den slet ikke enten at forbyde eller indskrænke den frie
Kritik. Hvis Protestantismen tilfulde indsaae, at Evangelietroen er ueensartet
med Videnskaben, behøvede Kirken slet ikke at ændse den kritiske Opløsning.

Skuffelsen holder sig derved, at det Historiske og det subjectivt Religiøse
virke sammen i Troens Eenhed, uagtet det Historiske, som Historisk, dog
alligevel er ueensartet med det subjectivt Religiøse.

Den religiøse Tro kan ikke drage et eneste Aandedræt uden at støtte sig til sin
historiske Forudsætning. Dette har Catholicismen meget rigtig seet, og derfor
indskjærper den det Historiske. Men nu skal Kirkelæren tillige være Theologie,
og Theologien sikkre det brede historiske Grundlag (af Skrift og Tradition) ved
historiske Vidnesbyrd og historisk Lærdom, skal ved videnskabelig Forskning
hæve den objective, historiske Vished til absolut, uimodsigelig afgjørende
Vished. Det kan Theologien ikke, det er langt over dens Kræfter; saa tyer man
til det kirkelige Magtsprog, og Kirken tyranniserer Videnskaben.

Istedetfor det kirkelige Magtsprog sætter Protestantismen den frie subjective
Tro (\textit{fides religiosa}). Den protestantiske Bevidsthed har meget rigtig
indseet, at det Historiske er ueensartet med det Religiøse, at Alt dog tilsidst
kommer an paa Aanden, den
% --- p. 19 ---
levende, Menigheden iboende Aand, og ikke paa en større eller mindre Sum af
historiske, længst forbigangne Begivenheder. Imod denne almindelige
Grundsætning gjøre vi heller ingen almindelig Indvending; kun kommer det an
paa, hvorledes den i Gjerningen anvendes.

Den falske Anvendelse er, at man i Troens Navn og paa Troens Vegne adskiller
det, som Troen ikke kan og ikke tør adskille. --- Man tager det Historiske for
sig og gjør det til Gjenstand for en kritisk-videnskabelig Behandling. „Bibelen
er en Samling af forskjellige Oldtidsskrifter, skrevne til forskjellige Tider
af forskjellige Forfattere under forskjellige Vilkaar“. Saaledes begynder „den
kritiske Indledning“. Herimod er Intet at indvende, hvis „den kritiske
Indledning“, idet den begynder sit historisk-kritiske Cursus over disse
„Oldtidsskrifter“, saa strax vil erklære, at den videnskabelige Undersøgelse
Intet, aldeles Intet betyder for Evangelietroen og Kirkelæren. Men det vil „den
kritiske Indledning“ vistnok ikke erklære. Den lader sig meget mere forlyde
med, at disse Undersøgelser ere vigtige, meget vigtige, ikke blot --- hvad vi
gjerne indrømme --- for den lærde Videnskabs egen Skyld, men ogsaa vigtige, ja
i alle deres Enkeltheder og hele Vidtløftighed høist vigtige for alle Kirkens
tilkommende Lærere, for de vordende Præster. „Den kritiske Indledning“ vil
altsaa (middelbart eller umiddelbart) have en Stemme med i Afgjørelsen af,
hvorledes den religiøse Opfattelse af den hellige Skrift bør være.

„Den kritiske Indledning“ tænker som saa: lad os, med Reservation af en vis
almindelig Tro, først faae det reent Historiske kritisk afgjort, saa kunne vi
jo altid siden deraf uddrage Slutninger for vor religiøse Overbeviisning. Men
dette er en aldeles falsk Forudsætning. Loven er, at den historiske Tro Punkt
for Punkt skal hævdes og bæres af den inderlige, stærke, subjective, religiøse
Tro; Loven er, at den historiske Tro øieblikkelig gaaer tilgrunde, naar den
forlades af det Religiøse. Denne Lov
% --- p. 20 ---
bliver krænket af den kritiske Indledning. Beviset herfor er, at Kritikens
objective eller rettere: halv objective og halv subjective Afgjørelse løber ud
i endeløse Tvivl.

Den historiske Tro er i sit Væsen relativ, og gaaer kun paa det Sandsynlige;
men det blot Sandsynlige kan idelig omsættes og svækkes i det Usandsynlige. Den
subjective religiøse Tro er det Afgjørende, det Ubetingede, det Overgribende,
der griber det Absolute. Men i den religiøse Tale seer det alligevel ud, som om
de begge vare lige absolute; thi den Religiøse adskiller dem ikke: kun i
Reflexionen blive de adskilte.

Forholdet kan altsaa udtrykkes saaledes:

\emph{Den subjective religiøse Tro hviler paa det Historiske som paa sit
Grundlag, men bærer selv det Grundlag, hvorpaa den hviler.}

Efter saaledes at have bestemt Forholdet i Almindelighed, gaae vi da nærmere
ind i det Enkelte, og betragte f. Ex. \emph{Theologernes Strid om den
apostoliske Troesbekjendelse.}

Dersom man nu blot vilde nøies med at forholde sig christeligt til
Christendommen og videnskabeligt til Videnskaben: saa kunde Kirken strax for
sit Vedkommende ende Striden ved at indestaae for sin Troesbekjendelses
historiske Gyldighed; men nu skal man jo (om ikke for Andet, saa for
Theologiens Skyld) tillige forholde sig videnskabeligt til den apostoliske
Troesbekjendelse: nu begynder igjen den kritiske Indledning. --- „Kan man nu
ogsaa virkelig troe saa ubetinget paa dette Symbol, som Kirken antager? Er det
da ordret fra Christus selv eller fra hans Apostle? Har det ikke meget mere
dannet sig efterhaanden? Er ikke et Led optaget imod Gnostikerne, et andet imod
Apollinaristerne? Er „Nedfarten til Helvede“ ægte? Lad os sammenligne de
forskjellige Aftryk og Recensioner. Er den Recension, der findes hos Rufinus,
dog ikke mærkværdig afvigende fra andre tilsvarende Recensioner?“
% --- p. 21 ---
Nu reiser sig en Mængde Vanskeligheder, en Mængde Betænkeligheder; nu strides
her som paa Liv og Død (thi her strides jo om Kirkens Grundvold); nu dele
Theologerne sig i tvende Rækker; nu gaaer det: med Grunde --- for og imod,
Vidner og Vidnesbyrd --- for og imod, Fortolkninger over disse Vidner og disse
Vidnesbyrd --- for og imod, Fortolkninger igjen over disse Fortolkninger ---
for og imod. En videnskabelig Damp af Meninger og Raisonnementer indhyller de
Stridende! Det er en forviklet Strid; thi Theologien fortsætter jo ikke blot
Striden, men tillige Historien om Striden; hvorhos der da igjen kan opkomme
Strid om, hvorledes Stridens Historie egentlig bør fremstilles og bedømmes: det
gaaer i det Uendelige!

Hvad skal da nu Kirken gjøre ved al denne Strid? Skal den forbyde den ved et
Magtsprog? Det vilde jo være catholsk, og er i Protestantismen desuden aldeles
umuligt. Skal Kirken da maaskee vente sig nogen Fordeel af Stridens endelige
Afgjørelse? Nei ingenlunde! Hvis Kirken i noget Øieblik støtter sig paa et
videnskabeligt Forsvar, maa den ogsaa finde sig i at lide under et
videnskabeligt Angreb. Vil Kirken dele Theologiens Seire, maa den ogsaa finde
sig i at dele Theologiens Nederlag; men dette strider jo baade mod Kirken og
Christendommen. Derimod skal Kirken ganske rolig stole paa sin Myndighed, lade
Striden gaae sin videnskabelige Gang, og kun sørge for, at Videnskaben holder
sig paa sit Eget.

Den videnskabelige Kritik er paa ingen Maade istand til at fratage Kirken dens
Troesbekjendelse. Indvendingerne, naar de drives allervidest, kunne dog ikke
drive det videre end til at paavise en Usandsynlighed med Hensyn til Symbolets
Ægthed. Naar Kritiken har sagt: denne Troesbekjendelses Oprindelse fra Christus
og Apostlene er mig i høieste Grad usandsynlig: saa har den tvivlende,
modsigende Kritik sagt det sidste Ord, og kan ikke sige
% --- p. 22 ---
Mere. Men ved denne foregivne Usandsynlighed ere to Ting at bemærke:

a) Den høieste Grad af Usandsynlighed, der ellers er afgjørende for enhver
anden historisk-kritisk Prøvelse (naar \textit{fides historica} tages i og for
sig), vil i dette Tilfælde ikke slaae til. I Meddelelsen af Troesbekjendelsen
er det Overnaturlige jo medvirkende. Troesbekjendelsen har sin Oprindelse fra
Inspirationen, og viser hen paa Miraklet. Men Miraklet hører jo selv til det
Usandsynlige. At forkaste Symbolet paa Grund af en Usandsynlighed er altsaa en
Modsigelse: Miraklet vil idelig svække og forstyrre det kritiske Angreb.

b) Angrebskritiken bliver desuden svækket paa en anden Maade. De, der forsvare
Symbolet, finde jo ligesaa mange Grunde og Vidnesbyrd for deres Antagelse, som
de, der angribe det. Saa er Striden i en uophørlig Fluctuation af
Sandsynlighed og Usandsynlighed; her er ingen Afgjørelse, Afgjørelsen er
religiøs: Tro eller Vantro. Kirken, der taler Troens Sprog, gjør altsaa Ret i
at hævde Bekjendelsen uden at ændse den videnskabelige Kritik.

Stille vi nu den Troende midt imellem Kritikens Tvivl og Kirkens Bekjendelse,
hvorledes vil han da forholde sig? Naturligviis forskjelligt til det
Forskjellige: christeligt til Christendommen og videnskabeligt til Videnskaben.

Den Troende kan gjerne deeltage i Theologernes kritiske Symbolstrid, og
raisonnere for og imod, ligesom alle de andre Kritikere (thi idet han forholder
sig videnskabelig til Videnskaben, er han jo ogsaa Kritiker). Men hvori
adskiller han sig som Troende da fra Theologen? Theologen forventer jo, at der
endelig engang med Tiden skal udbringes et stort christeligt-kirkeligt Resultat
af den lange kritiske Undersøgelse; men det forventer den Troende nu slet ikke.
I Forhold til Kirkens Alvor er den hele vidtløftige Strid ham kun som en
menneskelig underholdende Tankeleg. Den christelige
% --- p. 23 ---
Afgjørelse er i Troen, ikke i Theologien. Dette vil ogsaa fremgaae af en
Betragtning over Stridens indre Beskaffenhed.

De, der angribe Symbolet, fremføre t. Ex. 3 Grunde, 3 Beviser imod det (derved
faaer Usandsynligheden foreløbig Overvægt). Saa fremtræde Forsvarerne med 4
Grunde, 4 Beviser for det (og i samme Øieblik stiger Sandsynligheden). Efter
nogen Tids Forløb udfinde Angriberne 7 Grunde og 7 Beviser (Usandsynligheden
har atter Overvægten) o. s. v.

Den Religiøse, der seer og forstaaer, hvad der kommer ud af slige
Forhandlinger, har jo dog fuldkommen Ret, naar han siger sig selv: skal jeg
vente, til denne Disput faaer Ende, saa faae jeg nok at vente til Dommedag; men
skal jeg bryde af ved den 20 Grund, fordi der følger en Modgrund, eller ved det
21 Beviis, fordi der følger et Modbeviis, saa kan jeg jo ligesaa godt bryde af
ved den første og eneste afgjørende Grund, den nemlig, der ligger i Kirkens
Vidnesbyrd.

Hvo er saa her den meest Intelligente, den Troende eller Theologen? Jeg skulde
mene den Troende, og trøster mig desuden til at godtgjøre det ved et
mathematisk Exempel. Der gives, som bekjendt, uendelige Brøker; at ville ved
fortsat Division udtømme en saadan uendelig Qvotient i endelige Tal, er Beviis
paa mathematisk Ukyndighed. Men saa naiv er „den kritiske Indledning“ i al sin
Skarpsindighed, at den mener at kunne udtømme en uendelig Qvotient i endelige
Tal.

Dog, der er endnu en Vanskelighed tilbage, ikke for Troen, ikke for Kirken, men
for den raisonnerende Reflexion, for Theologien.

% NB: the quotation opened here at „lad os kun indrømme is never closed in the
% 1850 printing — it simply stops at „udeladt?“ on p. 24. Left as printed.
„Lad os indrømme“, siger man, „lad os kun indrømme, at den apostoliske
Troesbekjendelse efter sit Indhold, det vil da sige: efter sit almindelige
Indhold, er baade apostolisk og grundchristelig; men saa kan da heller Ingen
negte, at dette guddommelige Grundord er udtrykt i menneskelige Ord,
articuleret i endelige Led og
% --- p. 24 ---
endelige Sætninger. Disse Led og Sætninger kunne da vel ikke Ord til andet være
aldeles uforbederlige eller uforanderlige. Hver Gang Symbolet oversættes fra
det ene Sprog i det andet, undergaae de enkelte Sætninger og Udtryk jo dog
altid nogen Forandring. Symbolet er desuden langt fra at indeholde alle Troens
vigtige Hovedsætninger, og et Led meer eller mindre kan da umuligt være
afgjørende for Ens evige Salighed. Eller skulde en troende Christen ikke kunne
blive salig, fordi han var døbt paa en Troesbekjendelse, hvoraf Nedfarten til
Helvede tilfældigviis var udeladt?

Saaledes raisonnerer den theologiske Kritik. Men, hvor megen Anseelse dette
Raisonnement end kan forskaffe sig i Videnskaben, saa er det dog i Forhold til
Kirken et aldeles spidsfindigt og sophistisk Raisonnement.

Naar den kritiserende Theolog spørger, om da et enkelt Led eller et enkelt Ord
i Bekjendelsen gjør Noget til eller fra: maae vi uvilkaarlig tænke paa
Sophisten, der spørger, om et enkelt afrevet Haar gjør skaldet, eller et enkelt
Korn en Dynge. Istedetfor at indlade sig paa Sophismer, bør Kirken for sit
Vedkommende øieblikkelig standse det endeløse Raisonnement, og føre den
principløse Undersøgelse tilbage til sit Princip. Principspørgsmaalet er ikke
dette, hvormange Bogstaver eller hvormange Sætninger eller hvormange Led der
kunne tænkes udeladte af Troesbekjendelsen uden Skade for Saligheden; men
Principspørgsmaalet er, om en kritisk Adskillelse af det Guddommelige og
Menneskelige i Kirkens Grundord kan skee uden Fare for Troen. Dette Spørgsmaal
er allerede besvaret med Nei. I Bekjendelsens Grundord er Guds Ord ikke til at
adskille fra det Historiske. --- Saa lyder Kirkens Sprog: „Ligesom Ingen kan
æde af Brød i Almindelighed, men af \emph{dette} Brød, \emph{dette} indviede
Brød, som vi bryde; og ligesom Ingen kan drikke af en Kalk i Almindelighed, men
af \emph{denne} Kalk, \emph{denne} Velsignelsens Kalk, som vi velsigne: saale-
% --- p. 25 ---
des kan og Ingen døbes paa en christelig Bekjendelse i Almindelighed, eller i
et religiøst Samfund i Almindelighed, eller af en Ordets Tjener i
Almindelighed; men dette er Grundordets afgjørende Myndighed, at Den, der
optages i Menighedens Skjød, maa døbes paa \emph{denne} Bekjendelse, i
\emph{dette} Samfund, i \emph{denne} Kirke, af \emph{denne} Præst, paa
\emph{denne} Dag. Dette er Grundordets afgjørende Myndighed, at det Evige her
ener sig med det Timelige, det Guddommelige med det Menneskelige i uopløselig
Bestemthed; ogsaa her gjælder det: „hvad Gud haver tilsammenføiet, skal intet
Menneske adskille.“ Imod dette Kirkens Sprog er det theologiske Raisonnement
forfængeligt og magtesløst! Hvilken Kritiker har vel Myndighed til at borttage
et eneste Led af Kirkens Bekjendelse? og hvilken Theolog har vel Myndighed til
at indsætte et eneste Led i denne Bekjendelse? --- Evangelietroen forbinder,
hvad Theologien adskiller, og hvorfor? \emph{Fordi Evangelietroen er ueensartet
med Theologien.}
% NB: the outer quotation „Ligesom Ingen kan æde … (opened p. 24) is closed by
% the same “ that closes the inner „hvad Gud haver tilsammenføiet …“; the
% printer set only one mark. Hence a net +1 „ over the batch.

%% ============================================================
%%  TREDIE FORELÆSNING (pp. 26--37 = PDF 39--50)
%% ============================================================
\chapter*{Tredie Forelæsning}
\addcontentsline{toc}{chapter}{Tredie Forelæsning}
\label{ch:forel3}
\markboth{Tredie Forelæsning}{}

\begin{center}\itshape
I Opfattelsen af Troesbekjendelsen og den hellige Skrift maa den
videnskabelige Theologie conseqvent vedblive at arbeide Troen imod;
thi Evangelietroen er ueensartet med Theologien.
\end{center}

% --- p. 26 ---
At ville paa een Gang, i eet og samme Nu, forholde sig troende
og dog videnskabelig til sin Troesbekjendelse, er en Modsigelse.
Forsaavidt jeg er i Bekjendelsen, er jeg udenfor Videnskaben;
forsaavidt jeg er i Videnskaben, er jeg udenfor Bekjendelsen. At
være stor i Bekjendelsen er noget ganske Andet, end at være stor og
udmærket i Videnskaben. Fiskeren fra Bethzaida var slet Intet i
Videnskaben; men Fiskeren fra Bethzaida blev dog stor i Bekjendelsen,
der han vidnede for Mesteren, og sagde: \emph{Du er Christus, den
levende Guds Søn.} Og Christus, der hører hans Bekjendelse, og veed,
hvad en Bekjendelse vil sige, svarer ham da heller ikke i
Videnskabens Sprog, sætter sig ikke i et kritisk-videnskabeligt
Forhold til ham, men styrker hans Tro, glæder sig over hans
Bekjendelse og priser ham salig: \emph{Salig er du Simon Jonas Søn;
thi Kjød og Blod haver ikke aabenbaret dig det; men min Fader, som
er i Himlene.}“
% PRINTER'S DEFECT: the closing “ after „Himlene.“ has no opening „
% anywhere in this quotation --- the printer set only the closing mark.
% Transcribed as printed; hence a net -1 „ over this batch.

Men --- det er mange Aar siden den første Christusbekjendelse blev
hørt i Verden; i disse mange, mange Aar er man saa efterhaanden
kommen i Vane med at forholde sig videnskabeligt til Christendommen
og theologisk-kritisk til Bekjendelsen.

% --- p. 27 ---
Naar da nu den kritiske Indledning tager Bekjendelsen for sig,
holder den op for Videnskabens Lys, og beseer den fra alle Sider;
faaer den (foruden de allerede anførte Betænkeligheder) endnu en
Skrupel, hvilken vi her saameget mere maa paaagte, som denne
videnskabelige Skrupel tillige synes at være en religiøs Skrupel.

Man raisonnerer omtrent saaledes:

„Tør den Troende paa ingen Maade adskille Bekjendelsens
guddommelige Indhold fra det menneskelige Ord og Bogstav, da tør
han heller ingen Anelse have om, at dette menneskelige Ord kan være
nogensomhelst Forandring underkastet. Man dømme nu om den kritiske
Undersøgelse, som man vil: Den, der eengang er bleven berørt af
Tvivlen, kan ikke vilkaarlig bryde af og vende tilbage til den
første Eenfoldighed. Faaer den Religiøse først at vide, at der
gives mere end eet Udkast, mere end een Recension af det apostoliske
Symbolum, og skal han desuagtet bindes til Bekjendelsens Ord og
Bogstav: da maa han jo leve i bestandig Angst for, at han endnu ikke
har fundet den rette oprindelige Text; hvo kan da fortænke ham i, at
han ved fortsat Efterforskning søger at komme til den størst mulige
Vished om det Oprindelige, at han t. Ex. interesserer sig for at
vide, om „Nedfarten til Helvede“ er oprindelig ægte eller en senere
Tilsætning. Sæt, at dette Led var uægte, sæt, at Christus ikke
nedfoer til Helvede, sæt at Forestillingen om denne Nedstigen i de
Dødes Rige blot var at forstaae billedligt om Christi virkelige Død;
saa kommer den Troende, der hænger sig i Bekjendelsens Ord, jo til
at troe paa en Usandhed, kommer i sin blinde Iver for at fastholde
det Apostoliske just til at antage en uapostolisk Forestilling.“ ---

Denne Kritikens Omhu og Bekymring for Kirkelærens Reenhed er, saa
smukt og tækkeligt den end ved første Øiekast kan tage sig ud, dog
alligevel af en meget spidsfindig og sophistisk Natur. Af Frygt for,
at her muligviis kunde have indsneget sig
% --- p. 28 ---
Noget i det Historiske (i \textit{fides historica}), som ikke var
reent og ægte historisk; af Frygt for, at dette Noget saa igjen skal
fordærve Kirken og skade Troen, vil den theologiske Kritik
godhedsfuld forføre baade Kirken og den Troende til at gjøre en
Adskillelse, der tilintetgjør Troen, nemlig en theoretisk
Adskillelse af det Guddommelige og Menneskelige i Bekjendelsens
Grundord. Theologen vil altsaa hjælpe den Troende! ja vist vil han
hjælpe ham; men Hjælpen er noget tvetydig, Hjælpen er en saadan, som
om han vilde hjælpe ham af Vandet i Ilden.

Faren, for at et uægte Led skulde, trods Kirkens Aarvaagenhed, være
indkommet i Troesbekjendelsen, er i religiøs Henseende Intet imod
den Fare, at gaae ind paa en raisonnerende Adskillelse af det
Uadskillelige, og derved tabe Troen.

I Tillid til Kirkens og Evangeliets Ord gaaer den Troende med sin
Bekjendelse da trøstig gjennem Verden, ind i Evigheden, ind for den
Dommer, der kjender alle Ting og dømmer retfærdig. Og, naar da
Menneskets Sjæl er stedt for den Dommer, der kjender Alt, baade det
Historiske og det Ikke-Historiske og det Overhistoriske: da skal det
vel vise sig, hvor Faren er, og hvori den bestaaer. Livsfaren er da
ikke den, at det her i Lysets Klarhed pludselig skulde opdages, at
Bekjendelsen ved en menneskelig Feiltagelse havde faaet et Led for
mange. Nei, Faren er en ganske anden, Livsfaren er, at Mennesket
ikke har levet i Troen paa sin Bekjendelse, som han burde. Har han
kun levet i Troen, stridt for Bekjendelsen og stridt den gode Strid:
da vil Dommeren vist ikke bebreide ham, at der er et Led for mange i
den Bekjendelse, hvormed han møder, men tage ham til Naade, og prise
ham salig, som han priste sin elskelige Discipel salig, der han
første Gang hørte hans Bekjendelse. Men vee den Sjæl, der møder i
Tvivl med tvivlsom Bekjendelse. „Den kritiske Indledning“ kan jo dog
ikke hjælpe ham, naar Dommeren kræver til Regnskab. Eller hvad skulde
det vel nytte, om denne
% --- p. 29 ---
tvivlende Sjæl nu begyndte at undskylde sig og sige: „Herre, der vare
saamange Udgaver og saamange Recensioner af vor kirkelige
Troesbekjendelse dernede, saamange indbyrdes afvigende og
forskjellige Recensioner, at jeg tilsidst ikke vidste, hvilken jeg
skulde vælge; Kritiken raadte mig da til at adskille det
Guddommelige fra det Menneskelige saavel i Din Bekjendelse som i Dit
Evangelium. Og saa fulgte jeg Kritikens Raad, og saa kom jeg ind i de
videnskabelige Undersøgelser, og ved mine videnskabelige
Undersøgelser kom jeg saa videre og videre ind i Tvivlen, og saa
sygnede min Tro, thi jeg vidste tilsidst --- desværre! ikke mere,
hvad jeg ret skulde troe.“ Mon da ikke Samvittighedsdommeren vil
svare denne Tvivler, og sige: „Du Daare, hvi troede Du ikke paa
Kirkens Vidnesbyrd, som det ordret lyder? hvi vilde Du ikke forstaae,
at mit Ord er uadskillelig Eenhed af guddommeligt og menneskeligt
Ord, som jeg selv er uadskillelig Eenhed af Gud og Menneske?“ ---

Vil Nogen absolut misforstaae denne Udtalelse, da give han blot Sagen
den Vending, som om jeg herved lagde an paa, at gjøre Videnskaben
forhadt og bandlyse Kritiken; thi det er sandelig en stor
Misforstaaelse. Det er saa langt fra, at jeg ved dette Forsøg har
havt til Hensigt at afskaffe Kritiken, at jeg tvertimod har villet
forhjælpe den theologiske Videnskab til en høist fornøden og aldeles
uundværlig Kritik, den Kritik, hvorved man adskiller den virkelige
Fare fra den indbildte.

Da Luther stred for sin Nadverlære, var den raisonnerende Kritik
allerede i fuld Gang, og vilde af al Magt gjøre Skilsmisse imellem
det Guddommelige og det Menneskelige i Indstiftelsens Ord. Men Luther
stred for sin Text, thi han stred i Troen. „Disse Indstiftelsens
Ord“ --- siger han --- „gjøre, at Brødet er hans (Christi) Legeme,
og Vinen hans Blod. Den, som nu æder dette Brød, han æder Christi
sande Legeme; hvo som drikker denne Kalk, drikker Christi sande Blod,
han være
% --- p. 30 ---
værdig eller uværdig. Det skal man troe fast; thi de kjære Christne
skulle gjøre Gud den Ære at bekjende, at hvad Gud siger, det kan han
og gjøre, som St. Paulus skriver om Abraham, at han har gjort (Rom.
4, 21). Hvo som vil være en Christen, han skal ikke gjøre som vore
Sværmere og Mytteriestiftere, der bekymre sig om, hvorledes det kunde
være, at Brødet er Christi Legeme, og Vinen Christi Blod. De ville
maale og begribe Gud med deres Fornuft, og, fordi det ikke rimer sig
med deres Fornuft, mene de, Gud kan heller ikke gjøre det. Men hvad
er det, at man allerede længe bekymrer sig derom? Om man ogsaa
sønderrev sig selv, saa vilde man dog ikke kunne begribe Gud med
menneskelig Fornuft; thi vor Herre Gud er ikke en saadan Gud, som
lader sig maale, begribe og fatte af menneskelig Fornuft, og hans
Gjerninger og Ord ere heller ikke saadanne Gjerninger og Ord, som
skulde være menneskelig Fornuft underkastede.“

„Lad os sætte, at vor Text og Forstand ogsaa var uvis og mørk, som
den ikke er, saavel som deres (Zwinglianernes Text og Forstand); saa
har Du dog den herlige trodsende Fordeel, at Du med god Samvittighed
kan staae paa vor Text, og sige saaledes: Skal og maa jeg da have en
uvis, mørk Text og Forstand, saa vil jeg hellere have den, som er
talt af den guddommelige Mund selv, end jeg vil have den, som er talt
af en menneskelig Mund; og skal jeg være bedragen, saa vil jeg
hellere være bedragen af Gud, om det var muligt, end af Mennesker;
thi bedrager Gud mig, saa vil han vel forsvare det og give mig
Opreisning. Menneskene kunne ikke give mig Opreisning, men føre mig
til Helvede. Saadan Trods kunne Sværmerne ikke have; thi de kunne
ikke sige: Jeg vil hellere staae paa den Text, som Zwingli og
Oekolampad med Tvedragt udsige, end paa den Text, som Christus selv
eendrægteligen udsiger.“

„Derfor kan Du glad sige til Christus baade i din Død og ved den
yderste Dom saaledes: Min kjære Herre Jesu Christe!
% --- p. 31 ---
Der har reist sig Klammerie om dine Ord i Nadveren: Nogle ville, at
de skulle forstaaes anderledes, end de lyde. Da de intet Vist lære
mig, men alene forvirre og gjøre mig uvis, og paa ingen Maade ville
eller kunne bevise deres Text: saa er jeg bleven ved din Text, som
Ordene lyde. Er der noget Mørkt deri, saa har Du villet have det
mørkt; thi Du har ingen anden Forklaring givet derover eller befalet
at give. Ligesaa finder man i intet Skrift eller Sprog, at er skal
hedde \emph{betyde}, eller \emph{mit Legeme Legemes Tegn}.“

„Skulde der nu være et Mørke deri, saa vil Du vel holde mig det
tilgode, at jeg ikke træffer det, ligesom Du holdt dine Apostle det
tilgode, at de ikke forstode Dig i mange Stykker; som da Du forkyndte
om din Lidelse og Opstandelse, og de dog beholdt Ordene, som de lode,
og gjorde ikke andre deraf. Ligesom ogsaa din kjære Moder ikke
forstod, da Du sagde til hende (Luk. 2, 49): Jeg maa være i det, som
er min Faders, og hun dog eenfoldeligen beholdt Ordene i sit Hjerte
og gjorde ikke andre deraf. Saaledes er jeg ogsaa bleven ved disse
Ord: \emph{Dette er mit Legeme, dette er mit Blod}.“

Hvilken Modsigelse det er, at ville paa een Gang forholde sig baade
christeligt og videnskabeligt til Christendommen, bliver os (om
muligt) endnu tydeligere, naar vi see hen til Theologernes forviklede
Strid om den Hellige Skrift. Lad os først betragte Striden i dens
Almindelighed, for desto bedre at kunne vurdere det Princip, der
udvikler sig i dens Enkeltheder.

At Kirken har sin hellige Skrift; at den betragter denne Skrift som
en af Gud indblæst, inspireret Skrift, saa at altsaa Skriftens
menneskelige Ord tillige ere Guds Ord, skal Ingen kunne negte, der
ikke ligefrem vil negte Kirkens og Christendommens Tilværelse.

Kirken har altsaa sin hellige Skrift, og sætter en Qvalitetsforskjel
imellem denne hellige, inspirerede Skrift og enhver anden
% --- p. 32 ---
Skrift. Inspirationen er Skriftens Qvalitet. Denne Qvalitet bliver
indskjærpet med saa stort Eftertryk af Refomatorerne, at Gyldigheden
% "Refomatorerne" so printed, for "Reformatorerne" --- printer's error, kept.
af vor evangeliske Kirkes Bekjendelsesskrifter endog ligefrem gjøres
afhængig af den hellige Skrift. Den augsburgske Confession, der jo
dog betragtes som vor Kirkes Palladium, er bygget paa Skriften; den
skal altsaa ikke gjælde i og for sig selv alene; men den skal gjælde
\textit{quia et quatenus}, \emph{fordi} og \emph{forsaavidt} den stemmer
med det i den hellige Skrift aabenbarede Ord.

Theologerne, der i Grunden kives om Alt, stride om Alt og tvivle om
Alt, ere naturligviis endnu i megen Tvivl og Strid om, hvorledes
dette \textit{quia} og \textit{quatenus}, dette \emph{fordi} og dette
\emph{forsaavidt}, egentlig skal forklares.

Betragte vi Forholdet fra Kirkens og Evangelietroens Standpunkt, er
Forklaringen let at udfinde. Naar den evangeliske Kirke bekjender om
sine Bekjendelsesskrifter, at de kun skulle gjælde, \emph{forsaavidt}
(\textit{quatenus}) de stemme overeens med den hellige Skrift: da har
det dermed ingenlunde været Kirkens Mening at priisgive disse
Bekjendelsesskrifter til Theologiens kritisk-videnskabelige
Behandling. Vi vide, hvorledes det gik Melanchthon, der forsøgte paa
at forandre den 10 Artikel om Nadveren. Melanchthon var sin Tids
meest anseete Theolog, ja den Mand, der selv havde udarbeidet
Confessionen; men den i Kirken levende Tro satte sig imod Theologen,
og Theologen maatte give efter. Kirken kunde ikke i noget Øieblik
tvivle om, at jo Bekjendelsen maatte være sand og tilforladelig
(hvorledes skulde et Kirkesamfund vel ogsaa kunne fremkomme med en
ifølge dets egen Tilstaaelse tvivlsom Bekjendelse); Kirken erklærer,
at Confessionen staaer fast, \emph{fordi} den stemmer med Skriften. I
dette \textit{quia}, i dette \emph{fordi} ligger den ubetingede
Vished; men i dette \textit{quia}, \emph{fordi}, udtrykkes med det
samme den gjennemgribende Qvalitetsforskjel, der er imellem den
hellige Skrift og Confessionen. Skriften er inspireret, det er
Confessionen ikke. Denne Forskjel gaaer ikke nærmest paa Indholdet
% --- p. 33 ---
(thi det væsentlige Indhold maa jo i Grunden være det samme i
Confessionen som i Skriften, efterdi Confessionen jo har sit Indhold
fra Skriften); men Qvalitetsforskjellen viser hen til Forbindelsen
imellem det guddommelige Indhold og det menneskelige Ord. I Skriften
er denne Forbindelse absolut uopløselig, thi Skriften er hellig og
inspireret; i Confessionen derimod er Forbindelsen ikke absolut: det
samme Indhold kan udtrykkes i andre maaskee mere adæqvate Ord og
Sætninger. Enkelte Led kunne føies til eller, for at gjøre Meningen
endnu tydeligere, udtrykkes med andre Ord, naar Kirken engang under
andre Tidsvilkaar skulde finde sig opfordret dertil. Derved faaer
Udtrykket: \textit{quatenus}, \emph{forsaavidt}, sin Betydning. Om
Forholdet mellem Skriftens Ord og Skriftens Indhold gjælder derimod
intet \emph{forsaavidt}, intet \textit{quatenus}. I Skriften selv er
Forbindelsen af det Guddommelige og det Menneskelige (Indhold og
Form, Mening og Udtryk) i alle Maader saa fuldkommen, saa absolut,
saa ufeilbar og uforanderlig, at ingen Adskillelse lader sig tænke:
Skriften er hellig og inspireret; hvert Ord i den kanoniske Skrift er
et Guds Ord, indblæst af Gud.

Men naar denne Opfattelse gjøres gjældende i sin hele Strenghed,
forstaaer dog vel Enhver, at det er en religiøs og ikke en
videnskabelig Opfattelse. Modsigelsen, den skjærende Modsigelse, der
ligger i, at ville paa een Gang forholde sig religiøst og dog
videnskabeligt til den hellige Skrift, fremlyser af den orthodoxe
Theologies egen Historie.

Der er i vor Kirkes Historie et Afsnit, man med Rette har kaldet den
bibliolatriske Periode. Bibliolatrie, afgudisk Bibeltilbedelse, er
ret et betegnende Udtryk for det usalige Forsøg, at ville paa een
Gang, i samme Henseende og i lige Grad, forholde sig baade
objectivt-videnskabeligt og subjectivt-religiøst til den hellige
Skrift.

For den objective videnskabelige Betragter er Bibelen en
% --- p. 34 ---
udvortes Gjenstand; for den Religiøse er den en hellig Gjenstand: for
den videnskabeligt Religiøse er Bibelen altsaa en udvortes hellig
Gjenstand, d. e. en Gjenstand for afgudisk Tilbedelse.

Saa sidder da den videnskabelige, men dog rettroende, Protestant, og
seer tilbedende paa sin Bibel, ligerviis som den videnskabelige, men
dog rettroende, Catholik seer tilbedende op til sin Pave. Det nytter
slet ikke, at en protestantisk Kritiker underkaster Paven en
videnskabelig Bedømmelse og i denne Bedømmelse blotter hans Feil og
Synder: den rettroende Catholik troer ligefuldt paa Paven, den
rettroende Catholik skal nok forsvare Pavens objective Ufeilbarhed
med objective Grunde; thi den rettroende Catholik tilbeder jo Paven.
Ligesaa lidt hjælper det, at en catholsk Kritiker --- til Fordeel for
den ene saliggjørende Kirke --- vil begynde at raisonnere mod Bibelen
og blotte dens menneskelige Ufuldkommenheder: den orthodoxe
Protestant afviser videnskabeligt enhver videnskabelig Indsigelse;
den orthodoxe Protestant skal nok vedblive at forsvare Bibelens
objective Ufeilbarhed med objective Grunde; thi den orthodoxe
Protestant tilbeder jo Bibelen.

I Protestantismens orthodoxe Theologie blev den objectivt
videnskabelige Inspirationstheorie udviklet med bibliolatrisk
Consequens. Bibelen var ikke blot inspireret, men
grammatisk-lingvistisk inspireret til Punkt og Prikke. Den Hellig
Aand havde udtrykkelig dicteret de hellige Mænd hvert Ord, hvert
Bogstav, de skulde skrive. Man kaldte derfor ogsaa de hellige Mænd
„Guds Skrivere“, „Guds Penne“ (\textit{notarii Dei, calami Dei}).

I Inspirationstheorien skulde altsaa Evangelietroen og Theologien
have fundet deres væsentlige og uopløselige Foreningspunkt; men see!
i Inspirationstheorien laae just Spiren til deres fremtidige
Adskillelse.

Idet Evangelietroen fastholder Skriftens Qvalitet i den
% --- p. 35 ---
ved Inspirationen bevirkede Eenhed af det Guddommelige og
Menneskelige, er den tillige villig til at anerkjende Forskjellen og
opfatte Skriftens menneskelige Side, som den faktisk viser sig. Troen
lægger slet ikke an paa, at begribe Forholdet, men takker Gud, at den
kan tilegne sig Indholdet: Troen forholder sig subjectivt til det
inspirerede Grundord.

Men den orthodoxe Theologie, der forholdt sig objectivt
videnskabeligt, vilde nu tillige uddrage videnskabelige Slutninger af
den religiøse Forudsætning, og foranledigede just derved sit eget
Fald.

Den religiøse Eenhed af det Guddommelige og Menneskelige i Skriftens
Ord opfattede den lærde Theologie da for det Første som en
\textit{grammatikalsk} Eenhed. Saaledes opkom Striden om
Vokalpunkternes Inspiration i det Gl. Tst. (Buxdorferne --- Ldv.
Capellus). I denne Strid led Orthodoxien et stort Nederlag.

Den religiøse Eenhed af Skriftens guddommelige Indhold og
menneskelige Form opfattede den lærde Theologie dernæst som en
lingvistisk-æsthetisk Eenhed. „Var det N. Tst. indblæst af den Hellig
Aand, da maatte det N. Tst's Græsk“ --- saa meente Theologerne ---
„ogsaa være det reneste og fuldkomneste Græsk“
% the print has no full stop before this closing “ --- as set.
Saaledes opkom Striden om Græciteten i det N. Tst. (Seb. Pfochen,
Siegm. Georgi --- Thom. Gataker, Salmasius, Sturz o. Fl.). I denne
Strid led Orthodoxien et stort Nederlag.

Med Hensyn paa den orthodoxe Inspirationstheorie maae vi altsaa
udtrykkelig bemærke: a) at den opløste sig i sine egne Conseqvenser;
b) at den svækkede Kirkens og Troens Selvstændighed; c) at den banede
Veien for den negative Kritik.

Da Kritiken først havde gjennembrudt Orthodoxiens theologiske
Ringmuur, rev den efterhaanden Kirkens øvrige religiøse
Forudsætninger ned, alt eftersom den fandt, at de stode den i Veien.
Under det Foregivende, at de historiske Kjendsgjerninger dog i
% --- p. 36 ---
Grunden vare uvæsentlige for Troen, forstod den kritiserende
Theologie ved mangehaande Vendinger at afspise den almindelige
Fromhedsfølelse, og kunde saa desto friere sætte sig i et reent
videnskabeligt Forhold til Skriften.

Den endnu moderate, halvvidenskabelige og halvconseqvente Theologie
er et maskeret Fritænkerie; det resolute Fritænkerie er Theologiens
videnskabelige Consequens.

I Theologiens Strid med den moderne Bevidsthed lide Theologerne Dag
for Dag det ene Nederlag efter det andet, komme dybere og dybere ind
i Vanskelighederne, dybere og dybere ind i Tvivlen.

Vistnok kan den besindige Theologie endnu slaae mangt et Angreb
tilbage, som en altfor hidsig Kritik kan have rettet mod Ægtheden af
et eller andet bibelsk Skrift; men hvad er saa herved vundet? Kan
Theologen nogensinde bringe det dertil, at han tør sige: „Ægtheden af
dette omtvistede Skrift er afgjort, nu endelig een Gang for alle
afgjort“? --- Ingenlunde! en saadan Paastand er uvidenskabelig, og
strider ligefrem mod Betingelserne for den historiske Tro
(\textit{fides historica}). Theologen bringer det i det Allerhøieste
kun til en Standsning, en foreløbig Standsning: „dette Skrifts
Uægthed er endnu ikke afgjort!“ Men, naar Ægtheden heller ikke er
ganske afgjort, saa er den jo dog underkastet Tvivl; men, er den
underkastet Tvivl, er den jo tabt for Troen; thi Troen gaaer ud, hvor
Tvivlen gaaer ind.

Saa stride nu Fritænkeren og Theologen --- under den samme
Forudsætning, under det samme videnskabelige Princip: kun Sproget og
Stemningen er forskjellig. Fritænkeren er resolut: „Uægtheden er
afgjort!“ Theologen er vankelmodig: „Uægtheden er endnu ikke
afgjort.“ Fritænkeren har i Videnskabens Navn ombyttet Tro med
Vantro. Theologen har i Videnskaben tabt Troen, og er trods al
videnskabelig Anstrengelse ikke mere istand til at hæve sin Tvivl til
Tro. Fritænkeren forholder sig
% --- p. 37 ---
ikke-religiøst til Skriften; Theologen kan aldrig bringe det dertil,
at han, som videnskabelig Theolog, forholder sig religiøst.

Vil Theologen holde sig til Kirken, maa han aabenlyst bryde med
Videnskaben; vil han holde sig til Videnskaben, maa han aabenlyst
bryde med Kirken, dette er sidste Dilemma.

Skal Kirken hævde sin Selvstændighed, maa den i lige Grad forvare sig
mod Theologens Forsvar og Fritænkerens Angreb. Det negative Vaaben,
hvormed Kirken værger sig, er den Grundtanke, \emph{at Christendommen
er høiere end Videnskaben, og Evangelietroen ueensartet med
Theologien.}

%% ============================================================
%%  FJERDE FORELÆSNING (pp. 38--51 = PDF 51--64)
%% ============================================================
\chapter*{Fjerde Forelæsning}
\addcontentsline{toc}{chapter}{Fjerde Forelæsning}
\label{ch:forel4}
\markboth{Fjerde Forelæsning}{}

\begin{center}\itshape
For den videnskabelige Theologie er Bibelen et historisk Oldtidsdokument,
hvis heterogene Bestanddele maae underkaste sig den raisonnerende Kritik;
for den troende Opfattelse er Bibelens vidunderlige Indhold ophøiet over
Kritiken, efterdi det er Betragteren umuligt at forholde sig
kritisk-raisonnerende til Miraklet: Evangelietroen er ueensartet med
Theologien.
\end{center}

% --- p. 38 ---
Bibelen, Kirkens hellige Skrift, er paa mange Maader bleven misforstaaet,
misbrugt og mishandlet af menneskelig Uforstandighed; men neppe er den dog i
Principet nogensinde før bleven saaledes misforstaaet, misbrugt og mishandlet,
som af den theologiske Kritik i den nyere Tid.

Mange Sværmere have mistydet Skriften til deres egen Fordærvelse; men de have
dog med alt dette forholdt sig til den, som til en hellig Skrift. Mange
Fritænkere have haanet Skriften --- til deres egen Fordærvelse; men de have dog
med alt dette forholdt sig til den, som til en hellig Skrift. Thi er det ikke
netop det Helliges Særkjende, at det frister Sværmeren og vækker Vantroens
Spot? --- Hvorfor tyer Sværmeren til Bibelen? Fordi han af denne „hellige
inspirerede“ Bog vil bevise os, at han selv er inspireret. --- Hvorfor spotter
Fritænkeren over Bibelen? Fordi det nu er ham til saa stor Forargelse, at denne
gamle Bog med sine „mange Myther“ og „hellige Eventyr“ endnu skal ansees og
gjælde for at være inspireret.

% --- p. 39 ---
I Forhold til Sværmeren og Fritænkeren bevarer Bibelen dog altsaa endnu sin
oprindelige Qvalitet; men ikke saaledes i Forhold til den kritiserende Theologie.

„Den kritiske Indledning“ vil naturligviis hverken opføre sig som Sværmeren
eller som Fritænkeren. „Den kritiske Indledning“ er en stikkelig Person, der i
lærd Forsigtighed lidt efter lidt har lært sig til at staae upartisk
\textit{above all party}, og mener saa at bevare den religiøse Ligevægt i den
videnskabelige Alsidighed. „Paa den ene Side kan man da“ --- saa hedder det jo
--- „vistnok ikke negte, at der i Bibelen findes mange uægte Bestanddele, mange
mythiske Sagn og overtroiske Fortællinger, der ligefrem maae henføres til
Datidens Folkephantasie og indskrænkede Forestillingskreds. Men paa den anden
Side kan man da heller ikke negte, at der i Bibelen findes saamange store og
herlige Tanker, saamange ophøiede Exempler og opløftende Charakteertræk, at
Bogen i det Hele maa vække den største Ærefrygt hos ethvert følende Gemyt, ja,
at selv den Videnskabelige maa tilstaae, at denne Bøgernes Bog dog vist, i det
Hele taget, ikke kan være skreven uden en særegen Guds Medvirkning.“ ---
Omtrent saaledes raisonerer „den kritiske Indledning.“

Men hvad er det saa for et Forhold, hvori „den kritiske Indledning“ da sætter
sig selv til Bibelen? Ja, det er et høist tvetydigt Forhold, et Forhold, der
paa een Gang skal være baade religiøst og dog videnskabeligt, og just derfor
hverken er det Ene eller det Andet.

Og hvad bliver saa Bibelen til i dette tvetydige Forhold? Den bliver en
tvetydig Bog, en Bog, der baade er inspireret og dog ikke inspireret, en
besynderlig Bog, der baade er det Ene og det Andet, og derfor i Grunden hverken
det Ene eller det Andet.

I denne Tvetydighed taber altsaa Bibelen sin Qvalitet; og derfor siger jeg, at
den theologiske Kritiks Forhold til Kirkens hellige Skrift --- til Trods for
Sværmeriets Forvildelser og Fritænkeriets
% --- p. 40 ---
Misgreb --- maa bestemmes som det egentlige principielle Misforhold, som det
Misforhold, der kommer Usandheden nærmest.

Den hellige Skrift er ingen objectiv videnskabelig Bog, og man kan da som Følge
heraf heller ikke forholde sig objectivt og videnskabeligt til den. Hvad
Bibelen er, kan ikke een Gang for alle afgjøres i og for sig; det maa hver Gang
afgjøres i og ved det Forhold, hvori Læseren sætter sig til Bibelen.

\emph{Som jeg selv er, saa er min Bibel, og som min Bibel er, saa er jeg selv.}

Forholder jeg mig nu troende til Bibelen, saa er Bibelen ogsaa for mig en
Aabenbaringernes Bog, en hellig inspireret Skrift, der viser mig, hvad jeg er,
og hvorledes jeg er, og kaster Lys over min hele Tilværelse. Min Forstand kan
nu vistnok opdage forunderlige Modsætninger, ja Modsigelser, i denne hellige
Bog. Det Ueensartede er her sat sammen til Eenhed: det udvortes Historiske og
det indvortes Sjælelige, det Oververdslige og det Verdslige, Almagten og
Skrøbeligheden, det Hellige og det Syndige findes her tilsammen i synderlig
Blanding og saare forunderlig Vexelvirkning. Det kan forarge den Vantroe, men
er dog alligevel til sand Opbyggelse for den Troende. I den Troendes Liv er her
jo noget Tilsvarende: ogsaa her bliver jo det Ueensartede forbundet til Eenhed;
ogsaa her kommer jo det udvortes Historiske til det indvortes Sjælelige; ogsaa
her skal jo den Verdslige forholde sig til den Oververdslige, den Skrøbelige
til den Almægtige, den Syndige til den Hellige, Den, der er Støv og Aske, til
Ham, der er fra Evighed og indtil Evighed. Ogsaa i den Troendes Liv er der jo
Modsigelser. Forsaavidt den Troende nu lærer at forstaae de Modsigelser, der
ere i hans eget Liv, vil han ogsaa lære at forstaae de tilsyneladende
Modsigelser, der ere i den hellige inspirerede Skrift; og omvendt, forsaavidt
han lærer at forstaae Modsigelserne i Skriften, vil han med Aandens Bistand
ogsaa lære at forstaae Modsigelserne i sit
% --- p. 41 ---
eget Liv. Maatte hine hellige Mænd, endskjøndt de vare Guds Udvalgte, dog
vedblive at troe med Haab mod Haab; hvormeget mere skulde dette ikke gjælde om
den Religiøse, der dog ikke saaledes tør kalde sig Guds Udvalgte! Og, naar den
Religiøse seer, at der i denne Verden maa troes med Forstand mod Forstand; naar
han lærer at indsee, at denne Lov bliver strengere og strengere i Forhold til
som Troeslivet bliver mere udviklet og alsidig belyst af Reflexionen: skulde
han da ikke ogsaa forstaae, at hine hellige Mænd kun vandrede saaledes, som
Skrifterne vidne, fordi de vandrede i Tro og i Haab, fordi de vandrede som Guds
Udvalgte. Forholdet er altsaa dette: jo inderligere den Religiøse existerer i
Troen og forstaaer Troens Vilkaar, desto inderligere erkjender han, at Kirkens
hellige Skrift er en inspireret Skrift.

Blev end den Troende forøvrigt saa udmærket i menneskelig Lærdom og Indsigt, at
han, om det saa skulde være, kunde kritisere alle Verdens Vise og opkaste sig
til deres Læremester: aldrig vil han dog bringe det dertil, at han kritiserer
den hellige Skrift, og opkaster sig til dens Læremester; men det omvendte
Forhold vil tiltage meer og meer, at nemlig den hellige Skrift i stedse høiere
Forstand bliver hans Læremester, den Trøster og Tugtemester, der holder ud med
ham indtil Enden.

\emph{Som jeg selv er, saa er min Bibel, og som min Bibel er, saa er jeg selv.}

Den Troende kan nu foreløbig abstrahere fra sin Tro, gaae ind i Videnskaben, og
sætte sig i et objectivt, litterært, kritisk-videnskabeligt Forhold lil
% printer's defect: „lil“ for „til“ — the first glyph is a plain Fraktur l,
% distinct from the t of „til“ elsewhere on the page. Transcribed as printed.
Bibelen. Herimod er der vistnok fra Dannelsens Standpunkt Intet at indvende.
Skulde det være den historisk-philologiske Kritiker forbudt at forholde sig
kritisk-videnskabeligt til den hellige Skrift, saa maatte det jo ogsaa være
Philosophen forbudt at forholde sig speculativt-videnskabeligt til
Christendommen. Skulde Bibelkritiken ligefrem være forbudt, saa maatte
Religionsphilosophien jo ogsaa ligefrem være forbudt;
% --- p. 42 ---
men det følger slet ikke af det her opstillede Princip. Vi ivre ligesaa lidt
imod en sand videnskabelig Bibelkritik, som imod en sand videnskabelig
Religionsphilosophie. Der er Intet i Veien for, at Christendommen jo kan blive
Gjenstand for dialektiserende Speculation; men den speculative Philosoph maa da
vel vide, hvad det er, han gjør. Idet Philosophen sætter sig i et
dialektisk-speculativt Forhold til Christendommen, er han udenfor
Troesforholdet; thi Troesforholdet er qvalitativt forskjelligt fra det
speculative Vidensforhold. Som nu Forholdet forandrer sin Qvalitet og bliver et
ganske andet Forhold, forandrer ogsaa Christendommen sin Qvalitet, og bliver
til noget ganske Andet. For den Troende er Christendommen Guds underfulde
Foranstaltning til Synderes Frelse og Salighed; for den Speculerende er
Christendommen et storartet Phænomen, et Phænomen, hvori nye vidunderlige Ideer
have ladet sig tilsyne, og den Speculerende søger nu at udgrunde disse Ideers
væsentlige, metaphysiske Forhold til det hele System af Verdensideer, hvori vi
leve, aande og tænke. --- Som det gaaer med Religionsphilosophien, saaledes
ogsaa med Bibelkritiken. Det religiøse Troesforhold til den hellige Skrift er
qvalitativt forskjelligt fra det litterære-kritiske Forhold. Som nu Forholdet
forandrer sin Qvalitet, og bliver et ganske andet Forhold, forandrer ogsaa
Bibelen sin Qvalitet, og bliver en ganske anden Bog. For den Troende er Bibelen
en hellig inspireret Skrift; for den litterære Kritiker er den hverken meer
eller mindre end en Samling af religionshistoriske Oldtidsskrifter, et
Oldtidsdokument, der i Principet bliver at opfatte fra samme Synspunkt, som vi
t. Ex. opfatte Indianernes Veda, Persernes Zenda-Vesta, Muhamedanernes Koran.

Her frembyder sig nu en Mængde Opgaver for den lærde Videnskab: at sammenligne
Codices, Versioner og Læsemaader, at medele archæologiske og philologiske
% printer's defect: „medele“ for „meddele“, single d as printed.
Oplysninger o. s. v. At disse
% --- p. 43 ---
Opgaver ere af reent videnskabelig Art, sees deraf, at de kunne behandles af
Enhver uden Hensyn til Troesbekjendelse. Hvad vore lærde Kritikere og Philologer
i denne Henseende have udrettet, hvad en Wetstein, en Griesbach, en Lachmann,
en Gesenius, en Winer i denne Retning have præsteret, har sin Gyldighed og
videnskabelige Berettigelse uden Hensyn til det Troesstandpunkt, hvorpaa hine
Lærde have befundet sig. Om den kritiske Philolog tilhører den reformeerte
eller lutherske Kirke, om han er Orthodox eller Rationalist, ja om han
overhovedet er en Troende eller en Ikke-Troende, gjør her aldeles Intet til
Sagen. Den lærde Kritiker forsker jo ikke i Qvalitet af Troende, men i Qvalitet
af Vidende; her spørges kun om Talent og Kundskaber, ikke om Religiøsitet. Hvad
den lærde Forsker præsterer, skal prøves og bedømmes af de Sagkyndige i Skolen,
men ikke af de Troende i Menigheden.

Staaer det nu fast, at der gives et dobbelt Forhold til Bibelen, et religiøst og
et videnskabeligt, staaer det fast, at Bibelen som Følge heraf forandrer sin
Qvalitet med det forandrede Forhold: saa bliver det vel ogsaa nødvendigt at
besinde sig lidt nærmere paa denne Qvalitetsforandring, før man uden videre
betroer sig til den videnskabelige Veiledning, der nuomstunder tilbydes i „den
kritiske Indledning.“

Ifølge Bibelens dobbelte Qvalitet, som inspireret Skrift og litterært
Oldtidsdokument, maa Bibelkritiken nødvendigviis støde paa saadanne Opgaver,
der ligge paa Grændsen af det Religiøse og det Videnskabelige, og som vi derfor
indtil videre kunne kalde de blandede Opgaver.

Herhen regne vi: a) Afgjørelsen af, hvilke Bøger der henhøre til Kanon. Denne
Opgave har en religiøs Betydning; thi herpaa beroer det jo, hvormange Bøger man
skal antage for inspirerede: Spørgsmaalet om Kanon bliver altsaa et
Troesspørgsmaal. Men opfattet fra det historisk-kritiske Synspunkt
% --- p. 44 ---
bliver det tillige et lærd videnskabeligt Spørgsmaal, nemlig om Bøgernes Ægthed
(Authentie). Endvidere: b) Undersøgelsen om Bøgernes Troværdighed, forsaavidt
denne Troværdighed kan bygges paa Beretningernes indbyrdes Overeensstemmelse:
Spørgsmaalet om Harmonien imellem de 4 Evangelier er et blandet Spørgsmaal.
Endelig: c) Afgjørelsen af alle de Punkter, hvori de hellige Forfattere i
Henseende til det reent Historiske afvige fra Profanskribenterne.

Det er nu vistnok disse blandede Opgaver, der have foranlediget
Indledningskritiken til at sætte sig i hiint ovenfor omtalte
Tvetydighedsforhold til den hellige Skrift, et Forhold, hvorunder Troen
neutraliseres med Videnskaben. Men i dette Forhold kunne de blandede Opgaver
ikke løses; thi, hvis de blandede Opgaver løses i Tvetydighed, maa Løsningen jo
selv blive tvetydig.

Før man altsaa forsøger paa at løse disse Opgaver, maa man allerførst være enig
med sig selv om, hvorvidt det nu egentlig er under Troens, eller under
Videnskabens Forudsætning, at Løsningen skal foregaae. Thi, eftersom man vælger
det ene eller det andet af de qvalitative Forhold, vil ogsaa Løsningen faae et
aldeles forskjelligt Udfald.

I. \emph{Lad os nu for det Første vælge Troesforholdet.} Den religiøse Faktor
er altsaa den afgjørende, den historiske varierer. For at see, hvad heri
ligger, vælge vi t. Ex. Undersøgelsen om de fire Evangeliers Ægthed.

At den historiske Faktor (\textit{fides historica}) her varierer, vil Enhver
indrømme, der kjender lidt til Kritiken. Det er jo muligt, at vort græske
Matthæusevangelium ikke er skrevet af Matthæus, men af en anden med Evangelisten
selv i nærmere eller fjernere Berøring staaende Forfatter. Det er jo ligeledes
muligt, at Evangelisten Johannes ikke selv har redigeret det Evangelium, der
bærer hans Navn o. s. v. Her er nu Plads for en Mangfoldighed af Undersøgelser,
Hypotheser og Forklaringer.
% --- p. 45 ---
Den religiøse Faktor er derimod ikke varierende, men constant. For den
subjective Tro (\textit{fides religiosa}) staaer det urokkelig fast, at de fire
Evangelier (hvo der saa end har skrevet dem) ere inspirerede. Denne kirkelige
Forudsætning kan Troen ikke opgive, uden at den med det Samme opgiver sit
Princip.

Sætte vi nu begge Faktorer sammen, saa bliver det Betingede bestemt ved det
Ubetingede, det Relative ved det Absolute, og Opgavens Løsning kan da udtrykkes
saaledes: de fire inspirerede Evangelier ere skrevne af fire inspirerede
Forfattere; disse Forfattere benævnes efter Matthæus, Markus, Lukas og
Johannes. De græske Overskrifter \textit{Κατα Ματθαιον, Κατα Μαρκον} etc. passe
da her paa det Allernøiagtigste.

Den samme Fremgangsmaade bliver at anvende, naar Spørgsmaalet er om de fire
Evangeliers indbyrdes Overeensstemmelse (\textit{harmonia evangeliorum}).

Den historiske Faktor varierer. Der er imellem de fire Evangelier ikke blot en
mærkelig Overeensstemmelse, men tillige mærkelige Afvigelser og tilsyneladende
Modsigelser (Jesu Slægtregister hos Matthæus kan Kritiken t. Ex. ikke faae til
at stemme med Slægtregistret hos Lukas). Et lidet Ord, en lille Forandring, og
Harmonien vilde være der; men --- det lille Ord mangler, selv den heldigste
Forklaring er dog alligevel kun en subjectiv Formodning, og saa vexler igjen
Hypothese med Hypothese.

Den religiøse Faktor er constant. Som Beretningerne staae, skulle de blive
staaende med uforanderlig guddommelig Autoritet; her reflecteres ikke paa nogen
Uovereensstemmelse: den ene Beretning er for den Troende selv akkurat ligesaa
useilbar, som den anden.

Sættes nu begge Faktorer sammen i Evangelietroens Eenhed, kan Resultatet
udtrykkes saaledes: i de fire Evangelier er det guddommelige Indhold fremstillet
i en saadan menneskelig Form, at den væsentlige Overeensstemmelse imellem de
identiske Beretninger tillige har antaget Præget af tilsyneladende Afvigelser i
% --- p. 46 ---
enkelte uvæsentlige Punkter.\footnote{„Juristerne sige, at en
Capital-Forbrydelse sluger alle de mindre --- saaledes er det med Troen, dens
Absurditet sluger aldeles det Smaalige. Uovereensstemmelser, der ellers vilde
virke forstyrrende, forstyrre her ikke, og gjøre Intet til Sagen. Derimod gjør
det Meget til Sagen, om En gjennem en Smaaligheds Calcule vil bortlicitere
Troen til den høiest Bydende; thi det gjør saa meget til Sagen, at han aldrig
kommer til Troen.“ Johannes Climacus.} Men, da Evangelierne ere inspirerede, er
det den Troende umuligt, at udsondre det Guddommelige fra det Menneskelige;
begge Momenter maae altsaa vedblive at virke sammen, indtil Dagen kommer, der
opklarer Harmonien.

Hvad der allerede er sagt om Troesbekjendelsen, gjælder tillige om
Evangelierne, ja om alle de kanoniske Bøger i Kirkens hellige Skrift: Faren er
ikke den, at et enkelt Led eller en enkelt Fortælling skulde ved en menneskelig
Feiltagelse have indsneget sig; men Faren er, at der gjøres en kritisk
Adskillelse af det Uadskillelige, hvorved Inspirationen ophæves, og Troen
forsvinder.

Denne Løsning, der er saa fyldestgjørende for Evangelietroen, kan umulig være
fyldestgjørende for Theologien; thi Theologien er jo en Videnskab.

II. \emph{Vælge vi nu det videnskabelige Forhold,} og forsøge fra et reent
objectivt Synspunkt at løse de blandede Opgaver: tager Sagen strax en anden
Vending. Den religiøse Faktor falder da bort, og den historiske hjemfalder til
Kritiken. Her er nu intet Fast, intet Ubetinget: Mistanken og Tvivlen kjender
ingen Grændse. Forklaringen ender da med, at ikke et eneste Evangelium er ægte,
ja at disse Evangelier først have faaet deres nuværende Form halvandet eller to
hundrede Aar efter Christus.

Dette Resultat, der aldeles ingen Betydning har for den subjective Tro, er
derimod i kritisk Henseende afgjørende for den
% --- p. 47 ---
historiske Tro, og viser tydelig nok, hvorledes den historiske Faktor
forflygtiges, naar den først er traadt ud af den rette Forbindelse med den
religiøse. Men, naar den opløsende Kritik, efterat den ved denne
Fremgangsmaade har undergravet Aabenbaringstroens historiske Grundlag, saa
tillige mener at kunne undergrave Evangelitroen, er dette --- hvad vi allerede
have beviist\footnote{S. 20. S. 24. S. 28--29.} --- en stor Misforstaaelse. Den
Troende lader sig ikke forstyrre af den videnskabelige Erklæring, at de
bibelske Skrifter ikke kunne staae sig for Kritikens Domstol; han seer kun heri
en Bekræftelse paa den Sandhed, at Bibelens Indhold er et overvidenskabeligt
Indhold, et Indhold, hvortil han maa forholde sig religiøst.

\emph{Som jeg selv er, saa er min Bibel, og som min Bibel er, saa er jeg selv.}

Er Bibelen inspireret, maa den jo have sin Qvalitet i et Under; men har Bibelen
sin Qvalitet i et Under: da er det jo ogsaa med det Samme afgjort, at den i
denne Qvalitet ikke tillader nogen videnskabelig-kritisk Opfattelse.

Et videnskabelig-kritisk Forhold til et Under er Nonsens. Dette gjælder ikke
blot om et saadant indvortes og sjæleligt Under, som Inspirationen; men det
gjælder tillige om de mere udvortes og, saa at sige, mere i Øine faldende
Mirakler.

Miraklet forudsætter en Vexelvirkning af ueensartede Faktorer, af det
Overnaturlige og det Naturlige. Hvis denne Vexelvirkning er en Skuffelse, maa
Miraklet selv være en Skuffelse; men, hvis den er sand, maa Enhver dog let
kunne indsee, hvor umuligt det er, at sætte sig i et objectivt Forhold til et
virkeligt Mirakel, at blive Tilskuer ved et Mirakel, som ved en physisk
Begivenhed, eller forholde sig videnskabelig-kritisk til en Mirakelfortælling,
som til en anden historisk Beretning.

Betragte vi t. Ex. de Mirakler, der synes at gribe meest
% --- p. 48 ---
ind i det Sandselige, det Mirakel, hvorved Moses lader Kilden fremspringe i
Ørkenen, og det Mirakel, hvorved Christus bespiser de Hungrige i Ørkenen; er da
ikke Vexelforholdet af det Naturlige og det Overnaturlige kjendeligt nok?

Moses opløfter sin Stav og slaaer paa Klippen; det kan en Tilskuer see med
naturlige Øine, thi det gaaer jo naturligt til. Og Klippen „græder“, og Kilden
opvælder og strømmer, og Israeliterne komme at stille deres Tørst: det kan en
Tilskuer see med naturlige Øine, thi det gaaer jo naturligt til. Vil Tilskueren
ikke troe sine egne Øine, kan han endnu desuden opnaae en umiddelbar sandselig
Vished ved selv at bøie sig ned og drikke af Kilden. Men med alt dette er
Tilskueren dog ikke kommen Miraklet selv et Haarsbred nærmere. Her er i denne
Begivenhed Noget, som end ikke den meest videnskabelige Iagttager vilde kunne
opdage, og dog er det just dette Noget, hvorpaa det ved den hele Opfattelse
væsentlig kommer an, det er Aarsagsforbindelsen mellem Mosestaven og den
sprudlende Kilde. Betingelsen er naturlig; men den virkende Aarsag, hvad er
den: naturlig eller overnaturlig? Antager Tilskueren det Første, saa staaer
Begivenheden for hans Phantasie endnu kun som et blot \textit{mirabile};
antager han det Sidste, forvandler den sig til et \textit{miraculum}.

Hvilken af disse to Antagelser Iagttageren skal foretrække, kan ikke afgjøres
ved videnskabelig Kritik. Det virkelige Mirakel kan bringe Tilskueren paa et
Punkt, hvor han maa sige til sig selv: „skal jeg her troe mine Sandser, kan jeg
ikke troe min Forstand; skal jeg troe min Forstand, kan jeg ikke troe mine
Sandser.“ Men videre, end til dette Dilemma, kan Phænomenet ikke bringe nogen
Tilskuer: det objective Mirakel støder Betragteren fra sig, driver ham tilbage
i Inderlighed, og lader Afgjørelsen beroe paa den dømmende Subjectivitet.

Den hellige Skrift fortæller os nu ikke blot om de mange Undergjerninger, men
beskriver os tillige de Indtryk, disse for-
% --- p. 49 ---
skjellige Undergjerninger have gjort paa de forskjellige Øienvidner. Hvor
forskjellige ere ikke disse Indtryk! Nogle troe ubetinget, og frelses ved
Troen; men Troen er jo i Subjectiviteten. Andre studse og forfærdes; men
Forfærdelsens Yttring er jo betinget af Subjectiviteten. I Folkets Udraab: „en
stor Prophet er opvakt iblandt os!“ --- „Gud har i Naade seet til sit Folk!“
fornemme vi en Svæven imellem den blotte Studsen og den inderlige Tro; men
denne Svæven foregaaer jo i Subjectiviteten. Mængden er nysgjerrig, tankeløs,
fortabt i det Sandselige (saaledes de Mange, der kom til Moses, og drak af
Kilden i Ørkenen; saaledes de Mange, der kom til Christus, og aade af Brødene i
Ørkenen); men denne overfladiske Nysgjerrighed, denne Tankeløshed, denne
dyriske Sløvhed er jo en sørgelig Mangel ved Subjectiviteten. Nogle forarges;
af Forargelsens Lidenskab udspringer saa igjen Forargelsens Forklaring (som da
Pharisæerne sagde om Christus: „han uddriver Diævle ved Beelzebul“); men
Forargelsen er jo i Subjectiviteten.

Et virkeligt Mirakel er virkelig ogsaa til at gaae fra Forstanden over. Men
Menneskene gaae fra Forstanden paa to Maader, deels saaledes, at de ligefrem
blive vanvittige, deels omvendt saaledes, at de sætte Forstanden til paa et
Høiere og just derved undgaae Vanvid. Dette er maaskee endog, efter Bibelens
Mening, det allernøiagtigste Udtryk for Miraklets guddommelige Øiemed, at
nemlig Menneskene derved skulle bringes til at gaae fra Forstanden, uden dog
ligefrem at tabe Forstanden.

Den Troende, der troer Miraklet, gaaer jo virkelig fra Forstanden, forsaavidt
han troer --- mod Forstand. At han desuagtet bevarer sin sunde Forstand, kan
alene sees deraf, at han, hvad de blot forstandige Ting ellers angaaer, taler
og handler aldeles forstandig. Den Forargede, der forarges paa Miraklet, gaaer
ligeledes fra Forstanden, forsaavidt han i en Art lidenskabelig Afsindighed
forsøger at bortforklare det Uforklarlige; men, da den af-
% --- p. 50 ---
sindige Forklaring dog altid yder ham en vis subjectiv Lindring, bevarer han
alligevel Forstanden i alle blot forstandige Anliggender, og bliver følgelig
ikke vanvittig. Pharisæernes Forklaring af Christi Mirakel var vistnok en i sin
Art afsindig Forklaring; men Forargelsen gav sig Luft i den subjective
Lidenskabs subjective Udbrud, og de Forargede beholdt deres naturlige
Forstand.\footnote{Hvis de Mirakler, om hvilke Bibelen fortæller os, virkelig
gjentoge sig i vore Dage, vilde Theologerne --- ifølge Principet --- uden Tvivl
være iblandt de Første, der stode for Tour til at blive ligefrem vanvittige.
Denne Formation af Mirakelvirkningen er aldeles ubekjendt i den hellige Skrift.
Skriften kjender kun de subjective: \emph{Tro og Forargelse}, men Vanvidet
kjender den ikke. At nu Bibelen ikke kjender Saadanne, som ligefrem bleve
vanvittige over Miraklet, maa vel forklares deraf, at den overhovedet slet ikke
kjender Saadanne, som have forsøgt paa at forholde sig
videnskabelig-kritisk til Miraklet, for at faae Miraklet videnskabelig-kritisk
retfærdiggjort.}

Som det nu forholder sig med det, at blive Øienvidne til et Mirakel, forholder
det sig ogsaa med det, at skulle fortælle og beskrive et Mirakel. Miraklet kan
kun fattes og tilegnes i Troen; det kan kun fortælles, beskrives og rettelig
fremstilles af Saadanne, som staae paa Troens øverste Trin ɔ: af inspirerede
Forfattere.

I Mirakelfortællingerne ville da de to ueensartede Faktorer igjen vise sig i
deres Vexelforhold; Troens subjective Faktor vil igjen være den absolute, den
constante; den historiske Faktor den relative, den varierende.

Naar altsaa to Evangelister fortælle det samme Mirakel paa en noget forskjellig
Maade: vil Afvigelsen ikke træffe Underets religiøse Kjerne, men alene træffe
de phænomenale Biomstændigheder og psychologiske Nüanceringer.

Som det forholder sig med det, at skulle fortælle et Mira-
% --- p. 51 ---
kel, saa maa det jo vel ogsaa forholde sig med det, at skulle læse, opfatte og
tilegne sig en Mirakelfortælling.

\emph{Som jeg selv er, saa er min Bibel, og som min Bibel er, saa er jeg selv.}

Det er denne Hovedtanke, jeg efter Evne har stræbt at fastholde og udvikle i
mit Skrift: \emph{Evangelietroen og den moderne Bevidsthed}.

Der gives kun to rene, qvalitative Forhold til den hellige Skrift, den Troendes
og Fritænkerens.\footnote{En vistnok rettroende Præst og meget anerkjendt
Conventtaler har, som jeg seer, været i Færd med at udbrede den Opinion, at jeg
i mit Skrift: „Evangelietroen og den moderne Bevidsthed“ skulde være paa Veie
til at etablere tvende hinanden modsigende og dog lige berettigede Sandheder,
og at dette nok skulde vise sig, hvis denne min Anskuelse ellers fik Leilighed
til at klare sig ud. Med en saa flygtig og, som det forekommer mig, letsindig
Meningsyttring kan jeg naturligviis ikke videre indlade mig; kun tillader jeg
mig ærbødigst at bemærke, at hvis den ærede Taler havde, ikke blot kikket i
Bogens Fortale, men virkelig givet sig Stunder og læst Bogen selv: da vilde han
uden Tvivl have seet, hvorledes det ogsaa her til Slutningen gaaer den
underfundige „Caiphas“. Den Forklaring, Verdensaanden (jfr. VII. Jesus Christus
er Verdens Frelser) tilsidst giver over det moderne Dogma, at \emph{Alt maa
gaae naturlig til}, er, om end ikke just meget opbyggelig, saa dog idetmindste
tilstrækkelig til at udvise de Conseqvenser, hvori Fritænkerens formeentlige
\emph{Sandhed}, hvis den ellers faaer Leilighed til at klare sig ud, omsider
maa ende.} Den kritiserende Theologs Troesforhold er et uklart, tvetydigt og
ureent Forhold. Desaarsag har jeg i benævnte Skrift --- saa gjerne jeg end
ellers havde villet --- da heller ikke ret kunnet være saa heldig, at tale
Theologerne til Behag; thi det er virkelig min Overbeviisning, at
\emph{Christendommen er høiere end Videnskaben, og Evangelietroen ueensartet
med Theologien.}

%% ============================================================
%%  FEMTE FORELÆSNING (pp. 52--65 = PDF 65--78)
%% ============================================================
\chapter*{Femte Forelæsning}
\addcontentsline{toc}{chapter}{Femte Forelæsning}
\label{ch:forel5}
\markboth{Femte Forelæsning}{}

\begin{center}\itshape
Den theologiske Skriftfortolkning miskjender Aabenbaringens eiendommelige
Charakteer, og opløser Troens Indhold i Videnskabens Forklaring; den
religiøse Skriftfortolkning nedsætter Videnskaben til Middel for en
troende Tilegnelse af det aabenbarede Ord: Evangelietroen er ueensartet
med Theologien.
\end{center}

% --- p. 52 ---
Skal Bibelen, Kirkens hellige Skrift, ikke blive borte i Theologiens
„kritiske Indledning,“ maa Troesforholdet sættes ubetinget, og
Inspirationen hævdes.

Men selv da, naar Bibelen i Qvalitet af hellig, inspireret Skrift er
sluppen vel igjennem den theologiske Kritik; selv da, naar den er kommen
Læseren saa nær, at ingen lærd Hypothese kan trænge sig ind imellem det
læsende Øie og den forefundne Text; selv da, naar Bibelen ligger for os
som en opslagen Bog, hvori den Troende læser Aabenbaringens Ord; selv da
er her endnu en Vanskelighed at overvinde, en Vanskelighed, som, langt
fra at formindskes, tvertimod forøges ved Theologiens theologiske
Bistand: denne Vanskelighed ligger i \emph{Fortolkningen}.

Hvad nytter det, at en Troende læser Bibelen, Ord for Ord, naar han dog
ikke ret forstaaer, hvad det er, han læser! At læse uden at forstaae er
meningsløst; at læse sig ind i en Misforstaaelse er værre, end slet ikke
at læse: Bibelen maa altsaa tydes, forklares og udlægges, at dens
guddommelige Indhold
% --- p. 53 ---
kan blive tilbørlig forstaaet og rettelig anvendt: Bibelen maa fortolkes.

At fortolke et menneskeligt Skrift, et Skrift, hvori en menneskelig
overlegen Aand har nedlagt en Skat af originale Tanker, er allerede en
stor Kunst; endnu vanskeligere bliver Opgaven, naar slige Tanker og Ideer
ikke ere systematisk fremstillede i Videnskabens almeengyldige Form, men
gestalte sig i Phantastens Form, i en Mangfoldighed af Charakterer og
Situationer, i Livets Form, hvor Tanken farves i Lidenskaben, og Sproget
vexler med Stemningen.

Er det nu allerede en stor Kunst at leve sig tilbørlig ind i et originalt
menneskeligt Skrift og fortolke et saadant Skrift, som det bør fortolkes:
hvad skulle vi da sige om den Opgave, at leve sig tilbørlig ind i den
hellige inspirerede Skrift og fortolke denne Skrift, som den bør
fortolkes?

Maa det indrømmes, at Bibelen har en overmenneskelig og dog menneskelig
Oprindelse, et overmenneskeligt og dog menneskeligt Indhold; saa følger
ogsaa heraf, at Bibelens rette Fortolkning maa være ikke blot en
menneskelig Kunst, men tillige en overmenneskelig Gave.

Forsaavidt nu Bibelfortolkningen er en menneskelig Kunst, maa den kunne
uddannes ved Videnskab og bygges paa videnskabelige Regler; men,
forsaavidt den betinges af en overmenneskelig Gave, strækker Videnskaben
ikke til: Alt beroer paa „Aandens Vidnesbyrd.“

Er der noget Punkt i den christelige Erkjendelse, hvor Troen og
Videnskaben berøre hinanden, da maa det vel være Fortolkningen af den
hellige Skrift.

Det forholder sig ikke med Fortolkningsvidenskaben, som med
Indledningskritiken. „Den kritiske Indledning“ med alle dens Bryderier og
Qvaklerier kan den Troende kaste fra sig, hvad Øieblik han vil, disse
kunstige Hypotheser og hypothetiske Kunster
% --- p. 54 ---
vedkomme ham dog i Grunden slet ikke; Bibelen er nu eengang den, den er,
og som den er; vi have ingen anden hellig Skrift, saa gjælder Luthers
Trods: „Guds Ord det skal de lade staae, og dertil Utak have.“

„Den kritiske Indledning“ er da ogsaa først egentligt bleven uddannet i
det sidste Aarhundrede, det er en Videnskabsgreen, den nyere Theologie
fornemmelig har opelsket til egen videnskabelige Tidsfordriv. Anderledes
med Hermeneutiken og den Veiledning, Videnskaben herved tilbyder den
religiøse Skriftfortolker. Denne Veiledning tør Ingen forsmaae, der vil
gjøre sig fortrolig med Bibelens menneskelige Form og Udtryk.

Den Tro, der i dette Stykke foragter den menneskelige Veiledning, er
ingen sand og sund Evangelietro, men en selvraadig, egensindig,
vilkaarlig Tro, der let udarter til Sværmeri. „Dette Ene“ --- siger
Luther --- „ophøier den hellige Christenhed og Kirke, at man forstaaer
den hellige Skrift, dette Ene nedtrykker Christenheden og Kirken, at man
ikke forstaaer den hellige Skrift . . . . Alle andre Kunster og
Haandværker have deres Præceptores og Mestre, af hvilke man maa lade sig
belære efter Orden og Lov; --- kun den hellige Skrift og Guds Ord maa
underkaste sig Enhvers Hovmod, Indbildning, Modvillighed og Anmasselse,
maa lade sig mestre, dreie og tyde, som Enhver forstaaer og vil efter sit
Hoved; derfra saamange Sværmerpartier, Secter og Forargelser, saa Gud sig
forbarme.“

For altsaa at sætte en Dæmning for den subjective Vilkaarlighed, har man
--- og det med Rette --- behandlet Fortolkningskunsten videnskabeligt; og
Theologien har saa herpaa grundet sin „theologiske Hermeneutik.“

Men, alt som den videnskabelige Hermeneutik har organiseret sig
theologisk, har ogsaa her Bedraget indsneget sig, det theologiske Bedrag,
at det Christelige skulde være eensartet med det Videnskabelige, det
theologiske Bedrag, at man ved fortsat videnskabelig
% --- p. 55 ---
Anvendelse af videnskabelige Regler skulde trænge dybere og
stedse dybere ind i Skriftens guddommelige Indhold. En sørgelig
Vildfarelse! thi inden man veed et Ord deraf, har Videnskaben igjen faaet
Magt med Christendommen, og er i fuld Gang med at udhule Troens Kjerne.

Det er nemlig charakteristisk for Hermeneutiken, at den kan opstille
saadanne Regler, som maae billiges af Alle, der ville tænke over deres
Bibel, saadanne Fortolkningsregler, som vi fra Videnskabens almindelige
Synspunkt maae tillægge almindelig Gyldighed, uden at den dog derfor paa
nogen Maade seer sig istand til at beherske eller bestemme den særegne
Anvendelse, som den enkelte Fortolker i hvert enkelt Tilfælde kan og bør
gjøre af disse Regler.

\emph{Skriften} --- hedder det --- \emph{skal fortolke sig selv.} Denne
Fortolkningsregel blev allerede udtalt paa Rigsdagen i Speyer, 1529, og
er forsaavidt en ægte protestantisk Regel. Men, naar Regelen opstilles
som videnskabelig Forudsætning, kan man jo deraf uddrage de
allerforskjelligste, ja aldeles modstridende religiøse Slutninger. Den
orthodoxe Theolog, den heterodoxe, Rationalisten, den speculative
Theolog, Fritænkeren, de kunne allesammen enes om den protestantiske
Regel; men ved Regelens Anvendelse komme de jo dog alligevel til ganske
forskjellige Resultater!

\emph{Enkelthederne i Skriften} --- hedder det --- \emph{maae forklares
af Skriftens Heelhed: Skriftens Analogie er den øverste Regel for
Fortolkningen.}

En ægte protestantisk Regel! men, naar Regelen opstilles som
videnskabelig Præmis, uddrager jo Theologien heraf de meest modstridende
religiøse Slutninger. Den Orthodoxe, Rationalisten, den Speculative,
Fritænkeren, de kunne jo alle enes om Regelen i Almindelighed, og endda
udbringer jo Enhver sin særegne Anskuelse af Skriftens Heelhed.
% --- p. 56 ---
Denne Uenighed er ikke tilfældig eller blot forbigaaende, men dybt
begrundet i Sagens Natur. Opfattelsen af Skriftens Heelhed er netop
betinget ved Opfattelsen af de enkelte Skriftsteder i deres indbyrdes
Sammenhæng; men Opfattelsen af denne Sammenhæng er jo igjen betinget ved
det Eftertryk, man hver Gang vil lægge paa dette eller hiint enkelte
Skriftsted. Paa et Sted siger nu Christus: „Jeg og Faderen vi ere Eet;“
men paa et andet Sted: „Faderen er større end jeg.“ Paa nogle Steder
tales der i Skriften om den Hellig Aand, som om en personlig selvstændig
Aand (i Faderens, Sønnens og den Hellig Aands Navn); men paa andre Steder
benævnes Aanden blot som „Kraften fra det Høie,“ „den Høiestes Kraft.“
--- Paa nogle Steder lægges der en overordentlig Vægt paa den
menneskelige Frihed: „arbeider paa Eders Saliggjørelse med Frygt og
Bæven;“ men umiddelbart derpaa nedsættes saa igjen den menneskelige
Frihed og de menneskelige Gjerninger til et reent Intet, et Intet for
Gud: „Gud er selv den, der virker i Eder baade at ville og at udrette
efter sit Velbehag.“

Disse forskjellige Udsagn staae skarpt imod hinanden; det ene kan ikke
afledes af det andet, ikke indskrænkes ved det andet, ikke ligefrem
forklares af det andet; Forstanden kan ikke formidle dem, ethvert
Formidlingsforsøg har jo viist sig som en Indbildning, en Selvskuffelse.

Den religiøse Forstaaelse er altsaa i absolut Forstand afhængig af den
subjective Tro:

\emph{Som min Tro er, saa er min Skriftfortolkning, og som min
Skriftfortolkning er, saa er min Tro.}

Men den subjective Tro, der i Vexelforhold til Evangeliets Aand afgiver
det høieste og absolute Fortolkningsprincip, er aldeles ueensartet med
„Hermeneutikens“ objective Regler.

Ville nu Skriftfortolkerne desuagtet gjøre disse ueensartede Faktorer
eensartede ved at uddrage religiøse Slutninger af videnskabelige
% --- p. 57 ---
Præmisser, eller omvendt: videnskabelige Slutninger af
religiøse Præmisser: saa gjælder det da i Sandhed, som Christus siger, at
det er skjult for de Vise og Skriftkloge, hvad der er aabenbaret for de
Umyndige.

Vi behøve ikke at gaae langt, for at finde Exempler; de ligge os kun
altfor nær, og ere desværre kun altfor talrige. For blot at nævne eet
Exempel, vil jeg her minde om Theologen \textit{De Wette}.

Den for nylig afdøde Theolog, \textit{De Wette}, har visselig udmærkede
Fortjenester baade som Kritiker og som videnskabelig Skriftfortolker. Jeg
vilde kun vanære mig selv, hvis jeg her tillod mig at vanære hans Minde
ved at henregne ham blandt dem, der ringeagte Christendommen eller
behandle Skriften letsindig. Tvertimod! \textit{De Wette} var vistnok en
efter Nutidens Tænkemaade ligesaa alvorlig og retsindig, som flittig og
begavet Gransker. Men i religiøs Forstand --- det siger jeg uden
Forbeholdenhed --- i religiøs-ethisk Forstand er hans Aabenbaringsbegreb
aldeles forfeilet, og hans theologiske Granskning i Forhold hertil bygget
paa en Misforstaaelse.

Eller er det ikke en Misforstaaelse, naar \textit{De Wette} vil frakjende
Propheterne den Aandens Gave, som Christus og Apostlene tillægge dem? er
det ikke --- i religiøs Forstand --- en Misforstaaelse, naar han
fornemmelig vil have de prophetiske Forudsigelser begrundede paa svævende
Ønsker og ubestemte Forhaabninger?%
% printed footnote mark: *)
% PRINTER'S DEFECT: the opening „ before „Weil“ is never closed --- the
% footnote ends at „256.“ with no closing “ in the 1850 printing.
\footnote{„Weil die Idee in den Propheten durchaus vorherrscht, so sind
ihre Vorhersagungen zum Theil nur als Hoffnungen und Wünsche, . . . und
sind fast immer unbestimmt und schwebend. Lehrbuch der
historisch-kritischen Einleitung in die Bibel. 1. S. 256.}
Eller er det ikke en Misforstaaelse, naar \textit{De Wette} tillader sig
at kritisere Propheten Ezechiel, som man kritiserer en middelmaadig
Forfatter i den moderne Litteratur? naar han t.
% --- p. 58 ---
Ex. dadler denne Prophets Mangel paa Dybde og Aandsrigdom og paa store
Tanker? naar han beklager, at Ezechiels prophetiske Stiil er for lav og
for bred og nedsunken til en mat Prosa? naar han opholder sig over, at
Propheten ved allegoriske og symbolske Digtninger jo vistnok vil hæve sig
over det Sædvanlige, men dog sædvanlig forfalder til det Overlæssede, det
Søgte, det Forvirrede?%
% printed footnote mark: *)
\footnote{Was bei ihm (Ezechiel) am meisten auffällt, ist seine
levitische Gesinnung, vermöge deren er auf heilige Gebräuche einen hohen
Werth legt, und sich selbst in Bildung von Idealen nicht darüber erheben
kann. Daher auch sein Mangel an Tiefe und Reichthum des Geistes und an
großen Gedanken . . . . Die prophetische Rede ist bei ihm zur niedrigen,
weitschweifigen, matten Prosa herabgesunken; nur in symbolischen und
allegorischen Dichtungen will er sich über das Gemeine erheben, verfällt
aber gewöhnlich in das Ueberladene, Gesuchte und Verworrene. A. Skr. S.
282---83.}
Hvilken videnskabelig Overlegenhed, hvilken kritisk Selvtillid! Kan en
Theolog da saaledes mestre en Prophet?

Det er en sørgelig Misforstaaelse, foranlediget derved, at man vil
forholde sig videnskabelig-religiøst til Guds Ord, og fortolke de
inspirerede Forfattere som andre geniale Skribenter. Fortolkningen
forvandler sig til Kritik: man ophøier den ene Prophet og nedsætter den
anden. Og dog skal Aabenbaringsbegrebet hermed være ukrænket, og
Propheterne i en vis Maade allesammen ansees for inspirerede. Hvilken
Modsigelse, hvilken ukritisk Forvirring!

Ere Propheterne ikke inspirerede, tør Hermeneutiken og Kritiken virkelig
indestaae os for, at det egentlig slet ikke har Noget at betyde med
Bibelens Inspiration: hvorfor sige Theologerne det da ikke reent ud, og
give Fritænkeren Ret, som frie videnskabelige Tænkere? Men er Bibelen
inspireret, ere Propheterne tilsammen inspirerede af Guds Aand: hvorledes
kan da den theologiske Kritik være saa ukritisk, at anlægge den samme
Maalestok paa inspirerede Forfattere, som paa ikke inspirerede?
% --- p. 59 ---
Ere Propheterne inspirerede: da maa jo den Forskjel, der med Hensyn til
Stiil og Fremstilling muligviis kan findes imellem en Jesajas og en
Ezechiel i Forhold til Inspirationen selv være en aldeles uvæsentlig, en
forsvindende Forskjel.

Den theologiske Videnskab, der bilder sig ind, at tale med Alvor og
Værdighed om det Skjønne, det Høie, det Dybe, der formeentlig skal findes
hos den ægte Jesajas i Modsætning til Propheterne af en lavere
Skrivemaade, krænker i religiøs Forstand ligesaa vist Jesajas ved sin
formeentlige Roes, som den krænker en Ezechiel eller en Daniel ved sin
anmassende Dadel.

Propheten har slet ikke sin Qvalitet i Stilens æsthetiske Fortrin, men
Propheten har sin Qvalitet deri, at han er inspireret af Herrens Aand.
Man skal troe Prophetens Ord og gjøre derefter; men man skal ikke
raisonnere over hans Stiil eller strax erklære ham for uægte, fordi man
ikke videnskabeligt kan finde sig til Rette med hans Skrivemaade.

For nu endelig at fjerne denne theologiske Misforstaaelse med den deraf
flydende Forvirring, maa man i Hermeneutiken, som i Indledningskritiken,
skarpt adskille den videnskabelige Opfattelse fra den troende Tilegnelse,
og derhos tillige vedgaae, at disse Opfattelser ere saa ueensartede, at
ingen ligefrem Overgang kan gjøres fra den ene til den anden.

Naar Fortolkeren blot forholder sig videnskabelig til Bibelen, bliver
Bibelen philologisk og historisk at behandle som enhver anden
Religionsbog fra den fjerne Old. Dette Fortolkerens videnskabelige
Forhold til Kirkens hellige Skrift er paa det Allerbestemteste bleven
fremhævet af L.\ I.\ Rückert i Fortalen til hans Commentar over Pauli
Brev til Romerne, hvor det blandt Andet hedder: „Fortolkeren af det
N.\ T. . . . har, som saadan, slet intet theologisk System og tør intet
have, hverken et dogmatisk eller et Følelsessystem, han er, forsaavidt
han er Exeget, hverken orthodox eller heterodox, hverken Supranaturalist
eller
% --- p. 60 ---
Rationalist, eller Pantheist, eller hvilkesomhelst --- ister, der ellers
maatte gives; han er hverken from eller ugudelig, hverken sædelig eller
usædelig, hverken ømtfølende eller følesløs; thi han har blot den Pligt
at udforske, hvad hans Forfatter siger, for at overgive dette som reent
Resultat til Philosophen, Dogmatikeren, Moralisten, Asketen o.\ s.\ v.
--- Saa tør han da som Exeget heller ikke have anden Interesse end denne
ene, at forstaae Apostelen rigtig, opfatte hans Tanker reent uden fremmed
Indblanding, og forelægge dem ligesaa rene og ublandede for Læseren.
. . . For ham (--- den videnskabelige Exeget ---) maa det være
ligegyldigt, om Paulus taler Sandhed eller Løgn, om det er en sædelig
eller en usædelig Aand, der gaaer igjennem hans Breve, om hans Lære er
heldbringende eller grundfordærvelig; thi han (Fortolkeren) skal kun
fremstille, hvad han (ɔ: Forfatteren) siger, hvilken Aand der
gjennemtrænger ham, hvilken Lære han fører.“

Forud for Rückert (hvis Commentar over Romerbrevet udkom i Leipzig 1831)
havde allerede adskillige ældre Hermeneuter (G.\ L.\ Bauer,
Bretschneider, Griesbach o.\ Fl.) indskjærpet den videnskabelige
Forudsætningsløshed, som burde udmærke den fordomsfrie Bibelfortolker.
„Vilde man“ --- hedder det --- „forud antage det N.\ T. for Aabenbaring,
og ifølge deraf kun antage den Mening for den hermeneutisksande, som
synes at være en Aabenbaring værdig, saa vilde man bevæge sig i en
Cirkel; --- med lige Ret vilde man da udlægge ethvert andet
Aabenbarings-Dokument, som Koranen eller Zendbøgerne, paa denne Maade, og
derved kunne godtgjøre dets Indhold som logisk rigtigt og sandt.“ ---

Saaledes vare da Theologerne selv i deres Iver for en uhildet
videnskabelig Skriftfortolkning paa gode Veie til at opgive det
Theologiske i Theologien. Men, da Theologien jo netop bestaaer i den
Modsigelse, at skulle gjøre det Ueensartede eensartet, da Theologien ikke
kan opgive det Religiøse uden at opgive
% --- p. 61 ---
sin Charakteer, saa fulgte det af sig selv, at hiin videnskabelige
Conseqvens heller ikke lod sig gjennemføre.

Istedetfor at fastholde Principet og lade det Videnskabelige afklare sig
i Videnskaben, begyndte man naturligviis igjen med de sædvanlige
Raisonnementer om, ikke at gaae forvidt, forsøgte saa igjen at dække sig
bag Troeshensynet og gjøre nogle Indrømmelser til begge Sider, at
Theologien dog nogenlunde kunde bevare et theologisk Anstrøg.

Man kan vel ikke fortænke Theologien i, at den just ikke føler nogen
synderlig Drift til at slaae sig selv ihjel. Kan Theologien ikke bestaae
uden i Halvhed, Inconseqvens og Principløshed: hvad Under da, at de
besindige Theologer betragte Conseqvensen som en Overdrivelse og det
ubetingede Princip som en Eensidighed.

Saaledes skete det da, at mange ansete og fromme Theologer (hvoriblandt
Storr i Tübingen, Seiler i Erlangen, Stäudlin i Göttingen, Stein o.\ Fl.)
paa en ligesaa forsigtig som lemfældig Maade bestræbte sig for at kalde
Troens religiøse Forudsætning tillive igjen, og med stor Varsomhed
begyndte at mægle imellem „den historiske og den religiøse Interesse“.

Det første videnskabelige Forsøg paa at gjenoprette Forbundet imellem
Hermeneutiken og Troeslæren skyldes dog egentlig F.\ Lücke (nu Prof.\ i
Göttingen). Som Privatdocent i Berlin udgav Lücke et --- som vi indrømme
--- „til mange Sider opvækkende Ungdomsskrift (Grundriß der
neutestamentlichen Hermeneutik u.\ ihrer Geschichte, 1817);“ en nærmere
Udvikling af Begrebet om theologisk Hermeneutik findes af samme Forfatter
i „Studien und Kritiken“ for A.\ 1830. Opgaven for den bibelske
Hermeneutik bliver da bestemt derhen: „\emph{at construere de
hermeneutiske Principer saaledes, at det eiendommelige theologiske Moment
kan paa virkelig organisk Maade forenes dermed, og ligeledes, at danne og
fastsætte
% --- p. 62 ---
det theologiske Moment saaledes, at de almindelige Udlægningsprinciper
beholde deres fulde Gyldighed.}“

Hvilken Misviisning, denne for Resten velmeente Bestræbelse underligger,
kan vel ikke tilstrækkelig oplyses i denne Time, men vil dog vise sig
efterhaanden, som vi rykke videre frem i disse Foredrag. Enhver, der har
noget Indtryk af, hvad et videnskabeligt Fortolkningsprincip vil sige,
behøver blot at gjennemgaae de Commentarer, der hovedsagelig ere anlagte
og udførte efter ovenstaaende Grundsætning, jeg mener Commentarerne af
Lücke selv, saavel som af Tholuck og Olshausen.

Der er i disse Commentarer en dobbelt Mislighed, der rammer deels
Methoden, deels Principet.

Hvad Methoden angaaer, er det saa langt fra, at det religiøse Element her
skulde være sammenarbeidet til „organisk“ Eenhed med det videnskabelige,
at det Religiøse meget mere virker forstyrrende paa det Videnskabelige
og, omvendt, det Videnskabelige paa det Religiøse.

Med al den Lærdoms Vidtløftighed, der udmærker disse halvopbyggelige
Commentarer, mærker man snart, at det philologisk-historiske Apparat
optager en saadan Plads, at den religiøs-ethiske Udvikling ikke ret
trænger igjennem, men sædvanligviis maa nøies med nogle almindelige
Sætninger og indstrøede Bemærkninger. Undertiden indtræder (især hos
Olshausen) det omvendte Forhold, at nemlig de religiøse Forklaringer og
dogmatiske Excurser brede sig med en saadan Udførlighed, at det
Philologiske og Historiske lider derunder. Grundfeilen er, at de
ueensartede Opgaver neutralisere hinanden, at Opfattelsen bliver
tvetydig, og Interessen tvivlsom. Dette viser sig endnu tydeligere, naar
vi see hen til Principet selv.

I de nævnte Commentarer krydse de religiøse og de videnskabelige Opgaver
overhovedet hverandre saaledes, at det videnskabelige Princip idelig
svækker det religiøse. Dette sees ikke blot i
% --- p. 63 ---
Behandlingen af Miraklerne, hvor Fortolkeren nu paa saa mangfoldige
Maader maa komme i Forlegenhed med sine halvvidenskabelige Forklaringer
(hvorledes er t.\ Ex.\ Olshausen ikke kommen i Forlegenhed med sine
Fortolkninger af Sygehelbredelserne i Galilæa og sin halvmoderne Theorie
om Dæmonerne, en Forlegenhed, hvoraf Strauß har vidst at drage den
kritiske Fordeel); men det gjælder ogsaa om Fortolkningen af Christi Ord,
især i saadanne, som udtrykkelig ere bestemte til at ydmyge den naturlige
Forstand. For blot at minde om et enkelt Sted, lyder jo den haarde Tale
(Joh.\ 6): „Dersom I ikke æde Menneskens Søns Kjød og drikke hans Blod,
have I ikke Livet i Eder. Hvo, som æder mit Kjød og drikker mit Blod,
haver det evige Liv, og jeg skal opreise ham paa den yderste Dag.“

Slige Skriftsteder kunne ikke fortolkes videnskabeligt%
% printed footnote mark: *)
\footnote{Jfr.\ Evangelietroen og den moderne Bevidsthed S.\ 432---36.},
det er aldeles umuligt. For den videnskabelige Betragtning tabe de deres
Braad og dermed deres Kraft, og forvandle sig til blot figurlige
Talemaader. Istedetfor at indprænte sig „den haarde Tales“ religiøse
Charakteer og gjøre Troens Opgave kjendelig, geraader den theologiske
Fortolker ind i et Virvar af modstridende Forklaringer; og, naar han saa
føler, at han med alle disse Forklaringer dog i Grunden Intet kan
forklare, overfaldes han tilsidst af den Tvivl, at Johannesevangeliet
maaskee ikke engang kan holde sig.

I den videnskabelige Fortolkning gjælder det videnskabelige Forhold, og
derfor bør ogsaa Fortolkeren her være saa conseqvent, at forholde sig
videnskabeligt til sin Bibel. Det er dette objectivt-videnskabelige
Forhold, der saa klart og bestemt er blevet udtalt af L.\ I.\ Rückert.

I den religiøs-ethiske Fortolkning gjælder det religiøse Forhold, og
desaarsag bør ogsaa Fortolkeren her være saa conseqvent,
% --- p. 64 ---
at forholde sig religiøs-ethisk til den inspirerede Skrift. Det er denne
Conseqvens, jeg overalt har ladet den Troende fastholde i
„Evangelietroen og den moderne Bevidsthed.“

Hermed ere Sphærerne adskilte: Videnskaben og Troen have nu hver sit
Princip og bevæge sig hver om sit Middelpunkt.

De excentriske Kredse kunne nu forøvrigt berøre hinanden, gribe ind i
hinanden, ja belyse hinanden gjensidig, naar man kun ikke glemmer, at
Principerne ere ueensartede.

At der gives et Vexelforhold mellem den christelige og den humane
Dannelse, kan ikke negtes; men derfor maa den humane Dannelse ikke
forvexles med den christelige.

Den humane Dannelse, der fremvoxer af sit eget Princip, uden Indblanding
af det Christelige, trives, saa at sige, ved egne Midler, og naaer sin
ideale Høide i Videnskaben som Videnskab. Paa dette Trin finde vi den
lærd-videnskabelige Bibelfortolker.

Den christelige Dannelse vil altid forudsætte den humane til en vis Grad;
men med denne Forudsætning udvikler den sig for Resten af sit eget
Princip, og det saaledes, at det samme Menneske kan staae paa et
forholdsviis lavt Trin af human Dannelse og dog paa et forholdsviis meget
høit Trin af christelig Dannelse. Paa dette Standpunkt hensætte vi den
eenfoldige, men ulærde Bibelfortolker.

Endelig kan den humane Dannelse have naaet et vist videnskabeligt Trin i
et Menneske, der saaledes føler sig grebet af Christendommen, at han nu
nedsætter denne Dannelse til Middel for en principmæssig og altsaa
„uvidenskabelig“ Udvikling af det christelige Troesliv. Paa dette
Standpunkt hensætte vi den eenfoldige, men dog lærde Bibelfortolker.

Forsaavidt den lærde Fortolker forklarer den hellige Skrift religiøst,
maa han væsentlig forklare den paa selvsamme Maade, som den troende
Lægmand. Men, da den lærde Skriftfortolker dog har Lærdommen og den
videnskabelige Dannelse forud for
% --- p. 65 ---
Lægmanden, vil denne Dannelse og Lærdom sætte ham istand til at fortolke
Skriften saaledes, at han i menneskelig Forstand bliver en sand
Skriftklog, en Lærer for de Eenfoldige, der have Troen, men mangle
Lærdommen.

Men i guddommelig Forstand er Skriftfortolkningen en overmenneskelig
Gave, og just derved ophøiet over enhver menneskelig Kunst. I Kraft af
denne sin dobbelte Natur tilstræber Fortolkningen en dobbelt
Fuldkommenhed, en formel og en reel.

Hvad Fortolkningens formelle Fuldkommenhed angaaer, vil Den være den
bedste Fortolker, der har den meste Dannelse; men hvad Troeserkjendelsens
Gehalt angaaer, vil Den være den bedste Fortolker, hvis Existens har den
største Inderlighed i Troen. I første Henseende kan den Skriftkloge
belære den Læge; i sidste Henseende kan den eenfoldige Lægmand ligesaavel
belære den Skriftkloge, som omvendt belæres af ham. Det Humane i
Christendommen er, at den Ene saaledes kan lære den Anden; men det
Christelige, det absolut Guddommelige, er, at den Ene i Grunden slet ikke
kan lære den Anden.

Hvor vil Theologen nu hen med sin religiøs-videnskabelige Fortolkning?
Den religiøs-ethiske Skriftfortolkning staaer deels udenfor Videnskaben,
deels over Videnskaben, men aldrig i Videnskaben!

Er det nu virkelig sandt, at der ikke findes og ikke kan findes nogen
Mellemsphære, hvori den religiøs-ethiske Fortolkning skulde lade sig
aflede, bestemme eller beherske af den lærd-videnskabelige; saa er dette
jo et nyt Beviis for, at \emph{Christendommen er høiere end Videnskaben,
og Evangelietroen ueensartet med Theologien.}

%% ============================================================
%%  SJETTE FORELÆSNING (pp. 66--80 = PDF 79--93)
%% ============================================================
\chapter*{Sjette Forelæsning}
\addcontentsline{toc}{chapter}{Sjette Forelæsning}
\label{ch:forel6}
\markboth{Sjette Forelæsning}{}

\begin{center}\itshape
Den theologiske Skriftfortolker beraaber sig paa Bibelens Inspiration i
det Totale, men tør paa ingen Maade vedkjende sig dens Inspiration i det
Enkelte. Da der nu saaledes er Tvivl om alt det Enkelte, maa der
naturligviis ogsaa være Tvivl om det Totale: for Theologen er Bibelens
Inspiration en tom Talemaade. Den evangeliske Skriftfortolker, der ifølge
Existensens Lov overalt seer det Totale concentreret i det Enkelte,
opfatter hvert enkelt Skriftsted som inspireret, under den kirkelige
Forudsætning, at den hele Bibel er inspireret: Evangelietroen er
ueensartet med Theologien.
\end{center}

% --- p. 66 ---
At fortolke Bibelen, Kirkens hellige Skrift, er med andre Ord at udfinde,
fastsætte og fremstille det væsentlige Forhold mellem Mening og Udtryk, Aand
og Bogstav (πνεῦμα og γράμμα). Inspirationen viser os her sin dobbeltsidige
Betydning. I den inspirerede Skrift ere Indhold og Form uadskillelig
forbundne; Alt maa forstaaes efter Bogstaven: Den, der ringeagter Bogstaven,
krænker Aanden. Men Bogstavens Tjener, der kun seer Bogstaven, krænker jo
ogsaa Aanden; thi Bogstaven ihjelslaaer, men Aanden levendegjør.

Den hermeneutiske Opgave, at bestemme Forholdet mellem Aanden og Bogstaven,
har fra Theologiens første Begyndelse beskjæftiget Theologerne gjennem
Kirkens mange Aarhundreder, uden at man dog endnu har kunnet formidle den
Modsætning,
% --- p. 67 ---
der allerede tidlig gav sig et videnskabeligt Udtryk i Striden mellem den
alexandrinske og den antiochenske Skole.

Vistnok er den allegoriserende Retning med alle derunder hørende
Phantasterier bleven meer og meer tilbagetrængt af den stedse fremskridende
Videnskabelighed; men den grammatisk-historiske Fortolkning har dog desuagtet
ligesaa lidt været istand til at udgrunde Aanden, som at hævde Bogstaven.

I en vis Forstand kan den grammatiske Fortolker jo vistnok hævde Bogstaven,
forsaavidt han nemlig, ifølge den ovenfor anerkjendte Methode, kan udfinde
Textens bogstavelige Mening; men denne bogstavelige Mening bliver saa med det
Samme forvandlet til en reent menneskelig Mening. „Gud selv har jo dog
alligevel ikke tænkt denne endelige Tanke, ikke talt disse endelige Ord, ikke
construeret disse grammatikalske Sætninger!“ Men hvor bliver saa Aanden af
med samt den dybere guddommelige Mening?

Dette religiøse Spørgsmaal kan ikke besvares videnskabeligt. For den
videnskabelige Betragtning bliver Aanden i de enkelte Ord og Sætninger en
reent menneskelig Aand, Meningen en reent menneskelig Mening. Det Dybere, det
Overmenneskelige, den guddommelige Mening (forsaavidt en saadan overhovedet
antages) er ikke at finde paa noget enkelt Sted. Den dybere Mening kommer
først tilsyne, naar man sammenstiller en Mangfoldighed af Skriftsteder og
indordner dem under Ideens Almindelighed, det er med andre Ord: naar man seer
paa det Totale. Det Enkelte er menneskeligt, det Totale guddommeligt! saa er
det da ret charakteristisk, at Theologien i den nyere Tid har et saa
overvættes Hang til at forgude det Totale.

Forelægger man en theologisk Bibelfortolker et bestemt Sted af den hellige
Skrift, og spørger ham, om han nu ogsaa tør vedkjende sig dette Sted som et i
Virkeligheden inspireret Sted, saa at disse Ord, disse Udtryk, disse
Sætninger, alle ere gjennemtrængte af Guds Aand: da kommer han vist i stor
For-
% --- p. 68 ---
legenhed. Vil man derimod disputere med ham om Stedets Forhold til det hele
Skrift og det hele Skrifts Forhold til Skriftens Totalitet; skal han strax
være villig til at vedkjende sig Inspirationen. Vælg Evangeliet om Christi
Fødsel, dette enkelte bestemte Evangelium, som det findes hos Lukas; spørg
saa Theologen, om han i Videnskabens Navn tør vedkjende sig denne evangeliske
Fortælling som en af Gud inspireret, som en i alle Maader ufeilbar og
troværdig Fortælling: han kommer i stor Forlegenhed, og forsøger at undvige
Spørgsmaalet; men hvis Nogen derimod vil raisonnere med ham om Lukas i det
Totale, om de fire Evangelier i det Totale, om det N.\ Tst., ja om hele
Bibelen i det Totale: da skal Theologen være villig til at forsvare baade
Inspirationen og den dybere Mening og det guddommelige Vidnesbyrd, og hans
Forsvar skal blive desto kraftigere, jo mere det faaer Lov at brede sig i det
Totale.

Kun det Enkelte saarer; det Totale kommer Ingen for nær. Det Totale er i
Fjernhed, kun det Enkelte er det Nærværende; det Totale er i svævende
Ubestemthed, kun det Enkelte er bestemt og afgjørende; det Totale er for den
videnskabelige Betragters objective Phantasie, det Enkelte er for den Troende
i Tilegnelsens Inderlighed.

Modsætningen mellem den theologiske og den evangeliske Fortolkning kan altsaa
med Hensyn til Forholdet imellem Aand og Bogstav udtrykkes saaledes:

\emph{Ifølge den theologiske Fortolkning bliver det Enkelte sat som Moment i
det Totale, ifølge den evangeliske Fortolkning bliver det Totale sat som
concentreret i det Enkelte.}

For at belyse denne Modsætning nærmere vælger jeg et bestemt Exempel:

\emph{Naar Nogen slaaer dig paa din høire Kind, vend ham ogsaa den anden til}
(Matth.\ 5, 39).
% --- p. 69 ---
I disse Christi Ord er her ikke blot Eenhed af Aand og Bogstav, men her er
tillige Forskjel. Man kan gjerne indrømme den videnskabelige Fortolker, at
denne Befaling ikke kan tages ligefrem efter Bogstaven, uden at Bogstaven
ihjelslaaer Aanden; men det er da ogsaa vist, at kun den religiøs-ethiske
Fortolkning formaaer at hæve sig til det Punkt, hvor Aanden igjen forener sig
med Bogstaven, saa at Idealet selv sees i det Bogstavelige.

Hvis Nogen tager Ordene ligefrem efter Bogstaven, og altsaa forstaaer dem i
Aandløshed, vil Misforstaaelsen let vise sig; man behøver vistnok ingen
videnskabelig Veiledning for at indsee dette. Et Menneske kan, naar han faaer
et Slag paa den høire Kind, gjerne vende den venstre til, uden derfor at være
en Christen; ligesom omvendt: et Menneske kan være en sand Christen, uden at
han derfor vender den venstre Kind til, hver Gang han faaer et Slag paa den
høire. At den Troende nu skulde gjøre sig skyldig i en saa aandløs,
vilkaarlig Misforstaaelse af det bogstavelige Udtryk, er ingenlunde at
befrygte; derimod er det i høi Grad at befrygte, at den videnskabelige
Fortolker, der her kommer saa let til en rationel Adskillelse af Aand og
Bogstav, skal omforklare os Budet saa aandeligt, at han aldeles udsletter
Christendommens Mærke.

Med en saadan Omforklaring gaaer det som af sig selv. Har man først udfundet,
at Christi Bud ikke kan forstaaes bogstaveligt; hvad ligger da nærmere end at
uddrage en almindelig Moral af det billedlige Udtryk, og oplyse denne
almindelige Moral ved en Sammenstilling af flere herhen hørende Skriftsteder.
Lad nu den almindelige Moral være den, at en Christi Discipel skal vide at
beherske sig selv, vise sig forsonende og eftergivende, være villig til at
forlige sig med sin Modstander, overvinde det Onde med det Gode o.\ s.\ v.
Idet nu den videnskabelige Fortolker udlægger denne almindelige Moral i al
Almindelighed efter „Skriftens Analogie“, kommer han i Tanker om, at andre
Old-
% --- p. 70 ---
tidens Moralister og vise Mænd have fremsat lignende moralske Forskrifter i
lignende billedlige Udtryk; slige classiske Steder, være sig af Confucius
eller Pythagoras eller Xenophon eller Plato eller Cicero eller Horats eller
Seneca, eller hvo det nu er, bør da ogsaa citeres. Ved at citere slige Steder
viser Fortolkeren ikke blot sin store Belæsthed og sit lærde Overblik, men
Fortolkningen vinder tillige i Videnskabelighed: Christus træder i Række med
de ædleste Aander, der have levet paa Jorden, og de ædleste Aander træde
igjen i Række med Christus; saa udvides Troens Analogie, og Skriftens
Guddommelighed erkjendes i det Totale.

Anderledes, ja ganske anderledes forholder det sig med den religiøs-ethiske
Fortolkning. Den Religiøse, der i Troens Inderlighed stræber at forstaae
Christi Bud, begynder ikke med at sprede sin Opmærksomhed paa en
Mangfoldighed af Skriftsteder, giver sig ikke hen til at blade i Bibelen, for
at underholde sig selv og Andre med en interessant Sammenligning af
„ubilledlige“ og „billedlige“, „prosaiske“ og „poetiske“ Udsagn; endnu mindre
fortaber han sig i lærde Undersøgelser, for at opdage, hvad andre Viismænd
muligviis have sagt om det Samme; nei, den Religiøse samler sit Sind og
fæster sit Blik paa det ene Bud, der her just sigter paa ham som et Christi
Bud; den Religiøse taber sig i Tilegnelsen af Budets Enkelthed. Naar han da
læser Budet og forstaaer, at dette Bud angaaer ham selv i Særdeleshed (som
det jo angaaer hver Christi Discipel i Særdeleshed); naar han fatter og
forstaaer, at det som Christi Bud maa være ham et ubetinget Bud, et Bud,
hvorefter han skal leve: da vil det billedlige Udtryk strax lære ham, at
Budet har alle christelige Mærker.

\emph{„Naar Nogen slaaer dig paa din høire Kind:“} --- ja, hvad er da
naturligere, end at Du, den Mishandlede, søger at hævne det paa din
Modstander, hvad er da naturligere end at Du, den Forurettede, søger at
skaffe dig Ret, ja uvilkaarlig fristes
% --- p. 71 ---
til at tage dig selv til Rette? --- Aldrig er Selvkjærligheden farligere, end
i det Øieblik den gjør sig til Eet med Retsfølelsen og hidser til
Gjengjældelse. Som nu Selvkjærligheden blusser op, kommer Christi Bud og vil
slaae den ned. Det første naturlige Indtryk, det naturlige Menneske faaer,
naar Budet saa øieblikkelig og ubetinget vil tilintetgjøre den forædlede, i
Vredens Øieblik retfærdiggjorte Selvkjærlighed er --- Forargelsens Indtryk.

At man overhovedet, i Almindelighed og til en vis Grad skal fornegte sig
selv, at man overhovedet, i Almindelighed og til en vis Grad skal være
forsonlig, eftergivende og villig til at forlige sig med sin Modstander, det
kan Enhver forstaae, der har nogen sædelig Sands: det kan Israeliten
forstaae, det kan den ædle Hedning forstaae! men at det virkeligt skulde være
bogstaveligt Alvor med en saadan Selvfornegtelse, som Christi Bud her
paabyder; nei, det kan kun en troende Christen forstaae, naar han i ubetinget
Lydighed forsøger at gjøre derefter.

Hedningen forædler Selvkjærligheden, men fornegter den ikke (selv hos Moses
hedder det: Øie for Øie, Tand for Tand); kun i Christendommen bliver
Selvkjærligheden dømt, og Selvfornegtelsen sat som ufravigelig Betingelse for
at indgaae i Himmeriges Rige. Denne ubetingede Selvfornegtelse kan ikke
erhverves ved en blot menneskelig Resignation. En hedensk Viis, en Stoiker
vil ved fortsat Øvelse kunne bringe det til en forbausende Grad af
Selvbeherskelse; vil maaskee kunne bringe det dertil, at han aldrig overiler
sig, aldrig forraskes af Lidenskabens Udbrud; vil maaskee vinde en saadan
% printer's defect: a raised speck (broken comma?) stands after „Magt“; no
% punctuation is called for and none is transcribed.
Magt over sig selv, at han rolig modtager Slaget paa den høire Kind uden at
forandre en Mine, uden engang at finde det Umagen værd, at slaae igjen; men
her er ogsaa Grændsen for Resignationen. Gaaer man ud over denne Grændse, saa
viser sig det trodsige Overmod, den iiskolde Foragt, der endogsaa nedlader
sig til at lege med Uretfærdigheden; som, hvis Stoikeren ikke blot vilde
taale Slaget paa den høire
% --- p. 72 ---
Kind, men, til Beviis paa at han ikke agtede det, endog ligesom haane
Forhaanelsen, og vende den venstre Kind til.

En saadan Trods og Foragt, hvori Selvkjærligheden hemmelig triumpherer, en
saadan Trods og Foragt er dog saa himmelvidt forskjellig fra den christelige
Selvfornegtelse, at Mennesket i al Skrøbelighed hellere maatte gjengjælde det
uretfærdige Slag og derefter angre sin Overilelse, end saaledes forgude sin
egen Stolthed og forhærde sig i et umenneskeligt Hovmod.

Naar altsaa Christi Bud forudsætter, at Nogen kunde tilføie Dig en krænkende
Uret, og slaae Dig paa din høire Kind; da forudsætter Budet ogsaa, at et
saadant Slag smerter Dig, at Du som Menneske i Kjød og Blod bitterlig føler
Krænkelsen og Forhaanelsen, ja det forudsætter tillige, at det endog i samme
Øieblik stod i din Magt, at gjengjælde Uretfærdigheden ved at gjengjælde
Slaget. Naar da nu Budet desuagtet ikke alene forbyder Dig at slaae igien,
naar Nogen slaaer Dig paa din høire Kind, men endog udtrykkelig befaler Dig
at vende Modstanderen den venstre Kind til: saa er dette Bud den naturlige
Fornuft et forargeligt Bud, et Bud af unaturlig og umenneskelig Strenghed.

For den Troende derimod lyder dette Bud som et guddommeligt Kjærlighedsbud,
et Bud, der vel ikke kan opfyldes af det naturlige Menneske, men dog skal
opfyldes, fordi dets Opfyldelse ikke blot fordrer af Mennesket en ubetinget
Lydighed, men tillige viser hen til en overmenneskelig Bistand. Kun Christus,
der forjætter en saadan Bistand, kan i Alvor fremsætte et saadant Bud; men,
naar Budet i Aand og Sandhed opfattes som Christi Bud, vil ogsaa Idealet
erkjendes i det bogstavelige Udtryk. Vilde Christus indskjærpe
Selvfornegtelsen i et bestemt Exempel, hvor kunde der da findes et Exempel
mere talende end dette? Hvor afgjørende, hvor ubetinget maa Selvfornegtelsen
ikke være, naar Den, der faaer et Slag paa den høire Kind, strax og som uden
Betænkning skal vende den venstre Kind til! Og hvilken
% --- p. 73 ---
Inderlighedens Kraft maa det Menneske ikke besidde, der skal opfylde dette
Bud! thi Enhver, der øver Selvfornegtelse saaledes, at Verden forstaaer, at
det er Selvfornegtelse, og at Selvfornegtelsen er fornuftig, faaer dog den
Løn, at Verden erkjender ham for en dydig Mand (thi saa billig er dog den
selvkjærlige Verden, at den gjerne vil prise Selvfornegtelsen, naar blot en
Anden vil øve den); men Den, der uden Modstand og Klage vender den venstre
Kind til, naar han har modtaget Slaget paa den høire, hvad siger Verden om
ham? om ham siger Verden ikke, at han er en fornuftig og dydig Mand; men den
siger om ham, at han enten maa være et utaaleligt, trodsigt
% printer's defect: the 1850 print sets „Mennneske“ with three n's; transcribed
% as printed.
Mennneske, der foragter Alt, eller ogsaa en enfoldig Daare.

Saaledes har dette Bud da alle christelige Mærker, hvorpaa det bogstavelig
kan kjendes, at det i Aand og Sandhed maa være Christi Bud. Det Totale er her
concentreret i det Enkelte; den hele christelige Ethik med alle Kjærlighedens
Gjerninger er sammentrængt i det ene Bud: \emph{naar Nogen slaaer Dig paa din
høire Kind, vend ham ogsaa den anden til.} Vil Nogen forsøge at indrette sit
Liv efter dette Bud, da lover Evangeliet, at Guds Aand vil staae ham bi, og
Tilværelsen forhjælpe ham til principielt at opdage hele Christendommen; men
Den, der ikke vil og ikke kan leve derefter, kan heller ikke fatte, hvad der
ligger i Budet; han forstaaer det kun i Phantasie, i Mulighed (som vi i dette
Øieblik jo vel forstaae det i Phantasie og Mulighed); men i Virkelighed
forstaaer han det ikke.

At den religiøs-ethiske Fortolkning af det her omhandlede Skriftsted er
aldeles ueensartet med den videnskabelige, kan nu let bevises.

Forsaavidt Forfatteren forholder sig videnskabeligt til sin Text, og henfører
Betydningen af hvert enkelt Skriftsted til Totalitetens almindelige
Kategorier, kan han jo ikke engang i Forestillingen naae saavidt, at han
berører Existensen; thi den
% --- p. 74 ---
existentielle Tænkning adskiller sig netop derved fra den videnskabelige, at
den paa alle Punkter berører Enkeltheden, det enkelte Menneske, det enkelte
Bud, de enkelte Forhold og Betingelser, hvorunder Budet finder sin Anvendelse
o.\ s.\ v. Men selv om man vilde henregne den Fortolkning, der omsætter det
væsentlige Indhold i lutter existentielle Bestemmelser, til den
videnskabelige, og paastaae, at den Prøve, vi her have givet paa en
religiøs-ethisk Exegese, i Grunden ligesaa godt kunde kaldes videnskabelig,
som den philologisk-historiske: da er dette en reen Misforstaaelse.

Vel er det sandt, at den religiøs-ethiske Fortolkning, forsaavidt den
meddeles i et docerende Foredrag, har en vis Lighed med den videnskabelige,
den nemlig, at den doceres. Men forsaavidt den religiøs-ethiske Fortolkning
doceres, viser den paa alle Punkter ud over sig selv, erklærer sig
utilstrækkelig og ophæver sig selv. Den Existens, hvorom det docerende
Foredrag handler, er nemlig ingen virkelig, men kun en forestillet Existens.
Derfor kan man heller ikke slutte, at En er virkelig religiøs, fordi han har
en vis Færdighed i at fremsætte religiøs-ethiske Forklaringer. I Foredragets
Form er den religiøs-ethiske Forklaring altsaa dog endnu kun en forestillet
Forklaring, en Forklaring der blot adskiller sig fra den videnskabelige deri,
at den vender sig mod Existensen, fordrer Existensen, og gjør de existentielle
Opgaver kjendelige.

Den religiøs-ethiske Fortolkning har alene sin Sandhed i den virkelige
Existens; kun i den virkelige Existens er den for Alvor ueensartet med den
videnskabelige.

I den virkelige Existens er den Troende, der fortolker sin Bibel, ingen
forestillet Figur, men et virkeligt Menneske, og Forstaaelsen af Skriftens
Ord ingen forestillet, men en virkelig Forstaaelse. Den virkelige Forstaaelse
er kun i den virkelige Anvendelse; de Livsforhold og Tilfælde, i hvilke
Christi Bud finde deres Anven-
% --- p. 75 ---
delse, ere da heller ikke forestillede, men virkelige; men det virkelige
Tilfælde er hver Gang et enkelt bestemt Tilfælde.

Fastholde vi nu dette Synspunkt, saa følger heraf, at den religiøs-ethiske
Fortolkning ikke blot er forskjellig fra den videnskabelige, men tillige saa
ueensartet med den, at det endog i strengeste Forstand er det ene Menneske
umuligt, at fortolke Bibelen for det andet, efterdi det er det ene Menneske
umuligt, at sætte sig aldeles ind i det andet Menneskes Existens. Hvorvidt
den Troende i Virkeligheden forstaaer og tilegner sig Christi Bud: „naar
Nogen slaaer Dig paa din høire Kind, vend ham den \emph{anden til,}“ beroer
jo ene og alene paa, hvorledes han lever, men dette maa Enhver vide med sig
selv: Gudsforholdet er et Samvittighedsforhold.

Den Selvfornegtelse, Budet indskjærper, er ubetinget; men hvert Menneskes Liv
er Endelighedens Betingelser underlagt. Det bliver ikke Mennesket givet, i
hvert Øieblik at sætte det Ubetingede igjennem, og dog kræver det sig af ham
i hvert Øieblik. Det Ubetingede kan desuden paa mangfoldige Maader virke i de
forskjellige Betingelser. Et Menneske kan paa utallige Maader øve sig i
religiøs Selvfornegtelse, kan paa utallige Maader lægge denne Selvfornegtelse
for Dagen, stundom endog saaledes, at den inderligste Selvfornegtelse viser
sig i det Udvortes som anmassende Selvraadighed. Under slige Vilkaar er det
da aldeles umuligt for det ene Menneske at raade det andet: Enhver maa i
Troen raade sig selv. Naar dette Vendepunkt indtræder (og det indtræder, hver
Gang Existensen ret alvorlig trænger paa), forstummer al gjensidig
Underviisning; Enhver griber i sin egen Barm --- efter Forstaaelsen, Enhver
slutter Budet ind i sit eget Indre, og slutter sig, saa at sige, selv inde
med Budet, for paa eget An- og Tilsvar at forstaae, hvorledes dette Bud i
dette givne Tilfælde kan og bør anvendes.

Idet nu Existensens Alvor saaledes standser den gjensidige
% --- p. 76 ---
Underviisning, Formaning og Veiledning, og splitter Menneskene ad, at Enhver
i Afgjørelsens Stund maa raade sig selv; kan Meningen dog umulig være den, at
den egensindige, vilkaarlige Selvraadighed heraf skulde drage ny Fordeel.
Tvertimod, det er saa langt fra, at den Enkelte her henvises til vilkaarlig
Selvraadighed, at han nu først tilfulde forstaaer, hvor høilig han trænger
til Raad, hvor umuligt det er ham, at raade sig selv. Til den naturlige
Selvfølelse kan Mennesket ikke betroe sig, thi Selvfølelsen er i Forstaaelse
med Selvkjærligheden; men her gjælder det jo om Selvfornegtelse. Den
reflecterende Forstand kan heller ikke løse ham Opgaven; thi deels bliver
Forstanden bestukken af Selvkjærligheden, deels bliver den hildet og fangen i
sin egen Reflexion. Istedetfor at løse Opgaven vil Reflexionen kun gjøre den
endnu langt vanskeligere, idet den nemlig vil gjøre den dialektisk. Paa den
ene Side kræves en absolut Selvfornegtelse, en ubetinget Selvopoffrelse; paa
den anden Side kræves der, at Mennesket skal erkjende sin Grændse. Et
Menneske maa jo dog altid give Noget efter for de udvortes, forhaanden
værende Omstændigheder; man kan jo dog ikke ligesom med et Ryk sprænge alle
Existensens Betingelser, for at gribe det Ubetingede. For Reflexionen skiller
Budet sig ad, og deler sig i tvende indbyrdes modsigende Bud: grib det
Ubetingede! og: erkjend din Grændse! Denne Modsigelse kan ikke formidles ved
fortsat Reflexion: jo længere der reflecteres, desto større bliver
Vanskeligheden; desuden er her jo i Afgjørelsens Øieblik ikke Tid til at
reflectere, thi Øieblikket er kostbart: det gjælder om at beslutte og handle.

Bliver nu Mennesket saaledes henviist til sig selv uden dog at kunne raade
sig selv, saa vilde der jo slet ingen Raad være for ham, hvis ikke Guds Aand
vidnede med hans Aand i Samvittigheden, og indgav ham, hvad han i det
afgjørende Øieblik skulde gjøre. Saa siger Evangeliet: Den Aand, der lærer
Men-
% --- p. 77 ---
nesket det rette Valg i Afgjørelsens Stund; den Aand, der bestemmer
Selvfornegtelsen i Forhold saavel til de ydre Vilkaar som til det
paagjældende Menneskes Evner og Dannelse, den raadende og styrende Aand, der
opløfter Mennesket over sig selv i Selvfornegtelsens Øieblik, den Kraftens
Aand, der gjør, at den Forurettede bliver stærkere end sin Modstander, idet
han --- vidunderligt nok! --- bliver stærkere end sig selv; den Aand, der
gjør, at han, der faaer Slaget paa den høire Kind, kan i Sagtmodighedens og
Ydmyghedens Overlegenhed vende den venstre Kind til: kunde denne Aand vel
være nogen anden end Guds egen hellige Aand?

Hvad der i Videnskaben er betydningsløst og i Theologien en tom-høitidelig
Talemaade, \emph{at nemlig den Hellig Aand er den eneste rette Fortolker af
den hellige Skrift}, det bliver en virkelig Sandhed i den Troendes Existens;
thi den Hellig Aand er en Existensens Aand.

I den virkelige Existens vil Aanden altid svare til Troen, som Troen til
Aanden. Vov at troe, og Du skal faae Aandens Kraft til at gribe det
Ubetingede; vov at fornegte Dig selv, og Du skal ikke blot faae Kraft til at
bryde med Verdensklogskabens Hensyn, men ogsaa Taalmod og Viisdom til at
forstaae Endelighedens og Skrøbelighedens Vilkaar, thi Kraftens Aand er
tillige Raadets og Viisdommens Aand. Dette er Loven, den aabenbarede Guds
Lov, hvorefter hver menneskelig Existens sig haver at rette, at Den, der har
en stærk Tro, faaer en stærk Aand, og forstaaer Skriftens Ord i Aand og
Sandhed; men den, der kun har en svag Tro, faaer en svag Aand, og forstaaer
Ordet skyggeviis. Dette er Loven, Skriftfortolkningens ufravigelige Lov: Hver
Den, der hører Ordet og gjør derefter, erfarer Ordets Kraft og maa lignes med
en forstandig Mand, der bygger sit Huus paa Klippen; men hvis Nogen hører
Ordet og ikke gjør derefter, da kan han heller aldrig erkjende Or-
% --- p. 78 ---
dets Kraft, men skuffer sig selv ved en indbildt Forstaaelse og bliver at
ligne med den uforstandige Mand, der bygger sit Huus paa Sand.

Theologernes Mening, at den Hellig Aand nok skal indlade sig paa deres lærde
Forskninger og omsider komme dem til Hjælp med en sand videnskabelig
Fortolkning af Skriften, er hverken begrundet i Videnskabens eller i
Aabenbaringens Væsen, men er, reentud sagt, hverken meer eller mindre end en
theoretisk Confusion af ueensartede Opgaver. Denne Confusion yttrer sig, hver
Gang man seer paa det Enkelte; thi Theologen, som Theolog, har aldrig den
Hellig Aand i det Enkelte, men altid i det Totale, aldrig i Virkelighed, men
altid i Mulighed, aldrig i Existens, men altid i Phantasie.

Men den Aand, der saaledes er i det Totale, uden at komme tilsyne i det
Enkelte, er en nebuløs og svævende Aand, en Aand, der vist ikke forarger
noget Menneske, men visselig heller ikke frelser noget Menneske, en Aand, der
stifter Talemaade, som en Digter stifter Phantasie.

Den theologiske Videnskabs videnskabelige Forsvarere ville uden Tvivl finde
denne Paastand noget nærgaaende; hertil svarer jeg, at Paastanden kun er
nærgaaende, fordi den er sand, og skal til Overflødighed bevise, hvad jeg
siger.

Det er jo vitterligt, at ikke blot den orthodoxe Theolog, men ogsaa
Rationalisten, ikke blot den speculative, men ogsaa den uspeculative
Rationalist beraaber sig under Fortolkningen af den hellige Skrift paa
% the „e“ of „hellige“ is blotted with ink in this copy; the reading is
% nevertheless certain from the sense and the surrounding letters.
Aandens Vidnesbyrd i det Totale (\textit{testimonium spiritus sancti
theoreticum}). Heraf kan man nu uddrage tvende Slutninger: enten er den
Hellig Aand, der saaledes giver sig Vidnesbyrd i alle disse Fortolkninger,
selv baade orthodox, heterodox, speculativ, uspeculativ o.\ s.\ v.; eller
ogsaa maa den Hellig Aands Vidnesbyrd være ophøiet over disse endelige
Skolers endelige Modsætninger. Men hvilken af disse
% --- p. 79 ---
Antagelser man end foretrækker, maa den herskende Theologie gaae tilgrunde. I
første Tilfælde er Theologernes Hellig Aand jo en sig selv modsigende og
modstridende Aand; men en Videnskab, hvis Aand er i væsentlig Strid med sig
selv, kan umulig bestaae. I sidste Tilfælde er den Aand, paa hvis Vidnesbyrd
Theologerne beraabe sig, en overvidenskabelig Aand, en Aand, i hvis Væsen
Theologien altsaa selv ingen Andeel har, eftersom Theologien jo er og vil
være Videnskab.

Hvad den religiøs-ethiske Bibelfortolkning angaaer, er det ingenlunde
„Theologerne“, vi trænge til; men, forsaavidt en menneskelig Veiledning i
denne Retning overhovedet er mulig, er en existentiel Forfatter vistnok den
Veileder, vi høilig kunde trænge til, hvis vore theoretiske Indbildninger
blot ikke forhindrede os i at forstaae Veiledningen.

Ved en existentiel Forfatter tænker jeg mig da en Saadan, der ikke blot har
en ganske særegen Gave til at udvise Existensens Kategorier og de tilsvarende
Opgaver i alle mulige Forhold, men en Forfatter, der tillige har den
Selvfornegtelse at ville offre et Livs Kræfter paa at fremarbeide disse
Opgaver; en Forfatter, der i Spøgens Alvor end ikke er bange for at
modarbeide sig selv, for desto redeligere at tjene sin Idee; en Forfatter,
der, naar det saa skal være, hellere bærer paa Døgnets Misforstaaelser, end
opgiver Forstaaelsen med en evig Sandhed; en Digter i den umiddelbare
Forelskelse, hvormed han seer paa det Ideale; en Ethiker i den strenge
Forstaaelse, at det just er hans egen Existens, det Ideale vil fordre; en
religiøs Individualitet i den ydmygende Erkjendelse, at han saa alligevel
ikke kan --- hvad dog hvert Menneske burde kunne!

Saaledes, omtrent saaledes forestiller jeg mig den Forfatterexistens, der vel
allernærmest maatte være rustet til at føre Krig mod Nutidens
„Sandsebedrag,“ og gjøre Christendommen kjendelig for den moderne Reflexion.
% --- p. 80 ---
At Theologerne, jeg mener de ægte-videnskabelige, især de
ægte-speculativ-videnskabelige, naturligviis helst vilde oversee en saadan
Forfatterexistens, eller, hvis Phænomenet i Litteraturen nu var saa
paafaldende, at det, idetmindste for Conversationens Skyld, ikke længere lod
sig oversee, da fortolke det som --- en „uvidenskabelig“ Eensidighed! --- en
„mærkværdig“ Eensidighed! --- en „genial“ Eensidighed! --- en „ekstatisk“
Eensidighed! er endda saa temmelig let at forstaae --- i Conversationen.
Denne fleersidige Fortolkning beviser jo netop, at Theologerne maae have
deres egen fleersidige Hermeneutik.

Enhver fortolke det da efter sin Hermeneutik, som han vil: jeg fortolker det
saaledes, \emph{at Christendommen er høiere end Videnskaben, og
Evangelietroen ueensartet med Theologien.}

%% ============================================================
%%  SYVENDE FORELÆSNING (pp. 81--96 = PDF 94--96, then 99--111)
%% ============================================================
\chapter*{Syvende Forelæsning}
\addcontentsline{toc}{chapter}{Syvende Forelæsning}
\label{ch:forel7}
\markboth{Syvende Forelæsning}{}

\begin{center}\itshape
I Behandlingen af den hellige Historie --- „Jesu Liv“ --- vil Theologen
modarbeide Fritænkeren; men den videnskabelige Kritik, hvorpaa han gaaer
ud, viser tydeligt nok, at han i Principet er enig med Fritænkeren. Den
Troende, der i Opfattelsen af Evangeliernes Christus udelukkende holder
sig til den religiøs-ethiske Kritik, er i Principet enig med
Evangelisterne: Evangelietroen er ueensartet med Theologien.
\end{center}

% --- p. 81 ---
Som bekjendt, ere Tydsklands Theologer for nærværende Tid ivrigt
beskjæftigede med deres kritisk-videnskabelige Bearbeidelser af Jesu Liv
(Das Leben Jesu). „Den kritiske Indledning“ har til den Ende opbudt
Theologiens forskjellige Potenser, baade Hermeneutiken og Dogmatiken og
Ethiken og Apologetiken samt Polemiken og Ireniken til at arbeide med
forenede Kræfter, at man saaledes med forenede Kræfter kunde virke til
det fælles Maal, den store Opgaves Løsning. Planen er da saaledes sat i
Værk; men, det kan ikke negtes, Foretagendet har hidtil havt en ganske
besynderlig Skjæbne.

Den moderne Idee til et „Jesu Liv,“ der blev undfangen og født af en
berømt Theolog, havde nemlig den mærkelige Skjæbne, at den strax efter
Fødselen blev knæsat af en skarpsindig Philosoph, for desto bedre at
optugtes og dannes i den moderne Videnskabs egen Aand. Det var
Kritikeren David Strauß, der hentede Ideen fra Schleiermachers Katheder,
satte den i
% --- p. 82 ---
Skole hos Hegel, og derpaa siden efter lod den optræde i sit navnkundige
Værk: „Das Leben Jesu“ --- til Skræk og Forargelse for Theologer af alle
Farver. Aldrig har Theologien lidt et fuldstændigere Nederlag, og aldrig
har den negative Kritik vundet en fuldstændigere Seier, end netop i dette
videnskabelige Sammenstød. Mange theologiske Skrifter ere udgangne,
mange theologiske Angrebsplaner forsøgte, meget theologisk Blæk udgydt
til videnskabelig Gjendrivelse af den store Fritænker: --- Alt forgjæves!

% Fraktur I/J are near-identical; read as J. H. A. Ebrard (Johannes
% Heinrich August Ebrard), the glyph matching the J of „Jesu“.
Vistnok er Strauß i flere Henseender bleven corrigeret af sine mange
Modstandere; blandt Andre har Theologen J.\ H.\ A.\ Ebrard i sin
„Wissenschaftliche Kritik der evangelischen Geschichte“ (2den Deel) med
Held paaviist en Deel Vilkaarligheder i „den destructive Methode;“ men
det hjælper Altsammen ikke, saalænge det Grundlag, hvorpaa den
straußiske Kritik hviler, endnu ikke selv er destrueret. Sagen er den,
at Strauß dog har et Princip; medens Theologien derimod er og bliver
principløs.

Det straußiske Princip er nemlig intet andet end den moderne Bevidstheds
eget almindelig anerkjendte Grundprincip, det Princip: at sætte
Fornuftideen som det Absolute, forgude Naturlovene og hæve Videnskaben
til Religion. Hvormange Smaaforandringer, den religiøse Følelse og
subjective Humanitet end for Resten foretager med dette Grundprincip, saa
sees det dog igjennem alle Modificationer og Fortolkninger at være det
selvsamme.

Den straußiske Bestræbelse at hæve Videnskaben til Religion kan umulig
være ueensartet med den theologiske Bestræbelse at hæve Religionen til
Videnskab; thi begge Bestræbelser gaae jo ud fra den Forudsætning, at
Religion og Videnskab ere eensartede. Men, naar Theologerne i en saa
væsentlig Forstand staae paa selvsamme Standpunkt, som Strauß: hvorledes
skulle Theologerne da kunne gjendrive Strauß?

Vi see det tydeligt paa Ebrards Forhold til Strauß. Strauß gaaer ud fra
den Forudsætning, at enhver Mirakeltro er Over%
% --- p. 83 ---
tro, og beviser saa deraf, at alle Mirakelhistorierne ere Myther, i det
N.\ Tst.\ ligesaavel som i det Gl.; Ebrard gaaer ud fra den
Forudsætning, at Gud kan gjøre Mirakler, og beviser saa deraf, at
Mirakelfortællingerne i de fire Evangelier ere troværdige; Strauß er
Naturalist, Ebrard er Supranaturalist; Strauß er vantro, Ebrard troende;
Strauß river ned, Ebrard bygger op: forsaavidt staaer da Ebrard jo
vistnok paa et ganske andet Standpunkt end Strauß. Men, see vi nærmere
til, vil det snart vise sig, at Modsætningen endda ikke er saa stor, som
det ved første Øiekast synes, men at den, som Mediationen siger, ret godt
lader sig ophæve i en høiere Eenhed.

Trods al den Uenighed, der for Resten hersker imellem de tvende
Modstandere, ere de dog begge enige i et meget vigtigt og væsentligt
Punkt; de ere nemlig begge to enige deri, at man under Behandlingen af
„das Leben Jesu“ naturligviis maa staae paa Videnskabens Standpunkt.
Denne Forudsætning deler Theologen med Fritænkeren, og med denne
Forudsætning knuser Fritænkeren Theologen.

En Supranaturalist paa Videnskabens Standpunkt er en Modsigelse; en
videnskabelig Kritiker paa Mirakeltroens Standpunkt er en Modsigelse;
skal Christendommen forklares i Videnskaben, da er den allerede forklaret,
det er da allerede afgjort, hvad den er: den er en Mythologie; skal den
hellige Historie kritiseres af Videnskaben, da er den allerede
kritiseret, det er allerede afgjort, hvad den er: den er en Kjæde af
Myther, en stor Mythe.

„Det Guddommelige“ --- siger Strauß --- „kan ikke være skeet saaledes;
det saaledes Skete kan ikke have været guddommeligt --- saaledes vil
Differensen udtale sig, og, naar Fortolkningen stræber at bilægge den,
vil den stræbe efter enten at fremstille det Guddommelige som ikke skeet
saaledes, altsaa negte de gamle Kilder deres historiske Gyldighed, eller
vise, at det Skete
% --- p. 84 ---
ikke saaledes er guddommeligt, altsaa bortforklare det absolute Indhold
af disse Bøger. I begge Tilfælde kan Forklaringen gaae partisk og
upartisk tilværks; partisk, naar den er blind mod Bevidstheden om
Differensen mellem den nye Dannelse og de gamle Kilder, og kun bilder sig
ind at udfinde disse sidstes oprindelige Mening; upartisk, naar den klart
erkjender og aabenbart tilstaaer, at den seer Det, hine gamle Skribenter
fortælle, anderledes, end de selv have betragtet det.“

„Den gamle, især orientalske Verden“ --- hedder det endvidere --- „var,
ifølge sin overveiende religiøse Retning og sin ringe Kundskab om
Naturlovene, om Sammenhængen i den jordiske, endelige Væren, noget saa
Løst, at den var istand til at springe fra ethvert Punkt i samme over paa
det Uendelige, og ved enhver enkelt Forandring i Naturen og
Menneskeverdenen at betragte Gud som den umiddelbare Aarsag. Fra dette
Bevidsthedens Standpunkt er ogsaa den bibelske Historie skreven . . . .
Men den nyere Tid har en Række af de møisommeligste Granskninger, der ere
fortsatte gjennem Aarhundreder, at takke for den Indsigt, at Alt i Verden
hænger sammen ved en Kjæde af Aarsager og Virkninger, der ingen
Afbrydelse taaler . . . . Denne Overbeviisning er saaledes bleven
Bevidsthed hos den nyere Verden, at den Mening eller Paastand, at en
overnaturlig Aarsag, en guddommelig Indvirkning, har et eller andet Steds
grebet umiddelbart ind, bliver i det virkelige Liv betragtet ligefrem
enten som Uvidenhed eller Bedragerie; og den er skredet frem til Extrem i
den nyere Oplysnings Anskuelse, der, lige i Modsætning af den Bibelske,
enten ganske bortfjerner den guddommelige Aarsagelighed, eller dog skyder
den saa langt tilbage, at den kun i Skabelsesakten skulde være en
umiddelbar, det vil sige, at Gud kun forsaavidt skulde virke paa Verden,
som han ved Skabelsen har givet den en bestemt Indretning. Fra dette
Standpunkt af, der i Naturen og Historien seer et fast Væv af endelige
Aar%
% --- p. 85 ---
sager og Virkninger, kunne de bibelske Fortællinger, i hvilke dette Væv
paa utallige Steder er sønderrevet ved en Indgriben af den guddommelige
Aarsagelighed, umulig vise sig som Historie.“

Man seer heraf, hvor betydningsløst det er, at ville gjendrive Strauß i
det Enkelte og forsøge sig med videnskabelige Modbeviser mod hans
nedrivende Kritik. Er det i den profane Historie allerede langt lettere
for Kritiken at rive ned end at bygge op; hvor meningsløst bliver da ikke
det videnskabelige Forsøg ved kritiske Sandsynlighedsbeviser at ville
hævde Miraklerne i den hellige Historie. Ebrards theologiske Vovestykke
har da heller ikke ret kunnet vinde Theologernes Bifald. Man læse t.\ Ex.\
Bleeks Recension over Ebrards Skrifter, og man vil faae et meget
nedslaaende Indtryk. Den ligesaa skarpsindige, som grundlærde Bleek
gjennemgaaer Ebrards Paastande, og ømmer sig hvert Skridt, som En, der
gaaer paa Naale. „Neppe troligt! --- Usandsynligt! --- Ubeviisligt! ---
Høist usandsynligt!“ ere de Udbrud, der hvert Øieblik undslipper den
ellers saa velvillige Recensent. Imidlertid er Recensenten Bleek akkurat
hildet i den selvsamme Misforstaaelse, som Forfatteren Ebrard.
Misforstaaelsen er, at man vil bedømme Miraklerne videnskabeligt, og
anlægge Sandsynlighedens Maalestok paa Det, som efter sin Natur er og maa
være det Usandsynlige.

Striden med Strauß forvandler sig øieblikkelig til en Strid mellem
Religionen og Videnskaben. Det er Kritikerens Fortjeneste at have givet
Sagen et saa stærkt Stød, at den just derved er dreven hen til dette
afgjørende Vendepunkt. Evangelietroens Strid med den straußiske Kritik
er en Competencestrid, en Souverainnitetsstrid, en Strid om den øverste
Magt, den høieste Myndighed i Aandernes Verden.

Er Videnskaben den høieste Instans, saa har Strauß i Principet Ret, saa
er det, aabenhjertig talt, og reentud sagt, ude med
% --- p. 86 ---
Christendommen, saa ligger Religionen paa sit Yderste, og forventer kun
en blid Opløsning --- i den humane Cultur.

Er Christendommen derimod høiere end Videnskaben, har den sit eget
overnaturlige Princip og bevæger sig i en Sphære, i hvis Dybder ingen
menneskelig Viden trænger ind; da vil den ogsaa endnu, som før, herske
med Souverainens Myndighed i sit eget Rige, og gjøre hvert Menneske
ansvarligt for, hvorledes han forstaaer dens Bemyndigelse og dens Forhold
til Videnskaben.

Videnskaben, som Videnskab, kan umulig anmasse sig den øverste
Myndighed. Forsaavidt Videnskaben veed, hvad den veed, veed den med det
Samme, at der er Noget, den endnu ikke veed, og Noget, den i denne
Tilværelse aldrig faaer at vide. Videnskaben veed med sig selv, at den
udøver en storartet og velgjørende Indflydelse paa Menneskeheden, idet
den fører Slægterne frem fra Oplysning til Oplysning; men den veed ogsaa,
at den ikkun har, og ikkun kan have, en meget middelbar Indflydelse paa
de virkelige Menneskers virkelige Existens. Videnskaben veed med sig
selv, at de store Begivenheder i Verdenshistorien: Folkenes
Undertrykkelse, Revolutionernes Seire, Nationernes Krige, ikke afgjøres
ved rene Fornuftgrunde. Videnskaben veed med sig selv, at der ere
Magter, dunkle, forfærdende Magter i Tilværelsen, der aldrig bøie Nakken
under dens Herredømme; Videnskaben veed, at de nye Tanker, de nye Ideer,
de nye Midler, den ved hvert Skridt tilbyder Slægten, ligesaa vel maae
tjene de Onde som de Gode, de Uretfærdige som de Retfærdige. Videnskaben
veed med sig selv, at den Sandhedserkjendelse, den bærer og opretholder i
fortsat fremskridende Udvikling, aldrig naaer og aldrig kan naae saa
vidt, at den besvarer os de høieste, for Menneskets Existens
allervigtigste Spørgsmaal: er der en Gud? er der et evigt Liv?

Ingenlunde paastaaer jeg, at Videnskaben, som Videnskab, ligefrem skulde
føre til Grundfornegtelse eller til Fornegtelse af
% --- p. 87 ---
det evige Liv. Langtfra! Philosophien kan jo ikke leve uden Idee: naar
den miskjender Ideen, miskjender den sig selv; naar den opløser og
fornegter Ideen, opløser og fornegter den sig selv. For Philosophien er
Ideen som Gud, Gud som Idee: \emph{Ideen er guddommelig}; for
Philosophien er det Evige som Idee, Ideen som det Evige: \emph{Ideens Liv
er et evigt Liv}. Høiere end til denne Forklaring er Philosophien hidtil
ikke kommen, og kan heller aldrig komme.

Philosophiens Idee kan nu paa mangfoldige Maader udvikles, articuleres og
dannes i Begrebssystemernes dialektiske Formationer: anderledes af den
ene Philosoph end af den anden, anderledes af en Aristoteles end af en
Plato, anderledes af en Schelling end af en Hegel. Under disse
systematiske Formationer gjør Idee-Erkjendelsen jo vistnok væsentlige
Fremskridt, forsaavidt Ideen selv klarer sig ud med det udklarede Begreb.
Men Ideens Grundvæsen er og forbliver dog igjennem alle Systemerne det
ene og selvsamme Grundvæsen, ligesom Philosophien jo ogsaa i alle
Philosophier er og forbliver den ene og samme Philosophie. Plato's og
Aristoteles's Gud er ogsaa Schellings og Hegels Gud, det er: den absolute
Idee.

De store Fremskridt, Philosophien i den nyere Tid antages at have gjort
ved Christendommens Bistand, ere illusoriske. Hvad Philosophien har
laant af Christendommen, maa den give tilbage igjen; dette Laan kan
aldrig blive dens Eiendom: Christendommen er ueensartet med
Philosophien. Studeer den hegelske Philosophie, gjennemgaa dette store,
strenge, imponerende Begrebssystem Punkt for Punkt fra Begyndelsen indtil
Enden; forelæg saa Dig selv det Spørgsmaal, hvad Philosophien da nu
endelig her har udgrundet om Guds Tilværelse og Sjælens Udødelighed; hold
Begrebet fast, lad Dig ikke bedaare af Ideens objective
Phantasieglands: Du skal visselig indsee, at den hegelske Gud med alle
sine logiske Trinitetsbevægelser endnu er den gamle
% --- p. 88 ---
% The Greek is set in antiqua italic against the surrounding Fraktur.
aristoteliske Verdensgud, (\textit{τὸ πρῶτον κινοῦν, οὐ κινούμενον}),
hiin første Bevæger, der ikke bevæges, ɔ: den absolute Idee; ja Du skal
see, at Hegel endnu i det 19de Aarhundrede efter Christus ikke er kommen
et Haarsbred videre med Sjælenes Udødelighed, end Plato allerede var i
det 4de Aarhundrede før Christus.

Det Grundlag, hvorpaa Philosophien opfører sine Systemer, er Ideens
Grundlag; men dette Grundlag er tvetydigt, som Ideen selv. Hvad er vel
den philosophiske Idee? Et amphibolisk Væsen! et Væsen, der ikke selv
veed, hvorfra det kommer og hvorhen det vil; et Væsen, der snart viser ud
over sig selv opad mod den personlige Gud og det hiinsidige Rige, men
snart igjen fornegter det Hiinsidige, spotter den personlige Gud, og
erklærer sig selv for den virkelige Verdens Guddom.

Tvetydigheden beroer derpaa, at Ideen tager sig anderledes ud for den
saakaldte rene Tænkning i Forhold til Væsen og Mulighed, end for den
umiddelbare Erfaring i Forhold til Existens og Virkelighed.

Den metaphysiske Tænkning kommer ikke til Ro, før den kommer til Ro i
Gud, som den sig selv og Alt tænkende, sig selv og Alt villende, sig selv
og Alt forklarende Viisdom, af hvem, ved hvem og til hvem alle Ting ere:
Philosophien finder det forsaavidt baade tænkeligt og muligt, at den
absolute Aand er en i sig selv fri og selvbevidst Aand. Ligeledes kan
den philosophiske Idee-Udvikling lægges an i en saa subjectiv Retning, at
det menneskelige Jeg deri erkjender sin indre Uendelighed, sin
oprindelige Kraft til at beseire Tilværelsens, ja Dødens Skranker, og
hæve sig ved den moralske Verdensorden til evig Frihed i stedse stigende
Fuldkommenhed.

Men hvad Philosophien saaledes forjætter os i Forhold til Væsen og
Mulighed, kalder den igjen tilbage, naar Ideen betragtes og prøves i
Forhold til Existens og Virkelighed. Den philosophiske Idee kan ikke
frigjøre sig fra Materiens Skranker,
% --- p. 89 ---
og altsaa heller ikke opløfte sig over denne endelige Verden. Ideen er i
de existerende Ting, som Tingene ere i Ideen. Ideen har ingen
selvstændig Existens: dens Fornuft er Naturens Fornuft, dens Kræfter
Verdens Kræfter; den gjærer i Elementernes Grund, den vaagner til Liv i
Dyret, til Lys og Frihed i Mennesket. Individerne maae døe, Mennesket
saavel som Dyret, kun Ideen lever i alle Slægter: Ideen splitter sig ad i
den virkelige Existens; for den virkelige Existens er der altsaa ingen
personlig Gud og ingen personlig Udødelighed.

Heraf maa det nu vel kunne skjønnes, hvorledes den Religiøse, der vil
støtte sin Tro paa Philosophien og Videnskaben, bliver bedragen. En
personlig Gud, som den, Religionen forkynder, er Videnskaben en
Forargelse; en personlig Udødelighed, som den, Christendommen forjætter,
er Philosophien en Forargelse! Videnskaben holder sig, hvad
Virkeligheden angaaer, ved det Nærværende, det Naturlige; men Religionen
viser hen til det Hiinsidige det Overnaturlige.

Striden mellem den Troende og Fritænkeren kan da nærmere bestemmes
saaledes: Fritænkeren har Ret i den Paastand, at Religionen og
Videnskaben ikke passe for hinanden; men derimod har han stor Uret, naar
han heraf uddrager den Slutning, at Videnskaben kan afskaffe Religionen.
En Aabenbarings Mulighed kan Videnskaben ikke negte; en Aabenbarings
Virkelighed kan Videnskaben ikke bekræfte; en Aabenbarings Sandhed kan
Videnskaben ikke prøve.

Man har meent, at en saadan Adskillelse af Religion og Videnskab, som
den, jeg her foredrager, vilde afstedkomme stor Forvirring, da man jo saa
maatte troe iblinde, og ligesaa godt kunde falde paa at troe den falske
som den sande Aabenbaring. Lutter Udflugter! Man behage blot at besinde
sig lidt alvorligt paa Forskjellen imellem den subjective Erkjendelse og
den objective Viden.
% --- p. 90 ---
Man har meent, at Den, der antog en Aabenbaring uden først at prøve dens
objective Sandhed i objectiv Viden, forsømte en Sandhedspligt, hvorfor
Ethiken burde kræve ham til Regnskab. Et maadeligt Argument, aldenstund
den ethiske Sandhed jo ikke engang selv kan prøves i objectiv Viden.

At ville forud overbevise sig om en videnskabelig, f.\ Ex.\ en
mathemathisk eller en logisk Sandhed, før man antager og anvender den, er
i sin Orden; men at ville forud overbevises om en ethisk Sandhed, før man
antager og anvender den, er Nonsens, efterdi den ethiske Sandhed
refererer sig til Beslutning og Handling. Den rene videnskabelige
Sandhed lader sig godtgjøre ved reen videnskabelig Tænkning: jeg behøver
slet ikke at gjøre Andet end at tænke, for at faae Sandheden beviist. I
den ethiske Sandhed er Gjerningen derimod selv en Deel af Beviset, ja
Beviset selv; men, naar Gjerningen selv er en Deel af Beviset, ja Beviset
selv, saa er det dog vel en reen Daarskab at ville have det fuldstændige
Beviis forud for Gjerningen.

Har man end i dyb videnskabelig Alvor for længe siden glemt Manden, der
ikke vilde gaae i Vandet, før han havde lært at svømme; saa burde man dog
ikke fordybe sig saa dybt i den theologiske Videns Dybheder, at man
aldeles forglemmer, hvad Krigeren synger i „Wallensteins Lager:“

\begin{quote}
Und setzet Ihr nicht das Leben ein,\\
Nie wird euch das Leben gewonnen seyn.
\end{quote}

Er Christendommen høiere end Videnskaben, saa følger det vel af sig selv,
at den evangeliske Historie ogsaa maa være høiere end den videnskabelige
Kritik. Er Christendommen kun for den Troende, saa er den hellige
Historie jo ogsaa kun for den Troende; hvoraf da igjen følger: at
Fritænkeren ligesaa lidt kan gjendrive den Troende, som den Troende kan
gjendrive Fritænkeren.

I Phantasie og Mulighed kunne de vel forstaae hinanden, gaae ind i
hinandens Tankegang, og forsaavidt underhandle og
% --- p. 91 ---
disputere med hinanden; men i Virkeligheden kan den Ene ikke ved noget
objectivt Beviis overvinde den Anden eller drage den Anden over til sig;
Enhver af dem har sit eget existentielle Fodfæste: det er den haarde
Existens, der stiller dem ad. Fritænkeren existerer i Vantroen; hans
Opfattelse af Jesu Liv er aldeles conseqvent, under den Forudsætning, at
Videnskaben er høiere end Religionen; den Religiøse existerer i Troen,
hans Opfattelse af Jesu Liv er aldeles conseqvent under den
Forudsætning, at Religionen er høiere end Videnskaben. Imellem begge
staaer nu Theologen, eller rettere, han staaer slet ikke; han staaer
idetmindste ikke fast (thi han har jo intet existentielt Fodfæste), men
rokker og halter til begge Sider, saa mod Troen og saa mod Videnskaben.
Theologens Opfattelse af Jesu Liv er en principløs Opfattelse i lutter
Halvhed og Inconseqvens; denne Halvhed og Inconseqvens kalder jeg
Middelmaadighed.

En Fremstilling af Jesu Liv kan nu lægges an saaledes, at Fritænkeren med
sin „modnere Bevidsthed“, den Religiøse med sin „Evangelietro“ og
Theologen med sin „Middelmaadighed“ forenes om at gjennemgaae den
evangeliske Historie og at udtale sig hver i sin Aand over ethvert af
Evangelierne. En saadan Fremstilling bliver da at charakterisere som en
existentiel Fremstilling, forsaavidt den nemlig ikke gaaer ud paa en
videnskabelig Sammenstilling af forskjellige Theorier, men en levende
lidenskabelig Strid mellem forskjellige Charakterer.

Den existentielle Fremstilling af Jesu Liv maa i Ordets strengeste
Forstand være kritisk; men den Kritik, der her kommer til Anvendelse, er,
vel at mærke, ikke en videnskabelig, men en existentiel Kritik. Kritik
kommer jo af Krise, og det vil dog vel Ingen negte, at Existensen har
sine Kriser. Den Aand, der sætter de stærkeste Modsætninger i Bevægelse,
er da ogsaa den høieste Aand, en Aand, der bevirker de dybeste Kriser.
Naar da nu et Men%
% --- p. 92 ---
neske træder op og forkynder sig som Guds eenbaarne Søn: hvilke Kriser
maa en Saadan ikke fremkalde!

Saaledes fremtræder Christus: et Menneske i udvortes Ringhed, og dog et
Menneske med guddommelig Myndighed; et Menneske, der gjør de Seende
blinde, og dog et Menneske, der gjør de Blinde seende; et Menneske, der i
denne Verden ikke har Det, hvortil han kan hælde sit Hoved, og dog et
Menneske, der siger til den Døde: stat op! og til den Værkbrudne: dine
Synder ere Dig forladne! Hvo er saa dette Menneske? --- Den aabenbarede
Gud? Den aabenbare Bedrager? Her begynder Krisen.

% printed footnote mark: *)
Den evangeliske Historie bærer den existentielle Kritik i sig selv: den
viser os Christus i Forhold saavel til de Ikke-Troende, som til de
Troende.\footnote{Jfr.\ Evangeltr.\ og den mod.\ Bevidsthed.\ S.\
210---227.} De Ikke-Troende see ham strax som en Bedrager; de Troende see
ham som Guds Søn; den store Hob skifter Mening. Disse Kriser og denne
Kritik ledsage Christus uafbrudt gjennem hele hans offentlige Liv, fra
Galilæa til Judæa, fra Jerusalem til Nazaret, fra Capernaum til Golgatha.
Det var Mænd af Galilæa, der bleve hans første Disciple og troede ham paa
hans Ord, at de fra nu af skulde see Himmelen aaben, og Guds Engle fare
op, og fare ned over Menneskens Søn; det var Mænd af Galilæa, der i
Nazarets Synagoge forargedes paa ham, og stødte ham ud af Staden for at
styrte ham ned fra Skrænten. Det var Pharisæer og Skriftkloge, der
hadede ham som en Sværmer og Folkeforfører; det var en Pharisæer og
Skriftklog, der kom til ham i Nattens Stilhed som til en Israels Lærer.
Det var en Røver paa Korset, der bespottede ham, og det var en Røver paa
Korset, der bekjendte ham. Det var en Romer, der dømte ham i Jerusalem,
da Jøderne rasede: Korsfæst, Korsfæst! og det var en Romer, der
frikjendte ham paa Golgatha, der Gravene brast, og Jorden skjælvede.
% --- p. 93 ---
Det Forhold, hvori Christus ved sin første Optræden sætter sig til sine
Samtidige, er tillige bestemt som det principielle Grundforhold, hvori
han sætter sig til den hele Menneskeslægt gjennem alle Tider. Paa den
ene Side følge ham da de Forargede, der snart ville stene ham, snart
støde ham ned, og altid spotte ham, fordi hans Ord er dem en haard Tale,
som Ingen kan høre. Paa den anden Side følge ham de Troende, der
bekjende hans Navn, og ikke kunne forlade ham, fordi de i Tidernes
Forvirring ikke kjende nogen anden Udvei, end Petrus, der han sagde:
\emph{Herre, til hvem skulle vi gaae hen? Du haver det evige Livs Ord.}
Om disse Charakterer bevæger sig Mængden, den bølgende Mængde, der
kommer og gaaer; den uselvstændige Mængde, der snart er rørt til
Opvækkelse, snart ophidset til Forargelse; den urolige Mængde, der raaber
og raaber, idag: Hosianna! imorgen: Korsfæst!

Her have vi, som i et Skyggerids, de oprindelige Forudsætninger og
Betingelser, paa hvilke Evangelietroens existentielle Kritik fremdeles
maa lægges an. Og nu maa Verden gaae saa langt, den vil, og blive saa
gammel og saa oplyst, og saa klog og saa videnskabelig, den vil: den
kommer dog aldrig ud over disse Kriser, den kommer aldrig ud over denne
Kritik.

Den existentielle Fremstilling af Jesu Liv bør efter Evangelisternes
Anviisning vel ogsaa være pragmatisk. Men den Pragmatisme, hvorom det
her gjælder, er ikke den videnskabelige. I den videnskabelige
Pragmatisme bliver det Enkelte bestandig sat som flygtigt Moment i det
Totale. Herved faaer Endelighedens Schema en saa afgjørende Overvægt, at
Aanden forsvinder. Saa bliver det Hele udstykket efter de forskjellige
Reiser og de forskjellige Fester; saa kunne Theologerne ikke enes om
Chronologien; saa strider man om Synoptikerne, saa kives man om
Bjergprædikenen, saa raisonnerer man om „Paaskeberegnin%
% --- p. 94 ---
gen“; saa har den lærde Videnskab glemt Troen, og bevæger sig i lutter
Endelighed.

I den evangeliske Pragmatisme bliver derimod det Totale sat som
concentreret i det Enkelte. I hver enkelt forløsende Gjerning er den
hele Forløser selv tilstede, i hvert enkelt Evangelium den hele
evangeliske Historie. Den endelige Tidsfølge bliver herved nedsat til et
forsvindende Moment i Grupperingen af de centrale Begivenheder. De
overnaturlige Phænomener og overordentlige Gjerninger i Jesu Liv følge
paa hinanden saa pludselig, saa overraskende, og sprænge det Timeliges
ordinære Sammenhæng saa aldeles, at en endelig Pragmatisme bliver umulig.
% printed footnote mark: *)  --- the note runs on from p. 94 to p. 95.
Ingen Discipel kunde forholde sig objectiv til disse
Begivenheder,\footnote{„Historisk Øienvidne altsaa har den samtidige
Lærende let ved at blive; Ulykken er imidlertid, at det at vide en
historisk Omstændighed, ja at vide dem alle med Øienvidnets
Tilforladelighed, ingenlunde gjør Øienvidnet til Discipel, hvad man jo
kan see deraf, at denne Viden betyder for ham ikke andet end det
Historiske. Det viser sig strax her, at det Historiske i concretere
Forstand er ligegyldigt; vi kan lade Uvidenheden indtræde i Forhold
dertil, og lade Uvidenheden ligesom tilintetgjøre det Historiske; naar
blot Øieblikket endnu er tilbage, som Udgangspunkt for det Evige, er
Paradoxet tilstede. Dersom der var en Samtidig, som selv havde
indskrænket sin Søvn til den korteste Tid for at følge hiin Lærer, hvem
han fulgte uadskilligere, end den lille Fisk, som følger Haien; dersom
han holdt hundrede Spioner i sin Tjeneste, der overalt belurede hiin
Lærer, og med hvilke han selv confererede hver Aften, saa han vidste hiin
Lærers Signalement, indtil det Mindste, vidste, hvad han havde sagt, hvor
han havde opholdt sig til enhver Time paa Dagen, fordi hans Iver lod ham
betragte endog det Ubetydeligste som vigtigt, var en saadan Samtidig
Discipelen? Ingenlunde. Han kunde vadske sine Hænder, hvis Nogen vilde
sigte ham for historisk Upaalidelighed, men mere heller ikke. Dersom der
var en Samtidig, der havde opholdt sig i fremmede Lande, og først vendte
hjem, da hiin Lærer endnu kun havde een Dag eller to tilbage af sit Liv;
dersom den Samtidige atter ved Forretninger blev afholdt fra at faae hiin
Lærer at see, og først kom til i det Sidste, da han vilde opgive Aanden,
var denne historiske Uvidenhed til Hinder for, at han kunde være
Discipelen, naar Øieblikket for ham var Evighedens Afgjørelse?“ Johannes
Climacus. Philosophiske Smuler. S.\ 85---86.} ingen Tilskuer kunde
smaalig optegne hvert enkelt Træk,
% --- p. 95 ---
ingen Apostel kunde, saa at sige, holde Dagbog over Jesu Gjerninger (jfr.\
Joh.\ 20, 30---31. 21, 25).

Den existentielle Fremstilling af Jesu Liv kan vistnok ikke være
philosophisk, men maa dog anlægges saaledes, at den evangeliske Historie
sees i Idealets Lys, som en over Fortid og Nutid ophøiet, evig og hellig
Historie.

% PRINTER'S DEFECT: the print reads „Fremstilligen“ (missing n) for
% „Fremstillingen“. Transcribed as printed.
Fremstilligen af den hellige Historie har da den særegne Bestemmelse at
kaste Lys over Verdens Historie og gjøre Idealerne kjendelige, at Enhver
maa dømme sig selv, naar han dømmer om det aabenbarede, Hjerternes Tanker
aabenbarende, Ideal.

Er der intet saligt Liv, da er det vistnok større at indtage Staden, end
blot at herske over sit Sind, og da er ogsaa Christus en Forargelse; men
er der et saligt Liv: ja, da er jo Den større, der hersker over sit Sind,
end Den, der indtager Staden, og da er Christus Verdens Frelser. Er der
i Troen intet evigt Liv, da maa den Enkelte søge sit Ideal i det
verdslige Samfund, og det verdslige Samfund sit Ideal i Slægtens
Historie. Historiens Aand aabenbarer da Rigets og Magtens og Ærens
allerhøieste Ideal, naar Staden indtages, naar det heroiske Genie siger:
\textit{jacta alea esto}, og gaaer over Rubicon med sine Legioner. Men
er der i Troen et evigt Liv; maa den Enkelte kjøbe det med
Selvfornegtelse, bevare det i Taalmodighed, og fuldende det i
selvopoffrende Kjærlighed: da er den Tornekronede vel ogsaa høiere end
den Laurbærkronede; da gaaer Idealets Vei heller ikke over Rubicon til
Capitolium, men Veien gaaer over Kedrons Bæk til Gethsemane Have; og da
er Beslutningens høieste Løsen jo
% --- p. 96 ---
heller ikke et \textit{jacta alea esto}, men det høieste Løsen er i
Bønnens Ord: Fader, ikke min, men din Villie!

% PRINTER'S DEFECT: the print reads „er der er saligt Liv“ for
% „et saligt Liv“ (cf. p. 95, „er der et saligt Liv“). As printed.
Hver dømme herom, som han vil: han dømmer sig selv. Men, dette staaer
fast: er der er saligt Liv, saa er Christus Verdens Frelser; er Christus
Verdens Frelser, saa er Christendommen høiere end Videnskaben; men er
Christendommen høiere end Videnskaben, saa er \emph{Evangelietroen jo
ogsaa --- ueensartet med Theologien.}

%% ============================================================
%%  OTTENDE FORELÆSNING (pp. 97--111 = PDF 112--126)
%% ============================================================
\chapter*{Ottende Forelæsning}
\addcontentsline{toc}{chapter}{Ottende Forelæsning}
\label{ch:forel8}
\markboth{Ottende Forelæsning}{}

\begin{center}\itshape
Det halvphilosophiske Christusbillede, den speculerende Dogmatiker
fremtryller, bestaaer i en mystisk Sammensmeltning af det Videnskabelige
og det Religiøse, og er desaarsag saa uendelig forskjelligt fra
Evangelisternes Christus, at det endog falder udenfor den religiøs-ethiske
Synskreds: Evangelietroen er ueensartet med Theologien.
\end{center}

% --- p. 97 ---
Endskjøndt Theologerne naturligviis ikke kunne gaae ind paa, hvad de i
Forhold til Nutidens Oplysning maae benævne som „det krasse
Inspirationsprincip“, og desaarsag ere i megen Tvivl om den evangeliske
Histories Enkeltheder og de fire Evangeliers Ægthed: saa trøste de sig dog
med, at den fromme Betragtning hos enhver af Evangelisterne dog endnu
gjenfinder de samme historiske Hovedtræk, den samme religiøse
Grundanskuelse, det samme Christusbillede. Den theologiske Stemme er aldrig
mere bevæget og fuldtonende, end netop naar der tales om Hovedtanken, om
Grundanskuelsen, om Christusbilledet. Selv da, naar man kommer i
Forlegenhed med alt det Enkelte, kan man jo freidig holde paa
Grundanskuelsen; selv da, naar Videnskaben har forvandlet Evangelisternes
Christus til en halvphilosophisk Phantasie, seer man dog altid et
Christusbillede. --- Medens den kritiserende Fortolker fornemmelig
interesserer sig for de historiske Enkeltheder i Jesu Liv, interesserer den
speculerende Dogmatiker sig i Særdeleshed for Ideen, for Grundanskuelsen,
for Christusbilledet.

% --- p. 98 ---
I Ideen og Grundanskuelsen mener nu Theologen at forsone Christendommen med
Videnskaben, idet han af Christendommens eget Indhold udvikler en
speculativ Erkjendelse, en Erkjendelse, der i alle væsentlige Punkter
formeentligen skal tilfredsstille Troen.

For at opdage den Vildfarelse, hvorpaa dette i Kirken ældgamle
Mediationsforsøg beroer, ville vi her dvæle lidt nærmere ved den tidt
berørte Modsætning af den videnskabelige og den existentielle Erkjendelse.

I den videnskabelige Erkjendelse har den Erkjendendes særegne Stemning og
subjective Tilstand Intet at betyde; i den existentielle Erkjendelse beroer
Alt paa det erkjendende Subjects eiendommelige Tilstand og Stemning. I den
videnskabelige Erkjendelse maa den Erkjendende abstrahere fra sig selv,
forglemme sig selv, give sig hen i Betragtningens Ro, og forholde sig
almindelig til den almindelige Gjenstand. I den existentielle Erkjendelse
maa det erkjendende Subject gaae i sig selv, stille sig i et afgjørende
Forhold til Gjenstanden, hvorfor Enhver da ogsaa i dette Subjectivitetens,
Inderlighedens og Interesserethedens personlige Forhold med Rette kalder
Gjenstanden sin Gjenstand. Den videnskabelige Erkjendelse, der bevæger sig
i det Almindelige, er lidenskabsløs, og bør være det, efterdi den er
upersonlig; den existentielle Erkjendelse, der bevæger sig om det
Punktuelle, det Enkelte, det Egne, er lidenskabelig, og maa være det i
samme Grad, som den er personlig.

Heraf følger da, at den Erkjendelse, der vindes i Videnskaben, kun har en
ligefrem Betydning for Videnskaben, og maa forandre sin Charakteer, naar
den anvendes paa Existensen; og omvendt, at den Erkjendelse, der opnaaes ad
Existensens Vei, kun har en ligefrem Betydning for Existensen selv, og maa
forandre sin Charakteer, naar den oversættes i Videnskaben. Da nu disse
tvende Erkjendelser paa denne Maade ere saa ueensartede, at de ikke kunne
ombyttes med hinanden; saa følger igjen
% --- p. 99 ---
heraf, at de hverken kunne forholde sig ligefrem til hinanden eller
ligefrem corrigeres ved hinanden. Den videnskabelige Erkjendelse kan ikke
ligefrem corrigeres ved Existens, saa lidt som den existentielle ved
Videnskab.

Hvad her er sagt, har (hvad allerede Aristoteles paaviser) en stor
Betydning for den rationelle Ethik, men faaer dog en endnu høiere og ganske
anderledes afgjørende Betydning for den særegne Sphære, vi i strengere
Forstand maae kalde Evangelietroens Sphære. Hvad Christendomserkjendelsen
angaaer, maae vi paa det Bestemteste fastholde den Regel, at Videnskab ikke
kan corrigeres ved Tro, saa lidt som Tro kan corrigeres ved Videnskab.

Anvende vi nu denne Regel paa vort foreliggende Thema, da vil ogsaa
Vildfarelsen med det omtalte Mediationsforsøg uden Tvivl komme for Dagen.

For at udforme den christologiske Idee og construere det saakaldte
Christusbillede, gaaer den speculerende Dogmatiker tilværks paa følgende
Maade. De christologiske Hovedsætninger, paa hvilke den dogmatiske
Ideeudvikling skal begrundes, laaner han ikke fra nogen udenfor Kirken
staaende Philosophie, han henter dem allesammen fra Skriften og Kirkelæren.
Da nu enhver Hovedsætning forsaavidt maa siges at være en Troessætning,
mener Dogmatikeren hermed at tilfredsstille Troen, at bygge paa Troens
Grund, at erkjende i Troen og ud af Troen. Men hvilken Erkjendelse er det
saa, han her tilstræber? Ja det skal jo være en videnskabelig og, saavidt
muligt, endogsaa speculativ, idetmindste halvspeculativ Erkjendelse. Nu
begynder Vildfarelsen. De kirkelige Troessætninger omsættes i den objective
Videns Form, og behandles som almindelige Videnskabssætninger; idet disse
Videnskabssætninger sammenarbeides til en videnskabelig Totalitet,
udformes saa den christologiske Idee, og Christusbilledet viser sig.

At en saadan Christologie ingenlunde forsoner Troen med
% --- p. 100 ---
Viden, men meget mere viser sig som en tvetydig Sammensmeltning af
ueensartede Bestanddele, en Sammensmeltning, der i lige Grad strider imod
Troens og Videns Natur, vil uden Vanskelighed kunne bevises.

Hvad Dogmatikerens halvspeculative Christologie angaaer, maa det strax
vække Betænkelighed, at den gjør Fordring paa en dobbelt Anerkjendelse,
idet den baade vil anerkjendes som et Product af Troen og dog tillige
anerkjendes som et Product af Videnskaben; den er som et Afkom, der har
Troen til Moder og Videnskaben til Fader, og nu med denne sin dobbelte
Natur gjør en dobbelt Fordring, een paa Moderens og een paa Faderens Vegne.
Med denne dobbelte Fordring skulde man da nu ogsaa vente, at den
christologiske Idee var villig til at underkaste sig en dobbelt Prøvelse,
nemlig en Prøvelse i Troen og en Prøvelse i Videnskaben; men med denne
Prøvelse bliver det dog alligevel ikke til Noget. Lægge vi nemlig Troens
existentielle Maalestok an paa den christologiske Idee, saa undviger den
Prøvelsen ved at henvise til Videnskaben; ville vi dernæst anlægge
Videnskabens strenge Maalestok, saa undviger den igjen Prøvelsen, og
henviser os til Troen. Paa denne Maade kan Dogmatikerens christologiske
Idee i Grunden aldrig blive prøvet, med mindre det da skulde være for en
Domstol, der er høiere end baade Troen og Videnskaben; men hvor er saa
denne Domstol?

Den christologiske Idee tilfredsstiller ikke Videnskaben; thi hvad
Videnskaben skal anerkjende, maa den ogsaa kunne prøve; men den
christologiske Idee kan Videnskaben ikke prøve. Da denne Idee nemlig har
sit Udspring ikke fra det naturlige Fornuftsystem i det Nærværende, men fra
Aabenbaringens hiinsidige Rige, er den af mirakuløs Oprindelse og har en
Bestanddeel i sig, hvorom Videnskaben Intet kan sige, efterdi den ikke kan
prøve den, nemlig den overnaturlige Bestanddeel. --- Christus blev
Menneske, der han blev undfangen af den Hellig Aand og født
% --- p. 101 ---
af Jomfru Maria. Den halvspeculative Christologie kan nu vistnok fremsætte
Adskilligt i Almindelighed om, hvad den kalder Incarnationens „kosmiske
Under“; men om \emph{denne} Incarnation, om \emph{denne} Undfangelse og
\emph{denne} Fødsel kan Videnskaben dog alligevel Intet afgjøre, efterdi
Videnskaben ikke paa nogen Maade seer sig istand til at prøve Miraklet.
--- Christus er sand Gud og sandt Menneske. Den halvspeculative
Christologie kan nu fremsætte Adskilligt om Gudmennesket i Almindelighed,
om Christi dobbelte Natur og dobbelte Villie i Almindelighed; men om den
virkelige Eenhed af de to Naturer, hvorledes det gik til, at Christus
virkelig kunde lide og døe paa Korset, endskjøndt han dog selv i Lidelsens
og Dødens Øieblik vedblev at være den sande Gud, den Gud, der bærer Himmel
og Jord, derom kan Videnskaben ikke dømme, derom kan den hverken yttre sig
bekræftende eller benegtende. Men kan Videnskaben ikke bedømme
Christologiens Indhold, hvorledes skulde den saa kunne vedkjende sig den
christologiske Fremstilling. --- Christus opstod fra de Døde, opfoer til
Himmelen, og satte sig hos Faderens høire Haand. Der var en Tid, da
Videnskaben for Alvor meente, at kunne afgive en Dom om denne Christi
Ophøielse: ak, hvilke heftige Stridigheder have Theologerne ikke ført om
Christi Legems Allestedsnærværelse, om \textit{ubiquitas absoluta} og
\textit{ubiquitas respectiva}, og hvilke Forargelser have disse
Stridigheder ikke vakt! Naar slige theologiske Stridigheder ikke længer
kunne sætte Videnskaben i Lidenskab, skulde dette da ikke være et Beviis
for, at man nu endelig er kommen til Bevidsthed om, at slige Ting ligge
udenfor den menneskelige Videns Grændser; men hvad der ligger udenfor vor
objective Videns Grændse, kan jo dog umulig behandles videnskabelig.

Den christologiske Idee tilfredsstiller da heller ikke Troen, og det af den
simple Grund, at enhver Troessætning taber sin Charakteer, naar den
omsættes i Videns Form. Den saakaldte
% --- p. 102 ---
Forsoning af Tro og Viden begynder uvilkaarlig med den Grundfeil, at stille
Troen paa lige Trin med Viden, og bevæger sig saaledes i en Cirkel.

At tale om Viden som om et fritstaaende, af den Vidende uafhængigt Begreb,
er i sin Orden, da Viden er objectiv; men at behandle Troen som et
almindeligt, i sig begrundet, af den Troende uafhængigt Væsen, er
urimeligt, efterdi Troen er subjectiv. Troen er kun for den Troende; den
Troende er et virkeligt, i Subjectivitetens Inderlighed existerende
Menneske; Troens Kategorier er Existensens høieste Kategorie, dens
Erkjendelse existentiel. Begaaer man nu alligevel det Misgreb at abstrahere
fra den Troende, naar man taler om Troen, ligesom man abstraherer fra den
Vidende, naar man taler om Viden; saa faaer man rigtignok ud, at Troen er
objectiv ligesaavel som Viden; men derhos har man da med det Samme aflivet
Troen, thi Troen døer i Objectiviteten.

At Viden er objectiv, kjendes derpaa, at den grunder paa almeengyldige
Forudsætninger, paa objective Beviser, og vil Intet have med de subjective
Grunde at gjøre; at Troen er subjectiv, sees deraf, at den ubetinget
grunder i Subjectivitetens Inderlighed, og derfor slet ikke vil vide Noget
af de objective Grunde. Som nu Den, der mener at hjælpe paa Videnskaben ved
at indblande subjective Grunde i de objective Beviser, krænker og forqvakler
Videnskaben; saa krænker og forqvakler man igjen Troen, naar man forsøger
paa at understøtte Subjectivitetens afgjørende Grund ved objective Beviser.

For at retfærdiggjøre sin halvspeculative Christologie, er Dogmatikeren ---
charakteristisk nok --- i det Hele taget mere opmærksom paa Brevene end paa
Evangelierne, og, hvad Evangelierne angaaer, mere opmærksom paa Johannes
end paa Synoptikerne. Grunden hertil er indlysende nok. Evangelierne, især
de „synoptiske“, holde sig paa alle Punkter i Existensen, medens
% --- p. 103 ---
Brevene derimod, der fremsætte Læren i større Almindelighed, synes at komme
det videnskabelige Foredrag langt nærmere. I Brevene altsaa, især i de
paulinske Breve, har man da meent at finde Typen for en speculativ
Christologie, og Apostelen Paulus er paa denne Maade bleven opstillet som
Forbillede for speculative Dogmatikere.

At der nu virkelig hos Apostelen Paulus, ligesom ogsaa hos Evangelisten
Johannes, forekomme mange Sætninger og Udtryk, som den dogmatiserende
Speculation paa sin Viis kan tage til Indtægt, skal aldrig negtes; men
heraf følger endnu ikke, at disse johanneiske eller paulinske Sætninger
selv hvile paa speculativ Grund.

Den paulinske Christuslære er ligesaa lidt speculativ, som den johanneiske
Logoslære er speculativ. Hvad Johannes og Paulus i denne Retning udtale,
udtale de ikke i Videnskabens, men i Troens Navn, med Troens subjective
Forvisning, med Inspirationens guddommelige Myndighed. Naar det hedder i
Johannes' Evangelium: „I Begyndelsen var Ordet, og Ordet var hos Gud, og
Ordet var Gud. Det var i Begyndelsen hos Gud. Alle Ting ere blevne ved det,
og uden det er ikke een Ting bleven til af det, som er“ . . . .; naar det
endvidere hedder (V. 14): „Ordet blev Kjød og boede iblandt os“ o. s. v.
--- da lyde disse Sætninger vistnok speculativt i et speculativt Øre; men
see vi lidt nærmere til, vil det snart vise sig, at slige Udsagn hverken
ere eller kunne være speculativt meente.

Skulde den johanneiske Logoslære have speculativ Betydning, da maatte den
vel være anlagt saaledes, at den lod sig retfærdiggjøre mataphysisk; men
% PRINTER'S ERROR (verified on the page image at 300 dpi): the 1850 print
% really reads „mataphysisk“ for „metaphysisk“. Transcribed as printed.
dette er aldeles ikke Tilfældet. Den Sætning: \emph{Ordet blev Kjød}, er
ingen videnskabelig Sætning og kan aldrig blive det. „Ordet blev Kjød“: det
Overnaturlige paatog sig vor Natur, det Overhistoriske blev historisk; en
saadan Sammensætning af saadanne ueensartede Faktorer kan
% --- p. 104 ---
Videnskaben ikke magte, og hvad Videnskaben ikke kan magte, vil den heller
ikke bedømme. Skulde Videnskaben afgjøre Noget om den Sætning:
\emph{Ordet blev Kjød}, da maatte Loven for det Overnaturliges
Vexelvirkning med det Naturlige først være opdaget (ligesom
Naturvidenskaben jo har opdaget Love for de naturlige Tings Vexelvirkning);
men denne Lov er Videnskaben et reent X, en aldeles ubekjendt Størrelse:
% the X is set in the surrounding Fraktur fount, not antiqua, so no \textit{}.
her kan ingen videnskabelig Prøve anstilles; men hvad Videnskaben ikke kan
prøve, vil Videnskaben heller ikke bedømme. Skal den johanneiske Sætning
begrundes videnskabelig, da seer det kun sørgeligt ud, saavel med Johannes'
Evangelium, som med Christendommen; thi Begrundelsen er virkelig aldeles
umulig. Men er den videnskabelige Begrundelse i dette Tilfælde aldeles
umulig, saa er den vel ogsaa, hvad Troen angaaer, aldeles overflødig.

Den johanneiske Sætning er en existentiel Aabenbaringssætning. Naar
Johannes siger: „Ordet blev Kjød og boede iblandt os“, siger han det i
Troens Navn; Evangelistens Tro kan ikke adskilles fra Evangelisten selv:
Troen kan ikke adskilles fra den Troende. Idet Evangelisten taler ud af sin
Tro, taler han ikke om Troen i Almindelighed, men han taler som Troende til
Troende: Forholdet er existentielt, ikke videnskabeligt. De Andre, der
ville overbevise sig om Sandheden i Evangelistens Ord, kunne da heller ikke
overbevises ved objective Grunde i Almindelighed. Den eneste Grund, hvorpaa
Overbeviisningen kan bygges, er den, at Evangelisten er inspireret; men, at
Evangelisten er inspireret, kan heller ikke bevises ved objective Grunde:
det maa troes. Skal jeg troe paa Evangelistens Ord, da maa jeg troe, at
Evangelisten er inspireret; skal jeg troe, at Evangelisten er inspireret,
maa jeg ogsaa troe paa Evangelistens Ord: Forholdet imellem Evangelisten
selv og alle Dem, der ved Evangelistens Ord skulle komme til Troen, er
igjen existentielt, ikke videnskabeligt.

% --- p. 105 ---
Heraf indsees da tillige, hvad man skal dømme om de christologiske Steder,
der findes i Apostlenes Breve. Allerede den Maade, hvorpaa disse Breve ere
blevne til, maa jo i høieste Grad advare os mod videnskabelige
Misforstaaelser. Hvorledes kommer vel et saadant Brev istand? Der opkommer
Forvirring i en Menighed; selvraadige Mennesker, usædelige Personer, falske
Lærere vække Forargelse. Efterretningen herom naaer den fraværende Apostel,
og vækker hans Bekymring. Med den Bekymredes Nidkjærhed, med Blikket fæstet
paa de stedfindende Forargelser, med al den Ilfærdighed, som Øieblikkets
Fare gjør nødvendig, sætter Apostelen sig nu hen og skriver sit Brev. I
dette Brev, hvis Hensigt er at straffe, advare, formane, belære, opmuntre
og trøste, indslynges nu de christologiske Aabenbaringssætninger i den
praktiske Underviisning, saaledes at det Ene ligesom uvilkaarlig flyder af
det Andet; som f.~Ex., naar Paulus (Phil. 2, 1--13) vil indskjærpe
Philippenserne Ydmyghed, og han da (V 5--13) siger: „Thi det samme Sind
% PRINTER'S DEFECT: the second reference is printed „(V 5--13)“, without the
% period after V that „(Phil. 2, 1--13)“ has. Verified by zoom; left as printed.
være i Eder, som og var i Christo Jesu, hvilken, der han var i Guds
Skikkelse, ikke holdt det for et Rov at være Gud lig, men forringede sig
selv og tog en Tjeners Skikkelse paa, og blev Mennesker lig; og da han var
funden i Skikkelse som et Menneske, fornedrede han sig selv, saa han blev
lydig indtil Døden, ja Korsets Død. Derfor haver og Gud høit ophøiet ham og
skjenket ham et Navn, som er over alt Navn; at i Jesu Navn skal hvert Knæ
bøie sig, deres i Himmelen og paa Jorden og under Jorden, og hver Tunge
skal bekjende, at Jesus Christus er en Herre til Gud Faders Ære. Derfor,
mine Elskelige,“ o. s. v.

Men hvad Apostelen saaledes skriver i Øieblikkets Fare, i Bekymringens
Angst, i Kjærlighedens Iver, saa at hans Hjerte brænder; see, det gjøre
Theologerne saa i al videnskabelig Ro og Mag til Gjenstand for objectiv
videnskabelig Betragtning: mærker
% --- p. 106 ---
man da ikke Misforholdet! Hvad Troens Apostel skriver til Troende til
Troens Fremme i Troens Navn, see det mene Theologerne saa at kunne behandle
som halvspeculative Bidrag til en halvspeculativ Forstaaelse af det
halvspeculative Forhold mellem Christi guddommelige og menneskelige Natur
(κρύψις eller κένωσις?), mellem Gudmenneskets Fornedrelse og Gudmenneskets
% „eller“ is set in antiqua between the two Greek words; transcribed plain.
Ophøielse (\textit{status exinanitionis} og \textit{status exaltationis});
mærker man da ikke Misforholdet!

Maa nu allerede den apostoliske Brevform og de særegne Vilkaar, hvorunder
disse Breve ere blevne til, advare os imod den speculativ-videnskabelige
Opfattelse; hvormeget mere maae vi da ikke advares, naar vi betænke, at
disse Breve ere udstrømmede af Inspirationens Væld. Den Aand, der
inspirerer den videnskabelige Forfatter, er ikke den apostoliske. Hvis en
inspireret Apostel opstod i Nutiden, vilde han vist faae Andet at gjøre,
end at være Forfatter, og da lære os noget ganske Andet, end hvad vi kunne
lære af videnskabelige Bøger. Hvor Religionens Aand har en Kraft, som den
jo maa have i en Apostels Liv, er det Religiøse altfor nidkjært paa sig
selv og altfor iversygt paa sin Forkyndelse til, at den Inspirerede skulde
faae Ro til videnskabelige Betragtninger og æsthetisk Ideenydelse. Et
apostolisk Ord er en apostolisk Gjerning; en Apostel har sandelig ikke Tid
til at bygge Theorier; selv naar han skriver sit længste Brev, og meddeler
Læren i den udførligste Form, er hvert Ord grebet ud af hans eget Liv, og
Meddelelsen Punkt for Punkt existentiel.

I den halvspeculative Dogmatik betragtes nu Incarnationen og Inspirationen
som de to store Grundundere, der betinge og forklare hinanden gjensidig.
Den, der ikke antager Inspirationen, kan heller ikke fatte Incarnationen;
som man dømmer om Christi Apostle, saa dømmer man om Christus selv: det
sande Christusbillede maa sees i Belysning af den apostoliske Aand. Hvad
Dogmatiken saaledes siger i Almindelighed, maae vi ogsaa indrømme
% --- p. 107 ---
i al Almindelighed; men denne almindelige Indrømmelse fører os dog
vel endnu ikke ind i speculative Betragtninger?

Som Inspirationsprincipet er et overvidenskabeligt Existentialprincip, der
kun gjælder for den Troende i Troen, saaledes er ogsaa
Incarnationsprincipet et overvidenskabeligt Existentialprincip, og
Forholdet imellem begge et existentielt, ikke et speculativt Forhold.

Den existentielle Christuslære, der kan udledes af Apostlenes Breve, er i
alle Punkter overeensstemmende med den Christuslære, der ligger i
Troesbekjendelsen og Evangelierne; men den halvspeculative Christologie,
Theologerne abstrahere af Skriften og Kirkelæren, vil aldrig ret stemme med
Evangeliernes Christus.

Den theologiske Christologie vil nu t. Ex. enten slet ikke lade Christus
fristes, fordi hans Fristelse i Ørkenen forekommer den saa urimelig (i
dette Tilfælde trodser Dogmatikeren da ligefrem Evangelisterne); eller
ogsaa den videnskabelige Christologie vel indrømmer en Fristelse, men uden
at indlade sig nærmere paa den Fristelse, hvorom Evangelisterne tale. Den
evangeliske Fristelseshistorie er slet ikke lagt an paa videnskabelig
Behandling. Hvad en halvspeculativ Christologie har at sige os om Christi
Fristelse i Almindelighed, er derimod, naar det siges allerdybsindigst,
just beregnet paa at gjøre et videnskabeligt Indtryk. Men dette
videnskabelige Indtryk, hvor forskjelligt er det dog ikke fra det
forunderlige, frastødende, forargelige Indtryk, vi fornemme, naar vi med et
„kritisk Øie“ læse Evangeliet om Fristelsen i Ørkenen.

Den, der læser den halvspeculative Christologie, faaer jo et Indtryk, som
om den menneskelige Forstand nu var i den bedste Forstaaelse med den
aabenbarede Kjendsgjerning, som om det nu Altsammen, selv det
Allervidunderligste, selv Fristelsen i Ørkenen, var ganske nær ved at gaae
op i Videnskabens christologiske Idee; Den derimod, der hengiver sig til
Indtrykket af den evangeliske Fristelseshistorie, føler uvilkaarlig, at han
enten maa forkaste
% --- p. 108 ---
denne Fristelse som en absurd Phantasie, eller ogsaa tage Fornuften fangen
under Troens Lydighed.

Gjør nu Evangeliet paa alle Punkter det Indtryk, at der skal troes med
Forstand \emph{mod} Forstand, medens Dogmatiken paa alle Punkter bestræber
sig for at frembringe det Indtryk, som om Forargelsens Braad var forsvunden
i Videnskaben, saa at der egentlig nu paa Speculationens nærværende Trin
kun skulde troes \emph{med} Forstand: da er det vel indlysende nok, at
Evangeliet og Dogmatiken arbeide --- ikke med hinanden, --- men imod
hinanden.

Derfor have Theologerne da ogsaa, ikke blot af historisk-kritiske, men
tillige af metaphysisk-dogmatiske Grunde altid saa meget at rette ved de
fire Evangelister; det er saa saare sjeldent, at en Evangelist er saa
heldig, at kunne gjøre dem det ret tilpas. I den nyere Tid har Evangelisten
Johannes t. Ex. maattet høre ilde, fordi han, tvertimod Christologiens
videnskabelige Regel, lader Christus bede høit ved Lazari Grav. „Das
Unpassende dieser Wendung“ --- siger De Wette --- „muß man anerkennen . . .
. Denn in der That dürfen und müssen wir uns Jesu Bewußtseyn von Gott und
seiner Stellung zu ihm als ein ununterbrochenes, nicht an einzelne
Gebetsmomente geknüpftes denken.“%
% printed footnote mark: *)
\footnote{Jfr. Evangltr. og den mod. Bev. S. 438--58.}

Men sæt endog, at Dogmatikerens halvspeculative Christusbillede var i alle
Hovedtræk saa ligt Evangelisternes Christus, som en Copie kan ligne sin
Original, saa vare vi endda ligenær. Det dogmatiske Christusbillede falder
ind under en ganske anden Synskreds, og er aldeles ueensartet med den
evangeliske Christus. Den evangeliske Christus staaer i virkelig Existens,
og gjør i ethvert Øieblik Existensens Fordring til hver Den, der nærmer sig
for at betragte ham; men det dogmatiske Christusbillede staaer kun i
Betragtningen,
% --- p. 109 ---
og tilhører det æsthetisk-videnskabelige Ideegallerie, hvor Enhver,
der har nogen „dybere Sands“, kan gaae hen at nyde Beskuelsen. At male et
Christusbillede med Ord og Tanker er vistnok en større Kunst end at male
det med Linier og Farver; Opgaven er da heller ikke endnu bleven løst af
nogen speculativ-protestantisk Dogmatiker, som den allerede for længst er
bleven løst af Middelalderens catholske Malere.

Men selv om Opgaven løstes med forbausende Held; selv om der opstod et
dogmatisk-speculativt Genie, hvem det lykkedes at male os et
Christusbillede med en saadan philosophisk Dybsindighed, en saadan
Fromhedens Inderlighed, en saadan Reenhed i Tanken, en saadan Klarhed i
Udtrykket, at Christus nu virkelig stod for os i dette Billede som „den
guddommelige Helligheds umiddelbare Fremtræden paa Jorden“, saa at „den
høieste Kjærlighed og Hjertets Ydmyghed og Hengivenhed i Gud her saaes at
smelte sammen til Eet med det, baade Frygt og Ærefrygt indgydende, Udtryk
af en Retfærdighed, hvis Grundvold Intet kan ryste, og med det mystiske
Udtryk af en Magt, der kunde byde alle Kræfter og alle Magter at tjene
sig;“ --- sæt, at Opgavens Løsning lykkedes Dogmatikeren efter en saa
overordentlig Maalestok: saa vilde et saadant Christusbillede dog alligevel
være ligesaa langt fra at gjøre et specifisk christeligt Indtryk, som
Malerens Billede.

Ligesom Malerens Billede nemlig gjør et æsthetisk-kunstnerisk Indtryk, og
derfor glæder og tiltrækker alle dem, der have Sands for den hellige Kunst
uden Hensyn til, om de i Virkeligheden ere troende eller vantroe; saaledes
vilde det æsthetisk-speculative Christusbillede gjøre et
æsthetisk-speculativt Indtryk, og derfor i høi Grad glæde og tiltrække alle
Dem, der have Sands for speculativ Phantasie-Anskuelse; men dette
Speculationens Billede vilde dog ingenlunde staae for Beskueren som
„\emph{Modsigelsens Tegn}“, vilde dog ikke have Kraft til at
„\emph{aabenbare Hjerternes Tanker},“
% --- p. 110 ---
vilde dog ikke være „\emph{mægtigt i Ord og
Gjerning}“: vilde dog ikke formaae at gjøre et sandt evangelisk Indtryk som
„\emph{Verdens Frelser!}“

Kan man ikke forholde sig æsthetisk-speculativt til Evangelietroens
Christus%
% printed footnote mark: *)
\footnote{Man læse, hvad dette Cardinalpunkt angaaer, „Philosophiske
Smuler“ af Johannes Climacus.}, saa kan man heller ikke forholde sig
evangelisk-troende til Speculationens Christusbillede.

Vel kan det stundom hænde sig, at En og Anden træder frem og bevidner, at
det nu virkelig er den dogmatiske Christologie, der ret egentlig først har
begrundet og klaret ham hans Tro, ja at han først i det
dybsindig-speculative Christusbillede ret egentlig har seet, hvo Christus
er; men med slige Vidnesbyrd bør man dog alligevel være forsigtig. Det er
desværre kun altfor sandt, at vi Mennesker i det 19de Aarhundrede ere meget
tilbøielige til at forvexle vor æsthetisk-speculative Interesse for Ideen
med den religiøs-ethiske Interesse for Sagen; hvor let kan det da ikke
skee, at En virkelig mener at være bestyrket i Christendommen, skjøndt han
dog, egentlig talt, kun er bleven bestyrket i æsthetisk-speculativ
Kjærlighed til et æsthetisk-speculativt Christusbillede!

Ad mange Veie kommer Christendommen til Menneskene, ad mange Veie komme
Menneskene til Christendommen. Vi ville jo gjerne indrømme, at den
theologiske Christologie ogsaa paa sin Viis kan blive Anledning til, at et
Menneske gaaer i sig selv, og besinder sig paa Evangeliet; men derfor maae
vi dog ligefuldt gjøre Forskjel mellem Sagen selv og Sagens Anledning.

Man kan med Grund sige, at det forholder sig med de dogmatiske
Christusbilleder i den protestantiske Kirke, som med de æsthetiske
Madonnabilleder i den catholske: der fortælles Adskilligt om deres
Omvendelseskraft, men denne Omvendelseskraft er noget tvetydig. For at hæve
Tvetydigheden, beslutte man sig til
% --- p. 111 ---
at give Troen, hvad Troens er, og Videnskaben, hvad Videnskabens er.

Ved denne Fordeling vil Dogmatikens halvspeculative Christusbillede,
forsaavidt det bestaaer i en mystisk Sammensmeltning af Tro og Viden, jo
vistnok for en stor Deel falme og svinde; men Evangeliets Christus vil
seire og vinde, naar det i Sandhed erkjendes og i Gjerning bevises,
% PRINTER'S DEFECT: the line ends „i Sandhed er=“ and the next line begins
% „erkjendes“, so the print doubles the „er“. Resolved here to the intended
% „erkjendes“, in line with this file's silent joining of end-of-line hyphens.
\emph{at Christendommen er høiere end Videnskaben, og Evangelietroen
ueensartet med Theologien.}

%% ============================================================
%%  NIENDE FORELÆSNING (pp. 112--128 = PDF 127--143)
%% ============================================================
\chapter*{Niende Forelæsning}
\addcontentsline{toc}{chapter}{Niende Forelæsning}
\label{ch:forel9}
\markboth{Niende Forelæsning}{}

\begin{center}\itshape
Theologiens halvspeculative Dogmatik kan ikke retfærdiggjøre sin kirkelige
Betydning af Kirkens Historie, da Kirkens Historie tvertimod beviser,
hvorledes dogmatiske Fordomme i denne Henseende have skadet Kirken, og
hvilken gjennemgribende Modsætning her er imellem den theologiske og den
kirkelige Dogmatik. Den theologiske Dogmatik maa falde for den Sætning,
at Christendommen er et Paradox: Evangelietroen er ueensartet med
Theologien.
\end{center}

% [text to be added: pp. 112--128]

%% ============================================================
%%  TIENDE FORELÆSNING (pp. 129--143 = PDF 144--158)
%% ============================================================
\chapter*{Tiende Forelæsning}
\addcontentsline{toc}{chapter}{Tiende Forelæsning}
\label{ch:forel10}
\markboth{Tiende Forelæsning}{}

\begin{center}\itshape
Ifølge den nu gjældende Methode befinder den theoretiske Theologie sig i
et aabenbart Misforhold til den religiøse Praxis; for at hævde
Troesprincipet maa den præstelige Dannelse emancipere sig fra Theorien;
thi Evangelietroen er ueensartet med Theologien.
\end{center}

% [text to be added: pp. 129--143]

%% ============================================================
%%  ELLEVTE FORELÆSNING (pp. 144--158 = PDF 159--173)
%% ============================================================
\chapter*{Ellevte Forelæsning}
\addcontentsline{toc}{chapter}{Ellevte Forelæsning}
\label{ch:forel11}
\markboth{Ellevte Forelæsning}{}

\begin{center}\itshape
Den theoretiske Theologie maa forholde sig skeptisk til Christendommens
aabenbarede Sandhed; den praktiske Theologie, der i forestillet Existens
arbeider paa at løse de christelige Opgaver, maa ende med at opløse sig
selv: Evangelietroen er ueensartet med Theologien.
\end{center}

% [text to be added: pp. 144--158]

%% ============================================================
%%  TOLVTE FORELÆSNING (pp. 159--174 = PDF 174--189)
%% ============================================================
\chapter*{Tolvte Forelæsning}
\addcontentsline{toc}{chapter}{Tolvte Forelæsning}
\label{ch:forel12}
\markboth{Tolvte Forelæsning}{}

\begin{center}\itshape
Hvis den praktiske Theologie ikke forstaaer at hæve sit theoretiske Skin,
vil dette Skin endnu langt mere mislede den religiøse Dannelse, end den
aabenbare Theorie. Den theoretiske Theologie, der ikke kan fornegte sin
væsentlige Forskjel fra den religiøse Praxis, vil ved at fremme sine
videnskabelige Interesser i Videnskabens Aand vedblive at virke
befrugtende og dannende paa Alle, som i Troens Inderlighed erkjende, at
Christendommen er høiere end Videnskaben, og Evangelietroen ueensartet med
Theologien.
\end{center}

% [text to be added: pp. 159--174]

\end{document}
